% =====================================================================
%  Anatomía de un LLM — Fase 0 — Módulo 0.8
%  Combinatoria y probabilidad
%  Compilar: latexmk -r ../latexmkrc P08-probabilidad.tex
% =====================================================================
\documentclass[11pt,a4paper,twoside,openany]{book}
\usepackage{../estilo-apuntes}

\newcommand{\ApuntesCurso}{Anatomía de un LLM}
\newcommand{\ApuntesModulo}{Módulo 0.8}

% Anclas de hiperenlace únicas por tema para los enunciados
\makeatletter
\@for\nombre:=definicion,teorema,proposicion,lema,corolario,ejemplo,observacion\do{%
  \expandafter\xdef\csname theH\nombre\endcsname{\noexpand\thechapter.\noexpand\arabic{definicion}}}
\makeatother

% --- Notación propia del módulo -------------------------------------------
\newcommand{\N}{\mathbb{N}}
\newcommand{\PP}{\mathbb{P}}
\newcommand{\Nor}{\mathcal{N}}
\newcommand{\Xbar}{\overline{X}}
\newcommand{\dsim}{\mathbin{\triangle}}
\DeclareMathOperator{\Var}{Var}
\DeclareMathOperator{\Ber}{Ber}
\DeclareMathOperator{\Bin}{Bin}
\DeclareMathOperator{\Geom}{Geom}
\DeclareMathOperator{\Poi}{Poisson}
\DeclareMathOperator{\Cat}{Cat}
\DeclareMathOperator{\Unif}{Unif}
\DeclareMathOperator{\Expo}{Exp}

\hypersetup{pdftitle={Módulo 0.8. Combinatoria y probabilidad},pdfauthor={Anatomía de un LLM}}

\begin{document}
\frontmatter
% ------------------------------------------------------------------ Portada
\begin{titlepage}
\centering
\vspace*{2.5cm}
{\large\scshape Anatomía de un LLM}\par
\vspace{0.4cm}
{\normalsize Fundamentos matemáticos de los modelos de lenguaje}\par
\vspace{2.2cm}
\rule{\textwidth}{0.8pt}\par\vspace{0.5cm}
{\small\scshape Fase 0 \textperiodcentered\ Preliminares matemáticos}\par\vspace{0.3cm}
{\Huge\bfseries Módulo 0.8\par}\vspace{0.3cm}
{\LARGE Combinatoria y probabilidad\par}
\vspace{0.5cm}\rule{\textwidth}{0.4pt}\par
\vspace{1.6cm}
{\large Apuntes del curso}\par
\vfill
{\small Edición de octubre de 2026}\par
\end{titlepage}

% ------------------------------------------------------- Presentación
\chapter*{Presentación del módulo}
\addcontentsline{toc}{chapter}{Presentación del módulo}
\markboth{Presentación del módulo}{}

Un modelo de lenguaje es una distribución de probabilidad sobre secuencias de tokens. Para entender qué significa esa frase, y para calcular con ella, hacen falta dos herramientas: contar con seguridad (cuántas secuencias, cuántos subconjuntos, cuántos repartos) y razonar con probabilidades (condicionar, aplicar el teorema de Bayes, calcular esperanzas y acotar desviaciones). Este módulo repasa ambas desde el nivel de bachillerato con el rigor de un primer curso universitario de probabilidad. Prepara los módulos~1 («Modelos de lenguaje probabilísticos») y~5 («Probabilidad y estadística») sin sustituirlos: aquí se fijan las definiciones, se demuestran los resultados básicos y se practica mucho; allí se aplican a los modelos de lenguaje y se amplían con la estimación, la inferencia bayesiana y el teorema central del límite.

\section*{Objetivos}
Al finalizar el módulo, el estudiante debe ser capaz de:
\begin{enumerate}[label=\arabic*.]
  \item aplicar las reglas de la suma, del producto, del complemento, de la biyección y de la división, y el principio de inclusión--exclusión;
  \item distinguir y contar variaciones, permutaciones y combinaciones, con y sin repetición, y manejar los números combinatorios, el binomio y el multinomio de Newton;
  \item demostrar identidades combinatorias por doble conteo y mediante biyecciones;
  \item enunciar los axiomas de Kolmogorov, deducir de ellos las reglas de cálculo y resolver problemas con el modelo de Laplace, incluido el problema del cumpleaños;
  \item calcular probabilidades condicionadas y aplicar la regla de la cadena, la probabilidad total y el teorema de Bayes (pruebas diagnósticas, Monty Hall), y decidir y demostrar la independencia de sucesos;
  \item manejar variables aleatorias discretas y continuas, calcular esperanzas y varianzas, explotar la linealidad de la esperanza con indicadoras y reconocer las distribuciones Bernoulli, binomial, geométrica, Poisson, categórica, uniforme, exponencial y normal;
  \item demostrar las desigualdades de Markov y Chebyshev y la ley débil de los grandes números, y usarlas para acotar probabilidades y tamaños muestrales.
\end{enumerate}

\section*{Requisitos previos}
Matemáticas de bachillerato. Se usan la inducción y las técnicas de demostración del módulo~0.1, la teoría de conjuntos y funciones del módulo~0.2 (inyectividad, biyecciones, principio del palomar) y los sumatorios y la serie geométrica del módulo~0.3. El \cref{tema:6} requiere el cálculo integral de una variable (integrales definidas, cambio de variable e integración por partes) de los módulos de cálculo de esta fase; cuando se usa una serie o una integral no elemental se indica explícitamente.

\section*{Organización}
\begin{center}
\begin{tabular}{@{}cll@{}}
\toprule
Tema & Contenido & Fuentes principales \\
\midrule
1 & Principios de conteo y ordenaciones & \cite{mcs}, \cite{ross}, \cite{bh} \\
2 & Combinaciones, números combinatorios y doble conteo & \cite{mcs}, \cite{bh}, \cite{stat110} \\
3 & Espacios de probabilidad & \cite{bt}, \cite{ross}, \cite{mcs} \\
4 & Probabilidad condicionada, independencia y Bayes & \cite{bh}, \cite{bt}, \cite{mcs} \\
5 & Variables aleatorias discretas & \cite{bh}, \cite{bt}, \cite{ross} \\
6 & Variables continuas y ley de los grandes números & \cite{bh}, \cite{ross}, \cite{gs} \\
\bottomrule
\end{tabular}
\end{center}

Cada tema termina con una lista de lecturas y con ejercicios graduados en tres niveles: $\star$ comprensión, $\star\star$ cálculo y $\star\star\star$ demostración. Los ejercicios marcados con \emph{Fuente} están adaptados de cursos y libros de texto universitarios; el enunciado se ha redactado de nuevo y, en algunos casos, se han cambiado los datos. Las soluciones completas, paso a paso, están en el \cref{ap:soluciones}.

\section*{Dónde se usa en el curso}
\begin{center}\small
\begin{tabular}{@{}p{0.44\textwidth}p{0.48\textwidth}@{}}
\toprule
Concepto & Módulos del curso que lo emplean \\
\midrule
Reglas del producto y de la suma, $\abs{\V}^{T}$ secuencias & M01 (espacio de secuencias), M12 (bytes, UTF-8) \\
Números combinatorios, pares y subconjuntos & M12 (recuento de pares en BPE), M22 (pass@$k$) \\
Coeficiente multinomial & M01 (verosimilitud de un corpus), M05 (distribución multinomial) \\
Espacio muestral, sucesos, axiomas & M01 (modelo de lenguaje como ley de probabilidad), M05 \\
Problema del cumpleaños, cota de la unión & M16 (colisiones de \emph{hash}, MinHash) \\
Probabilidad condicionada y regla de la cadena & M01 (factorización autorregresiva), M13--M15 \\
Teorema de Bayes, independencia condicional & M05 (inferencia bayesiana), M08 (clasificadores), M01 (hipótesis de Markov) \\
Distribución categórica y muestreo & M01 y M05 (salida softmax), M19 (decodificación) \\
Esperanza, linealidad, indicadoras & M05, M06 (entropía como esperanza), M22 (métricas) \\
Varianza de una media, $\sigma^{2}/n$ & M05 (Monte Carlo), M07 (gradiente estocástico) \\
Normal y tipificación & M05 (inicialización), M11 (GELU, normalización) \\
Markov, Chebyshev y ley de los grandes números & M05 (teoremas límite), M01 (entropía cruzada empírica) \\
\bottomrule
\end{tabular}
\end{center}

\section*{Sobre las fuentes}
El texto es de redacción propia. La parte de combinatoria sigue el capítulo~14, «Cardinality Rules», de Lehman, Leighton y Meyer~\cite{mcs}, libro del curso 6.042J del MIT~\cite{mit6042}, y el capítulo~1 de Ross~\cite{ross}; las demostraciones por doble conteo siguen el estilo de las \emph{story proofs} de Blitzstein y Hwang~\cite{bh}, libro del curso Stat~110 de Harvard. La probabilidad sigue a Bertsekas y Tsitsiklis~\cite{bt}, texto del curso 6.041 del MIT, a Blitzstein y Hwang~\cite{bh} y a Ross~\cite{ross}, con ejemplos de~\cite{mcs} y de Grinstead y Snell~\cite{gs}. Los ejercicios universitarios proceden de las hojas \emph{Strategic Practice and Homework} de Stat~110~\cite{stat110}, de los \emph{Problem Sets} de MIT 6.041~\cite{mit6041f10} y MIT 18.05~\cite{mit1805} y de los ejemplos de~\cite{mcs}.

\section*{Notación}
\begin{center}
\begin{tabular}{@{}ll@{}}
\toprule
Símbolo & Significado \\
\midrule
$\abs{A}$, $[n]$ & cardinal de un conjunto finito; $[n]=\{1,\dots,n\}$ \\
$n!$, $n^{\underline{k}}$ & factorial; factorial descendente $n(n-1)\cdots(n-k+1)$ \\
$\binom nk$, $\binom{n}{k_1,\dots,k_r}$ & número combinatorio; coeficiente multinomial \\
$\Omega$, $A^{c}$, $A\dsim B$ & espacio muestral; suceso contrario; diferencia simétrica \\
$\PP(A)$, $\PP(A\mid B)$ & probabilidad; probabilidad condicionada \\
$I_A$ & variable indicadora del suceso $A$ \\
$p_X$, $f_X$, $F_X$ & función de masa, densidad y función de distribución de $X$ \\
$\E[X]$, $\Var(X)$, $\sigma_X$ & esperanza, varianza, desviación típica \\
$\Ber(p)$, $\Bin(n,p)$, $\Geom(p)$ & distribuciones discretas \\
$\Poi(\lambda)$, $\Cat(\bm p)$ & distribuciones de Poisson y categórica \\
$\Unif(a,b)$, $\Expo(\lambda)$, $\Nor(\mu,\sigma^{2})$ & distribuciones continuas \\
$\varphi$, $\Phi$ & densidad y función de distribución de la normal estándar \\
$\Xbar_n$ & media muestral $\frac1n\sum_{i=1}^{n}X_i$ \\
\bottomrule
\end{tabular}
\end{center}

\tableofcontents

\mainmatter
% =====================================================================
\chapter{Principios de conteo y ordenaciones}
\label{tema:1}
% =====================================================================

Contar bien es el primer paso para calcular probabilidades en espacios finitos. Este tema reúne las reglas básicas del recuento (suma, producto, complemento, biyección y división), el principio de inclusión--exclusión para dos y tres conjuntos y el recuento de ordenaciones: variaciones y permutaciones, con y sin repetición. La presentación sigue el capítulo~14, «Cardinality Rules», de Lehman, Leighton y Meyer~\cite{mcs}, texto del curso 6.042J del MIT; el capítulo~1 de Ross~\cite{ross} (§1.2 y §1.3); la sección~1.4, «How to count», de Blitzstein y Hwang~\cite{bh}, y la sección~3.1 de Grinstead y Snell~\cite{gs}.

\section{Reglas de la suma, del producto y del complemento}

Se recuerda que, para un conjunto finito $A$, $\abs{A}$ denota su número de elementos y que $[n]=\{1,2,\dots,n\}$. Contar un conjunto consiste en hallar $\abs{A}$; casi siempre se hace descomponiendo $A$ en piezas más sencillas o poniéndolo en correspondencia con otro conjunto ya contado.

\begin{proposicion}[Regla de la suma]\label{prop:t1-suma}
Si $A_1,\dots,A_m$ son conjuntos finitos disjuntos dos a dos, entonces
\[
\abs{A_1\cup\dots\cup A_m}=\abs{A_1}+\dots+\abs{A_m}.
\]
\end{proposicion}

\begin{proof}
Para $m=2$: si $A_1=\{a_1,\dots,a_p\}$ y $A_2=\{b_1,\dots,b_q\}$ son disjuntos, la lista $a_1,\dots,a_p,b_1,\dots,b_q$ enumera $A_1\cup A_2$ sin repeticiones, luego tiene $p+q$ elementos. El caso general se obtiene por inducción sobre $m$, aplicando el caso $m=2$ a $A_1\cup\dots\cup A_{m-1}$ y $A_m$, que son disjuntos.
\end{proof}

\begin{proposicion}[Regla del producto]\label{prop:t1-producto}
Si $A_1,\dots,A_k$ son conjuntos finitos, $\abs{A_1\times\dots\times A_k}=\abs{A_1}\cdots\abs{A_k}$.
\end{proposicion}

\begin{proof}
Para $k=2$, $A_1\times A_2=\bigcup_{a\in A_1}\bigl(\{a\}\times A_2\bigr)$ es una unión disjunta de $\abs{A_1}$ conjuntos, cada uno con $\abs{A_2}$ elementos; por la regla de la suma, $\abs{A_1\times A_2}=\abs{A_1}\abs{A_2}$. El caso general se demuestra por inducción, identificando $A_1\times\dots\times A_k$ con $(A_1\times\dots\times A_{k-1})\times A_k$.
\end{proof}

En la práctica se usa una versión más flexible, en la que el conjunto de opciones de cada etapa puede depender de las elecciones previas, pero no su número.

\begin{teorema}[Regla del producto generalizada]\label{teo:t1-producto-gen}
Sea $S$ un conjunto de sucesiones $(s_1,\dots,s_k)$ tal que hay $n_1$ valores posibles de $s_1$ y, para cada elección de $s_1,\dots,s_{i-1}$, hay exactamente $n_i$ valores posibles de $s_i$ ($i=2,\dots,k$). Entonces $\abs{S}=n_1n_2\cdots n_k$.
\end{teorema}

\begin{proof}
Por inducción sobre $k$. Si $k=1$, $\abs{S}=n_1$. Supóngase cierto para sucesiones de longitud $k-1$ y sea $S'$ el conjunto de los prefijos $(s_1,\dots,s_{k-1})$ de elementos de $S$; por la hipótesis de inducción, $\abs{S'}=n_1\cdots n_{k-1}$. Agrupando los elementos de $S$ según su prefijo, $S$ es la unión disjunta, sobre los prefijos $p\in S'$, de los conjuntos $\{(p,s_k)\in S\}$, cada uno con $n_k$ elementos. Por la regla de la suma, $\abs{S}=\abs{S'}\,n_k=n_1\cdots n_k$.
\end{proof}

\begin{proposicion}[Regla del complemento]\label{prop:t1-complemento}
Si $A\subseteq U$ y $U$ es finito, $\abs{U\setminus A}=\abs{U}-\abs{A}$.
\end{proposicion}

\begin{proof}
$U$ es la unión disjunta de $A$ y $U\setminus A$; basta aplicar la regla de la suma.
\end{proof}

\begin{ejemplo}[Contraseñas]\label{ej:t1-contrasenas}
Una contraseña está formada por 6 caracteres, cada uno una letra minúscula (26) o un dígito (10). Por la regla del producto hay $36^{6}=2\,176\,782\,336$ contraseñas. Las que contienen \emph{al menos} un dígito se cuentan mejor por complemento: las que no tienen ninguno son $26^{6}=308\,915\,776$, luego hay $36^{6}-26^{6}=1\,867\,866\,560$. Intentar contarlas directamente («elegir la posición del dígito, el dígito y el resto») cuenta varias veces las contraseñas con más de un dígito.
\end{ejemplo}

\section{Reglas de la biyección y de la división}

\begin{proposicion}[Regla de la biyección]\label{prop:t1-biyeccion}
Si existe una biyección $f\colon A\to B$ entre conjuntos finitos, entonces $\abs{A}=\abs{B}$.
\end{proposicion}

La demostración es inmediata: si $A=\{a_1,\dots,a_n\}$ sin repeticiones, entonces $f(a_1),\dots,f(a_n)$ son distintos (inyectividad) y agotan $B$ (sobreyectividad). Esta regla permite contar un conjunto contando otro~\cite[§14.1]{mcs}.

\begin{ejemplo}[Subconjuntos y cadenas de bits]\label{ej:t1-bits}
Sea $A=\{a_1,\dots,a_n\}$. La aplicación que asigna a cada $X\subseteq A$ la cadena $b_1\cdots b_n$ con $b_i=1$ si $a_i\in X$ y $b_i=0$ en otro caso es una biyección entre $\mathcal{P}(A)$ y $\{0,1\}^{n}$. Por las reglas de la biyección y del producto, $A$ tiene $2^{n}$ subconjuntos.
\end{ejemplo}

\begin{proposicion}[Regla de la división]\label{prop:t1-division}
Sea $f\colon A\to B$ una función sobreyectiva entre conjuntos finitos tal que cada $b\in B$ tiene exactamente $d$ preimágenes ($f$ es \emph{$d$ a~1}). Entonces $\abs{A}=d\,\abs{B}$.
\end{proposicion}

\begin{proof}
Los conjuntos $f^{-1}(\{b\})$, con $b\in B$, son disjuntos, cada uno tiene $d$ elementos y su unión es $A$. Por la regla de la suma, $\abs{A}=\sum_{b\in B}d=d\abs{B}$.
\end{proof}

\section{Inclusión--exclusión para dos y tres conjuntos}

Cuando los conjuntos no son disjuntos, la suma de sus cardinales cuenta dos veces los elementos comunes y hay que corregirla.

\begin{teorema}[Inclusión--exclusión]\label{teo:t1-ie}
Para conjuntos finitos $A$, $B$ y $C$:
\begin{align*}
\abs{A\cup B}&=\abs{A}+\abs{B}-\abs{A\cap B},\\
\abs{A\cup B\cup C}&=\abs{A}+\abs{B}+\abs{C}-\abs{A\cap B}-\abs{A\cap C}-\abs{B\cap C}+\abs{A\cap B\cap C}.
\end{align*}
\end{teorema}

\begin{proof}
$A\cup B$ es la unión disjunta de $A$ y $B\setminus A$, y $B$ es la unión disjunta de $A\cap B$ y $B\setminus A$. Por la regla de la suma, $\abs{A\cup B}=\abs{A}+\abs{B\setminus A}=\abs{A}+\abs{B}-\abs{A\cap B}$. Para tres conjuntos se aplica la fórmula anterior a $A\cup B$ y $C$, y se usa la distributiva $(A\cup B)\cap C=(A\cap C)\cup(B\cap C)$:
\begin{align*}
\abs{A\cup B\cup C}&=\abs{A\cup B}+\abs{C}-\abs{(A\cap C)\cup(B\cap C)}\\
&=\abs{A}+\abs{B}-\abs{A\cap B}+\abs{C}-\bigl(\abs{A\cap C}+\abs{B\cap C}-\abs{A\cap B\cap C}\bigr),
\end{align*}
que es la fórmula del enunciado.
\end{proof}

La versión para $m$ conjuntos se demuestra en el \cref{tema:2}, una vez disponibles los números combinatorios.

\begin{ejemplo}\label{ej:t1-divisores}
¿Cuántos enteros de $\{1,\dots,100\}$ son divisibles por 2 o por 3? Sean $A_2$ y $A_3$ los múltiplos de 2 y de 3. Entonces $\abs{A_2}=50$, $\abs{A_3}=\lfloor 100/3\rfloor=33$ y $A_2\cap A_3$ son los múltiplos de 6, con $\lfloor100/6\rfloor=16$ elementos. Hay $50+33-16=67$.
\end{ejemplo}

\section{Variaciones y permutaciones}

Se trata ahora de contar \emph{selecciones ordenadas} de elementos de un conjunto de $n$ elementos. La terminología española distingue según se permita o no repetir.

\begin{definicion}[Variaciones y permutaciones]
Sea $A$ un conjunto con $n$ elementos y $k\ge 0$ un entero.
\begin{enumerate}
  \item Una \emph{variación con repetición} de orden $k$ es una sucesión $(a_1,\dots,a_k)$ de elementos de $A$, posiblemente repetidos; equivale a una función $[k]\to A$.
  \item Una \emph{variación sin repetición} de orden $k\le n$ es una sucesión de $k$ elementos \emph{distintos} de $A$; equivale a una función inyectiva $[k]\to A$.
  \item Una \emph{permutación} de $A$ es una variación sin repetición de orden $n$, es decir, una ordenación de todos sus elementos; equivale a una biyección $[n]\to A$.
\end{enumerate}
Se escribe $n!=1\cdot2\cdots n$, con $0!=1$, y $n^{\underline{k}}=n(n-1)\cdots(n-k+1)$ (\emph{factorial descendente}, con $n^{\underline{0}}=1$).
\end{definicion}

\begin{teorema}\label{teo:t1-variaciones}
Si $\abs{A}=n$, el número de variaciones con repetición de orden $k$ es $n^{k}$; el de variaciones sin repetición de orden $k\le n$ es $n^{\underline{k}}=n!/(n-k)!$, y el de permutaciones es $n!$.
\end{teorema}

\begin{proof}
En el primer caso, cada una de las $k$ posiciones admite $n$ valores y se aplica la regla del producto. En el segundo, $a_1$ admite $n$ valores; elegidos $a_1,\dots,a_{i-1}$, distintos, $a_i$ admite exactamente $n-(i-1)$ valores, sean cuales sean los anteriores. Por la regla del producto generalizada (\cref{teo:t1-producto-gen}), hay $n(n-1)\cdots(n-k+1)$ sucesiones. El tercer caso es el segundo con $k=n$.
\end{proof}

\begin{ejemplo}[Permutaciones circulares]\label{ej:t1-circular}
¿De cuántas formas pueden sentarse $n\ge1$ personas alrededor de una mesa redonda si solo importa quién tiene a quién a su derecha? Cada permutación $(a_1,\dots,a_n)$ determina una disposición circular (sentando a $a_1,a_2,\dots$ en sentido horario a partir de una silla fija), y las $n$ rotaciones de una permutación dan la misma disposición. La aplicación «permutación $\mapsto$ disposición» es, por tanto, $n$ a~1, y por la regla de la división hay $n!/n=(n-1)!$ disposiciones~\cite[§14.4.2]{mcs}.
\end{ejemplo}

\begin{ejemplo}[Funciones entre conjuntos finitos]
Sean $\abs{A}=k$ y $\abs{B}=n$. Hay $n^{k}$ funciones $A\to B$ y $n^{\underline{k}}$ inyectivas. Si $k=n$, las inyectivas son biyectivas y hay $n!$. Por ejemplo, con $k=3$ y $n=5$ hay $125$ funciones, de las que $5\cdot4\cdot3=60$ son inyectivas.
\end{ejemplo}

\begin{observacion}[Conexión con el curso]
Con un vocabulario de $K$ tokens hay $K^{T}$ secuencias de longitud $T$ (variaciones con repetición), y el modelo de lenguaje del módulo~1 asigna una probabilidad a cada una. Con $K=50\,000$ y $T=3$ ya hay $1{,}25\cdot10^{14}$ secuencias, lo que explica por qué una distribución sobre secuencias no puede almacenarse como tabla y se factoriza con la regla de la cadena (\cref{tema:4}). En el módulo~12, el alfabeto de bytes tiene $2^{8}=256$ símbolos y UTF-8 codifica cada carácter con 1 a 4 bytes; el número de cadenas de bytes de cada longitud se cuenta con la regla del producto.
\end{observacion}

\begin{observacion}[Errores frecuentes]
(1) Sumar cuando hay que multiplicar, o al revés: la regla de la suma corresponde a «o» entre casos \emph{disjuntos}; la del producto, a «y» entre etapas sucesivas. (2) Aplicar la regla de la suma a casos que se solapan: hay que usar inclusión--exclusión. (3) Contar «al menos uno» eligiendo primero uno y luego el resto, lo que cuenta varias veces los casos con más de uno; es mejor pasar al complemento. (4) Olvidar si el orden importa: $(a,b)$ y $(b,a)$ son variaciones distintas pero el mismo subconjunto. (5) Aplicar la regla del producto generalizada cuando el número de opciones de una etapa \emph{depende} de las elecciones previas.
\end{observacion}

\begin{lecturas}
\item Lehman, Leighton y Meyer~\cite{mcs}, cap.~14 «Cardinality Rules», §14.1 a §14.4.
\item Ross~\cite{ross}, cap.~1, §1.2 «The Basic Principle of Counting» y §1.3 «Permutations».
\item Blitzstein y Hwang~\cite{bh}, §1.4 «How to count».
\item Grinstead y Snell~\cite{gs}, §3.1 «Permutations».
\item Bertsekas y Tsitsiklis~\cite{bt}, §1.6 (clase~4 de MIT 6.041~\cite{mit6041f10}, «Counting»).
\end{lecturas}

\section*{Ejercicios del tema 1}
\addcontentsline{toc}{section}{Ejercicios del tema 1}

\begin{ejercicio}{1}{eje:1-reglas}
Indique qué regla (suma, producto, complemento, biyección o división) aplicaría en cada caso y dé el resultado: (a) número de menús con un primero entre 4, un segundo entre 5 y un postre entre 3; (b) número de formas de elegir \emph{un} plato, sea primero, segundo o postre; (c) número de cadenas de 10 bits que no son todas ceros; (d) número de subconjuntos de un conjunto de 10 elementos.
\end{ejercicio}

\begin{ejercicio}{1}{eje:1-vocabulario}
Un vocabulario tiene $K=50\,000$ tokens. (a) ¿Cuántas secuencias de 3 tokens hay? (b) ¿Cuántas sin tokens repetidos? (c) ¿Qué fracción de las secuencias de (a) tiene algún token repetido? Dé el resultado con tres cifras significativas.
\end{ejercicio}

\begin{ejercicio}{2}{eje:1-contrasenas}
Una contraseña tiene entre 6 y 8 caracteres, cada uno una letra minúscula (26) o un dígito (10), y debe contener al menos un dígito. ¿Cuántas contraseñas válidas hay?
\fuente{adaptado de Lehman, Leighton y Meyer~\cite{mcs}, §14.2.4 «Counting Passwords».}
\end{ejercicio}

\begin{ejercicio}{2}{eje:1-divisibles}
¿Cuántos enteros de $\{1,\dots,1000\}$ son divisibles por 2, por 3 o por~5? ¿Cuántos no son divisibles por ninguno de los tres?
\end{ejercicio}

\begin{ejercicio}{2}{eje:1-palindromos}
Una \emph{palabra} es cualquier cadena de letras de un alfabeto de 26. (a) ¿Cuántas palabras de 2 letras son palíndromos? ¿Y de 3 letras? (b) Compare la proporción de palíndromos entre las palabras de 2 letras y entre las de 3. (c) Generalice: ¿cuántos palíndromos de longitud $n$ hay?
\fuente{adaptado de Harvard Stat 110~\cite{stat110}, \emph{Strategic Practice}~1, problema~1(b).}
\end{ejercicio}

\begin{ejercicio}{2}{eje:1-mesa}
Ocho personas se sientan alrededor de una mesa redonda; dos disposiciones son iguales si una es rotación de la otra. (a) ¿Cuántas disposiciones hay? (b) ¿En cuántas se sientan juntas dos personas concretas, Ana y Blas? (c) ¿Y en cuántas \emph{no} se sientan juntas?
\fuente{adaptado de Lehman, Leighton y Meyer~\cite{mcs}, §14.4.2 «Knights of the Round Table».}
\end{ejercicio}

\begin{ejercicio}{2}{eje:1-cifras}
¿Cuántos números de cuatro cifras (de 1000 a 9999) tienen todas sus cifras distintas? ¿Cuántos de ellos son pares?
\end{ejercicio}

\begin{ejercicio}{2}{eje:1-funciones}
Sean $A=\{1,2,3,4\}$ y $B=\{x,y\}$. (a) ¿Cuántas funciones $A\to B$ hay? (b) ¿Cuántas son inyectivas? (c) Use el complemento para contar las sobreyectivas. (d) Cuente, mediante inclusión--exclusión, las funciones sobreyectivas de un conjunto de 5 elementos en uno de~3.
\end{ejercicio}

\begin{ejercicio}{2}{eje:1-sombreros}
Se reparten al azar los sombreros de $n$ personas, uno a cada una; un reparto es una permutación de los $n$ sombreros. (a) ¿En cuántos repartos recibe cada persona su sombrero? (b) ¿En cuántos reciben su sombrero las $m$ primeras personas ($m\le n$)? (c) Si $2m\le n$, ¿en cuántos reciben las $m$ primeras personas sombreros que pertenecen a las $m$ últimas?
\fuente{adaptado de MIT 6.041~\cite{mit6041f10}, \emph{Problem Set}~3, problema~1, apartados (a) a~(c).}
\end{ejercicio}

\begin{ejercicio}{2}{eje:1-encuesta}
En un grupo de 60 estudiantes, 30 cursan Cálculo, 25 Álgebra y 20 Programación; 10 cursan Cálculo y Álgebra, 8 Cálculo y Programación, 7 Álgebra y Programación, y 3 las tres. ¿Cuántos no cursan ninguna? ¿Cuántos cursan exactamente una?
\end{ejercicio}

\begin{ejercicio}{3}{eje:1-producto-ind}
Demuestre por inducción sobre $k$ que el número de funciones inyectivas de $[k]$ en un conjunto de $n$ elementos es $n^{\underline{k}}$ para $0\le k\le n$, sin usar la regla del producto generalizada. Deduzca que no hay ninguna si $k>n$ (principio del palomar).
\end{ejercicio}

\begin{ejercicio}{3}{eje:1-fibonacci}
Sea $a_n$ el número de cadenas de $n$ bits sin dos unos consecutivos. Demuestre que $a_1=2$, $a_2=3$ y $a_n=a_{n-1}+a_{n-2}$ para $n\ge3$, justificando cada paso con las reglas de la suma y de la biyección. Calcule $a_8$.
\end{ejercicio}

\begin{ejercicio}{3}{eje:1-ie-cuatro}
A partir del \cref{teo:t1-ie} para tres conjuntos, demuestre la fórmula de inclusión--exclusión para cuatro conjuntos finitos $A,B,C,D$.
\end{ejercicio}

\begin{ejercicio}{3}{eje:1-collares}
(a) Demuestre con la regla de la división que el número de formas de sentar a $n\ge3$ personas alrededor de una mesa, si además se identifican dos disposiciones que son simétricas (una es la otra leída en sentido contrario), es $(n-1)!/2$. (b) Demuestre que el número de formas de repartir $2n$ personas en $n$ parejas (sin orden entre parejas ni dentro de ellas) es $(2n)!/(2^{n}n!)$.
\end{ejercicio}

% =====================================================================
\chapter{Combinaciones, números combinatorios y doble conteo}
\label{tema:2}
% =====================================================================

Este tema cuenta \emph{selecciones no ordenadas}: subconjuntos, multiconjuntos y repartos. Aparecen los números combinatorios, el binomio y el multinomio de Newton, la técnica del doble conteo y la forma general del principio de inclusión--exclusión. Se siguen las secciones §14.5 a §14.7 de Lehman, Leighton y Meyer~\cite{mcs}; las secciones §1.4 a §1.6 de Ross~\cite{ross}; las secciones §1.4 y §1.5, «Story proofs», de Blitzstein y Hwang~\cite{bh}, cuyo estilo de «demostración mediante una historia» se adopta aquí, y la sección~3.2 de Grinstead y Snell~\cite{gs}.

\section{Combinaciones y números combinatorios}

\begin{definicion}[Combinaciones]
Sea $A$ un conjunto con $n$ elementos y $0\le k\le n$. Una \emph{combinación} (sin repetición) de orden $k$ de $A$ es un subconjunto de $A$ con $k$ elementos. Su número se denota $\binom{n}{k}$ y se llama \emph{número combinatorio} («$n$ sobre $k$»). Por convenio, $\binom{n}{k}=0$ si $k<0$ o $k>n$.
\end{definicion}

\begin{teorema}\label{teo:t2-binom}
Para $0\le k\le n$,
\[
\binom{n}{k}=\frac{n^{\underline{k}}}{k!}=\frac{n!}{k!\,(n-k)!}.
\]
\end{teorema}

\begin{proof}
Sea $V$ el conjunto de variaciones sin repetición de orden $k$ de $A$ y $C$ el de subconjuntos de $k$ elementos. La aplicación $V\to C$ que olvida el orden, $(a_1,\dots,a_k)\mapsto\{a_1,\dots,a_k\}$, es sobreyectiva y cada subconjunto tiene exactamente $k!$ preimágenes, sus $k!$ ordenaciones. Por la regla de la división (\cref{prop:t1-division}), $\abs{V}=k!\,\abs{C}$, y $\abs{V}=n^{\underline{k}}$ por el \cref{teo:t1-variaciones}.
\end{proof}

\begin{proposicion}[Propiedades básicas]\label{prop:t2-props}
Para enteros $n\ge1$ y $k$:
\begin{enumerate}
  \item simetría: $\binom{n}{k}=\binom{n}{n-k}$;
  \item fórmula de Pascal: $\binom{n}{k}=\binom{n-1}{k}+\binom{n-1}{k-1}$;
  \item absorción: $k\binom{n}{k}=n\binom{n-1}{k-1}$.
\end{enumerate}
\end{proposicion}

\begin{proof}
Se dan demostraciones combinatorias; todas pueden comprobarse también con la fórmula del \cref{teo:t2-binom}.

(i) La aplicación $X\mapsto A\setminus X$ es una biyección entre los subconjuntos de $k$ elementos y los de $n-k$ elementos (es su propia inversa).

(ii) Se fija $a\in A$. Los subconjuntos de $k$ elementos de $A$ se dividen en dos clases disjuntas: los que no contienen a $a$, que son los subconjuntos de $k$ elementos de $A\setminus\{a\}$, en número $\binom{n-1}{k}$, y los que contienen a $a$, que se corresponden biyectivamente (quitando $a$) con los subconjuntos de $k-1$ elementos de $A\setminus\{a\}$, en número $\binom{n-1}{k-1}$. Se aplica la regla de la suma.

(iii) Se cuentan de dos formas los pares $(X,p)$ con $X\subseteq A$, $\abs{X}=k$ y $p\in X$ (un comité de $k$ personas con un presidente que pertenece a él). Eligiendo primero el comité y luego el presidente: $\binom{n}{k}\cdot k$. Eligiendo primero el presidente y luego los otros $k-1$ miembros entre las $n-1$ personas restantes: $n\binom{n-1}{k-1}$.
\end{proof}

La fórmula de Pascal permite construir el \emph{triángulo de Pascal}, en el que cada número es la suma de los dos que tiene encima:
\[
\begin{array}{c}
1\\
1\quad 1\\
1\quad 2\quad 1\\
1\quad 3\quad 3\quad 1\\
1\quad 4\quad 6\quad 4\quad 1\\
1\quad 5\quad 10\quad 10\quad 5\quad 1
\end{array}
\]

\section{Binomio y multinomio de Newton}

\begin{teorema}[Binomio de Newton]\label{teo:t2-binomio}
Para $x,y$ reales y $n\ge0$ entero,
\[
(x+y)^{n}=\sum_{k=0}^{n}\binom{n}{k}x^{k}y^{n-k}.
\]
\end{teorema}

\begin{proof}
Al desarrollar el producto $(x+y)(x+y)\cdots(x+y)$ de $n$ factores, sin conmutar, se obtiene la suma de las $2^{n}$ palabras que resultan de elegir $x$ o $y$ en cada factor. Una palabra con $k$ letras $x$ vale $x^{k}y^{n-k}$, y hay tantas como formas de elegir los $k$ factores de los que se toma $x$, es decir, $\binom{n}{k}$.
\end{proof}

El módulo~0.3 demuestra este teorema por inducción; la demostración anterior es la combinatoria, que explica \emph{por qué} aparecen los números combinatorios.

\begin{corolario}\label{cor:t2-sumas}
Para $n\ge 0$, $\sum_{k=0}^{n}\binom{n}{k}=2^{n}$, y para $n\ge1$, $\sum_{k=0}^{n}(-1)^{k}\binom{n}{k}=0$.
\end{corolario}

\begin{proof}
Basta tomar $x=y=1$ y $x=-1$, $y=1$ en el \cref{teo:t2-binomio}. La primera identidad tiene también una demostración combinatoria: ambos lados cuentan los subconjuntos de un conjunto de $n$ elementos, el izquierdo agrupándolos por tamaño~\cite[\emph{Strategic Practice}~1, problema~5]{stat110}.
\end{proof}

\begin{definicion}[Coeficiente multinomial]
Sean $n=k_1+\dots+k_r$ con $k_i\ge0$ enteros. Se define
\[
\binom{n}{k_1,k_2,\dots,k_r}=\frac{n!}{k_1!\,k_2!\cdots k_r!}.
\]
\end{definicion}

\begin{teorema}[Permutaciones con repetición]\label{teo:t2-multinomial}
El número de sucesiones de longitud $n$ formadas con $k_1$ copias del símbolo $s_1$, $k_2$ del símbolo $s_2$, \dots, $k_r$ del símbolo $s_r$ es $\binom{n}{k_1,\dots,k_r}$. Equivalentemente, es el número de formas de repartir $n$ objetos distintos en $r$ cajas distinguibles, con $k_i$ objetos en la caja $i$.
\end{teorema}

\begin{proof}
Una sucesión queda determinada eligiendo las $k_1$ posiciones de $s_1$, después las $k_2$ posiciones de $s_2$ entre las $n-k_1$ restantes, y así sucesivamente. Por la regla del producto generalizada,
\[
\binom{n}{k_1}\binom{n-k_1}{k_2}\cdots\binom{n-k_1-\dots-k_{r-1}}{k_r}
=\frac{n!}{k_1!\,(n-k_1)!}\cdot\frac{(n-k_1)!}{k_2!\,(n-k_1-k_2)!}\cdots=\frac{n!}{k_1!\cdots k_r!},
\]
porque el producto es telescópico. La segunda formulación es equivalente: la sucesión indica en qué caja va cada objeto.
\end{proof}

\begin{ejemplo}
La palabra \texttt{BANANA} tiene 6 letras: una \texttt{B}, tres \texttt{A} y dos \texttt{N}. Tiene $\binom{6}{1,3,2}=\frac{720}{1\cdot6\cdot2}=60$ anagramas distintos. Este recuento es el de la «regla del contable» de~\cite[§14.6.2]{mcs}.
\end{ejemplo}

\begin{teorema}[Multinomio de Newton]\label{teo:t2-multinomio}
Para reales $x_1,\dots,x_r$ y $n\ge0$,
\[
(x_1+\dots+x_r)^{n}=\sum_{\substack{k_1+\dots+k_r=n\\ k_i\ge0}}\binom{n}{k_1,\dots,k_r}x_1^{k_1}\cdots x_r^{k_r}.
\]
\end{teorema}

\begin{proof}
Como en el binomio: el desarrollo es la suma de las $r^{n}$ palabras de longitud $n$ sobre $\{x_1,\dots,x_r\}$, y las que contienen $k_i$ veces $x_i$ son $\binom{n}{k_1,\dots,k_r}$ por el \cref{teo:t2-multinomial}.
\end{proof}

\section{Combinaciones con repetición}

\begin{definicion}
Una \emph{combinación con repetición} de orden $k$ de un conjunto de $n$ elementos es una selección no ordenada de $k$ elementos en la que se permite repetir; equivale a un vector $(x_1,\dots,x_n)$ de enteros $x_i\ge 0$ con $x_1+\dots+x_n=k$, donde $x_i$ es el número de veces que se elige el elemento $i$.
\end{definicion}

\begin{teorema}[Estrellas y barras]\label{teo:t2-estrellas}
El número de soluciones enteras de $x_1+\dots+x_n=k$ con todos los $x_i\ge0$ es $\binom{n+k-1}{k}=\binom{n+k-1}{n-1}$. El número de soluciones con todos los $x_i\ge1$ es $\binom{k-1}{n-1}$ (para $k\ge n\ge 1$).
\end{teorema}

\begin{proof}
A cada solución $(x_1,\dots,x_n)$ se le asocia la cadena formada por $x_1$ estrellas, una barra, $x_2$ estrellas, una barra, \dots, y $x_n$ estrellas: tiene $k$ estrellas y $n-1$ barras. Recíprocamente, toda cadena de $k$ estrellas y $n-1$ barras determina una única solución (los $x_i$ son los tamaños de los bloques de estrellas separados por las barras). Es una biyección, y las cadenas se cuentan eligiendo las posiciones de las $n-1$ barras entre $n+k-1$: $\binom{n+k-1}{n-1}$. Para la segunda afirmación, el cambio $y_i=x_i-1$ transforma biyectivamente las soluciones con $x_i\ge1$ en las soluciones de $y_1+\dots+y_n=k-n$ con $y_i\ge0$, que son $\binom{k-1}{n-1}$.
\end{proof}

\begin{table}[h]
\centering
\begin{tabular}{@{}lcc@{}}
\toprule
Selección de $k$ entre $n$ & sin repetición & con repetición \\
\midrule
ordenada (importa el orden) & $n^{\underline{k}}=\dfrac{n!}{(n-k)!}$ & $n^{k}$ \\[8pt]
no ordenada & $\dbinom{n}{k}$ & $\dbinom{n+k-1}{k}$ \\
\bottomrule
\end{tabular}
\caption{Los cuatro modelos básicos de selección. El caso con repetición y sin orden es el único en que la fórmula no es «evidente».}
\label{tab:t2-modelos}
\end{table}

\section{Doble conteo y demostraciones biyectivas}

Una \emph{demostración por doble conteo} establece una identidad $a=b$ mostrando que ambos lados cuentan el mismo conjunto de dos maneras distintas; una \emph{demostración biyectiva} la establece construyendo una biyección entre un conjunto de $a$ elementos y otro de $b$ elementos. Ya se han usado ambas en la \cref{prop:t2-props}. El esquema conviene escribirlo siempre de forma explícita: (1) describir el conjunto que se cuenta (una «historia»); (2) contarlo de una forma, obteniendo el lado izquierdo; (3) contarlo de otra, obteniendo el lado derecho~\cite[§1.5]{bh}.

\begin{ejemplo}[Caminos en una cuadrícula]\label{ej:t2-caminos}
Un camino reticular de $(0,0)$ a $(a,b)$, con $a,b\ge 0$ enteros, es una sucesión de pasos unitarios hacia la derecha ($D$) o hacia arriba ($U$). Todo camino tiene $a$ pasos $D$ y $b$ pasos $U$, y queda determinado por la posición de los pasos $D$ entre los $a+b$; hay $\binom{a+b}{a}$. Clasificando los caminos de $(0,0)$ a $(a,b)$ según su último paso se obtiene de nuevo la fórmula de Pascal: $\binom{a+b}{a}=\binom{a+b-1}{a-1}+\binom{a+b-1}{a}$.
\end{ejemplo}

\begin{proposicion}[Identidad de Vandermonde]\label{prop:t2-vandermonde}
Para enteros $m,n,k\ge0$, $\displaystyle\binom{m+n}{k}=\sum_{j=0}^{k}\binom{m}{j}\binom{n}{k-j}$.
\end{proposicion}

La demostración se pide en el \cref{eje:2-vandermonde}: un grupo de $m$ mujeres y $n$ hombres, y comités de $k$ personas clasificados por el número de mujeres.

\section{Inclusión--exclusión general}

\begin{teorema}[Inclusión--exclusión]\label{teo:t2-ie}
Para conjuntos finitos $A_1,\dots,A_m$,
\[
\Bigl|\bigcup_{i=1}^{m}A_i\Bigr|=\sum_{j=1}^{m}(-1)^{j+1}\sum_{1\le i_1<\dots<i_j\le m}\abs{A_{i_1}\cap\dots\cap A_{i_j}}.
\]
\end{teorema}

\begin{proof}
Se comprueba que cada elemento $x$ de la unión contribuye exactamente~1 al lado derecho (los que no están en la unión contribuyen~0 a ambos lados). Sea $r\ge1$ el número de conjuntos $A_i$ que contienen a $x$. Una intersección de $j$ conjuntos contiene a $x$ si y solo si los $j$ conjuntos están entre esos $r$, lo que ocurre para $\binom{r}{j}$ de las intersecciones de $j$ conjuntos. La contribución de $x$ es, pues,
\[
\sum_{j=1}^{r}(-1)^{j+1}\binom{r}{j}=1-\sum_{j=0}^{r}(-1)^{j}\binom{r}{j}=1-0=1,
\]
por el \cref{cor:t2-sumas}.
\end{proof}

\begin{ejemplo}[Desarreglos]\label{ej:t2-desarreglos}
Un \emph{desarreglo} de $[n]$ es una permutación $\sigma$ sin puntos fijos: $\sigma(i)\neq i$ para todo $i$. Sea $A_i$ el conjunto de permutaciones con $\sigma(i)=i$. Una intersección de $j$ de ellos fija $j$ puntos y permuta libremente los demás, luego tiene $(n-j)!$ elementos. Por el \cref{teo:t2-ie}, el número de permutaciones con algún punto fijo es $\sum_{j=1}^{n}(-1)^{j+1}\binom{n}{j}(n-j)!$, y el de desarreglos es
\[
D_n=n!-\sum_{j=1}^{n}(-1)^{j+1}\frac{n!}{j!}=n!\sum_{j=0}^{n}\frac{(-1)^{j}}{j!}.
\]
Como $\sum_{j\ge0}(-1)^{j}/j!=e^{-1}$, la proporción de desarreglos $D_n/n!$ tiende a $1/e\approx0{,}368$. Por ejemplo, $D_4=24\bigl(1-1+\tfrac12-\tfrac16+\tfrac1{24}\bigr)=9$.
\end{ejemplo}

\begin{observacion}[Conexión con el curso]
En el módulo~22, la métrica \emph{pass@$k$} para evaluar código generado se estima a partir de $n$ muestras de las que $c$ son correctas: es la probabilidad de que un subconjunto aleatorio de $k$ muestras contenga alguna correcta, $1-\binom{n-c}{k}\big/\binom{n}{k}$, un cociente de números combinatorios~\cite[§2.1]{chen}. En el módulo~1, la verosimilitud de un corpus depende solo de los recuentos de cada token, y el número de secuencias con unos recuentos dados es un coeficiente multinomial. En el módulo~12, el algoritmo BPE cuenta pares de símbolos adyacentes: una secuencia de longitud $L$ tiene $L-1$ pares adyacentes, mientras que un conjunto de $K$ símbolos admite $K^{2}$ pares ordenados posibles.
\end{observacion}

\begin{observacion}[Errores frecuentes]
(1) Usar $\binom{n}{k}$ cuando importa el orden, o $n^{\underline{k}}$ cuando no importa. (2) Olvidar dividir por el número de ordenaciones de grupos \emph{indistinguibles}: repartir 12 personas en tres equipos de 4 sin nombre no es lo mismo que en los equipos A, B y C (\cref{eje:2-equipos}). (3) Contar «al menos un par» eligiendo el par y luego el resto libremente, lo que repite manos. (4) Confundir combinaciones con repetición con variaciones con repetición: el número de multiconjuntos de tamaño $k$ \emph{no} es $n^{k}/k!$. (5) En una demostración por doble conteo, no decir explícitamente qué conjunto se cuenta.
\end{observacion}

\begin{lecturas}
\item Lehman, Leighton y Meyer~\cite{mcs}, §14.5 «Counting Subsets», §14.6 «Sequences with Repetitions» y §14.7 «Counting Practice: Poker Hands».
\item Ross~\cite{ross}, §1.4 «Combinations», §1.5 «Multinomial Coefficients» y §1.6 «The Number of Integer Solutions of Equations».
\item Blitzstein y Hwang~\cite{bh}, §1.4 «How to count» y §1.5 «Story proofs».
\item Grinstead y Snell~\cite{gs}, §3.2 «Combinations».
\end{lecturas}

\section*{Ejercicios del tema 2}
\addcontentsline{toc}{section}{Ejercicios del tema 2}

\begin{ejercicio}{1}{eje:2-calcular}
Calcule $\binom{10}{3}$, $\binom{10}{7}$, $\binom{52}{5}$ y el coeficiente de $x^{3}$ en $(1+x)^{7}$. Explique sin hacer cuentas por qué $\binom{10}{3}=\binom{10}{7}$.
\end{ejercicio}

\begin{ejercicio}{1}{eje:2-modelo}
Identifique el modelo de la \cref{tab:t2-modelos} y cuente: (a) podios (oro, plata, bronce) entre 8 atletas; (b) comisiones de 3 personas entre 8; (c) secuencias de 3 tokens de un vocabulario de 8; (d) bolsas de 3 caramelos de 8 sabores, con sabores posiblemente repetidos.
\end{ejercicio}

\begin{ejercicio}{2}{eje:2-poker}
Una mano de póquer es un subconjunto de 5 cartas de una baraja francesa (13 valores, 4 palos). Cuente las manos con (a) póquer (cuatro cartas del mismo valor); (b) full (tres de un valor y dos de otro); (c) doble pareja (dos parejas de valores distintos y una quinta carta de un tercer valor). Exprese cada resultado como fracción del total de manos.
\fuente{adaptado de Lehman, Leighton y Meyer~\cite{mcs}, §14.7 «Counting Practice: Poker Hands».}
\end{ejercicio}

\begin{ejercicio}{2}{eje:2-equipos}
(a) ¿De cuántas formas pueden repartirse 12 personas en tres equipos de tamaños 2, 5 y 5 (los equipos no tienen nombre)? (b) ¿Y en tres equipos de 4?
\fuente{adaptado de Harvard Stat 110~\cite{stat110}, \emph{Homework}~1, problema~2.}
\end{ejercicio}

\begin{ejercicio}{2}{eje:2-anagramas}
Un anagrama de una palabra es una ordenación distinta de sus letras. Cuente los anagramas de \texttt{MISSISSIPPI} y los de \texttt{PROBABILIDAD}. ¿En cuántos anagramas de \texttt{MISSISSIPPI} aparecen las cuatro \texttt{S} juntas?
\end{ejercicio}

\begin{ejercicio}{2}{eje:2-coeficientes}
Halle el coeficiente de $x^{3}y^{5}$ en $(2x-y)^{8}$ y el de $x^{2}y^{3}z$ en $(x+y+z)^{6}$. ¿Cuántos términos distintos tiene el desarrollo de $(x+y+z)^{6}$?
\end{ejercicio}

\begin{ejercicio}{2}{eje:2-ecuaciones}
Cuente las soluciones enteras de $x_1+x_2+x_3+x_4=10$ (a) con $x_i\ge0$; (b) con $x_i\ge1$; (c) con $x_i\ge0$ y $x_1\ge2$.
\end{ejercicio}

\begin{ejercicio}{2}{eje:2-caminos}
Cuente los caminos reticulares (\cref{ej:t2-caminos}) (a) de $(0,0)$ a $(5,6)$; (b) de $(0,0)$ a $(5,6)$ que pasan por $(2,3)$; (c) de $(0,0)$ a $(5,6)$ que \emph{no} pasan por $(2,3)$.
\fuente{adaptado de Harvard Stat 110~\cite{stat110}, \emph{Strategic Practice}~1, problema~3 (con otros datos).}
\end{ejercicio}

\begin{ejercicio}{2}{eje:2-estaciones}
Se anota la estación del año (primavera, verano, otoño o invierno) en que nació cada una de 7 personas. ¿Cuántas de las $4^{7}$ asignaciones posibles de estaciones incluyen las cuatro estaciones? ¿Qué fracción del total representan?
\fuente{adaptado de Harvard Stat 110~\cite{stat110}, \emph{Strategic Practice}~2, sección «Inclusion-Exclusion», problema~1.}
\end{ejercicio}

\begin{ejercicio}{2}{eje:2-desarreglos}
Cinco amigos organizan un amigo invisible: cada uno saca al azar el nombre de otro de una bolsa con los cinco nombres. (a) ¿Cuántos sorteos hay en los que nadie se saca a sí mismo? (b) ¿Qué fracción de los $5!$ sorteos representan? Compárela con $1/e$.
\end{ejercicio}

\begin{ejercicio}{2}{eje:2-ositos}
Una bolsa de gominolas contiene entre 30 y 50 ositos, cada uno de uno de 5 sabores. Dos bolsas se consideran iguales si tienen el mismo número de ositos de cada sabor. ¿Cuántas bolsas distintas hay? Exprese la respuesta como diferencia de dos números combinatorios, no como una suma larga.
\fuente{adaptado de Harvard Stat 110~\cite{stat110}, \emph{Homework}~1, problema~7(b).}
\end{ejercicio}

\begin{ejercicio}{3}{eje:2-palo}
Demuestre mediante una historia (doble conteo) que, para enteros $n\ge k\ge0$,
\[
\binom{k}{k}+\binom{k+1}{k}+\dots+\binom{n}{k}=\binom{n+1}{k+1}.
\]
\emph{Indicación:} ordene $n+1$ personas por edad y clasifique los grupos de $k+1$ según quién es el mayor del grupo.
\fuente{adaptado de Harvard Stat 110~\cite{stat110}, \emph{Homework}~1, problema~7(a).}
\end{ejercicio}

\begin{ejercicio}{3}{eje:2-sumak}
(a) Usando la absorción (\cref{prop:t2-props}), demuestre que $\sum_{k=0}^{n}k\binom{n}{k}=n2^{n-1}$ para $n\ge1$. (b) Demuestre que, para $n\ge2$, $\sum_{k=2}^{n}k(k-1)\binom{n}{k}=n(n-1)2^{n-2}$, contando de dos maneras los comités con un presidente y un vicepresidente distintos.
\fuente{el apartado (b) está adaptado de MIT 6.041~\cite{mit6041f10}, \emph{Problem Set}~3, problema~5.}
\end{ejercicio}

\begin{ejercicio}{3}{eje:2-vandermonde}
Demuestre la identidad de Vandermonde (\cref{prop:t2-vandermonde}) por doble conteo y deduzca que $\sum_{j=0}^{n}\binom{n}{j}^{2}=\binom{2n}{n}$.
\end{ejercicio}

\begin{ejercicio}{3}{eje:2-pares}
Sea $A$ un conjunto con $n\ge1$ elementos y $a\in A$ fijo. (a) Demuestre que la aplicación $X\mapsto X\dsim\{a\}$ (añadir $a$ si no está, quitarlo si está) es una biyección entre los subconjuntos de $A$ de cardinal par y los de cardinal impar. (b) Deduzca, sin usar el binomio, que $\sum_{k}(-1)^{k}\binom{n}{k}=0$ y que $A$ tiene $2^{n-1}$ subconjuntos de cardinal par.
\end{ejercicio}

% =====================================================================
\chapter{Espacios de probabilidad}
\label{tema:3}
% =====================================================================

Este tema da la definición axiomática de probabilidad, deduce de los axiomas las reglas de cálculo y estudia el caso más importante para empezar, el modelo de Laplace, en el que calcular una probabilidad se reduce a contar. Se siguen la sección §1.2 de Bertsekas y Tsitsiklis~\cite{bt} (clase~1 de MIT 6.041~\cite{mit6041f10}), el capítulo~2 de Ross~\cite{ross} (§2.2 a §2.5), las secciones §1.2, §1.3 y §1.6 de Blitzstein y Hwang~\cite{bh} y el capítulo~16, «Events and Probability Spaces», de~\cite{mcs}, que incluye el problema del cumpleaños (§16.4).

\section{Experimentos, espacio muestral y sucesos}

\begin{definicion}[Espacio muestral y sucesos]
Un \emph{experimento aleatorio} es un procedimiento cuyo resultado no puede predecirse con certeza. El \emph{espacio muestral} $\Omega$ es el conjunto de todos sus resultados posibles. Un \emph{suceso} es un subconjunto $A\subseteq\Omega$; se dice que $A$ \emph{ocurre} si el resultado obtenido pertenece a $A$. Los sucesos $\{\omega\}$ se llaman \emph{elementales}.
\end{definicion}

Las operaciones con conjuntos del módulo~0.2 se traducen en operaciones con sucesos: $A\cup B$ es «ocurre $A$ o $B$»; $A\cap B$, «ocurren $A$ y $B$»; $A^{c}=\Omega\setminus A$, «no ocurre $A$»; $A\subseteq B$, «si ocurre $A$, ocurre $B$». Dos sucesos son \emph{incompatibles} (o mutuamente excluyentes) si $A\cap B=\emptyset$. Las leyes de De Morgan dan, por ejemplo, $(A\cup B)^{c}=A^{c}\cap B^{c}$: «no ocurre ninguno de los dos».

\begin{ejemplo}\label{ej:t3-omega}
Al lanzar dos dados distinguibles, $\Omega=\{(i,j):i,j\in[6]\}$ tiene 36 resultados. El suceso «la suma es 7» es $\{(1,6),(2,5),(3,4),(4,3),(5,2),(6,1)\}$. Un modelo de lenguaje que genera dos tokens de un vocabulario $\V$ tiene espacio muestral $\V^{2}$, y «el primer token es \texttt{el}» es el suceso $\{\texttt{el}\}\times\V$.
\end{ejemplo}

\section{Axiomas de Kolmogorov}

\begin{definicion}[Espacio de probabilidad]\label{def:t3-axiomas}
Una \emph{ley de probabilidad} sobre $\Omega$ es una función $\PP$ que asigna a cada suceso $A$ un número real $\PP(A)$ y que cumple:
\begin{enumerate}
  \item \emph{no negatividad}: $\PP(A)\ge0$ para todo suceso $A$;
  \item \emph{normalización}: $\PP(\Omega)=1$;
  \item \emph{aditividad numerable}: si $A_1,A_2,\dots$ son sucesos incompatibles dos a dos, $\PP\bigl(\bigcup_{i\ge1}A_i\bigr)=\sum_{i\ge1}\PP(A_i)$.
\end{enumerate}
El par $(\Omega,\PP)$ se llama \emph{espacio de probabilidad}.
\end{definicion}

\begin{observacion}
Si $\Omega$ es finito o numerable, todo subconjunto puede ser un suceso. Si $\Omega$ es no numerable, como un intervalo de $\R$, no es posible en general asignar probabilidad a todos los subconjuntos de forma coherente, y los sucesos se restringen a una familia adecuada (una $\sigma$-álgebra). En este módulo esa sutileza no interviene: todos los sucesos que aparecen (intervalos, uniones de intervalos, regiones del plano) son admisibles~\cite[§1.6]{bh}.
\end{observacion}

\section{Propiedades que se deducen de los axiomas}

\begin{proposicion}\label{prop:t3-props}
Sean $A$, $B$ y $A_1,\dots,A_n$ sucesos. Entonces:
\begin{enumerate}
  \item $\PP(\emptyset)=0$;
  \item \emph{aditividad finita}: si $A_1,\dots,A_n$ son incompatibles dos a dos, $\PP\bigl(\bigcup_{i=1}^{n}A_i\bigr)=\sum_{i=1}^{n}\PP(A_i)$;
  \item $\PP(A^{c})=1-\PP(A)$;
  \item si $A\subseteq B$, entonces $\PP(B\setminus A)=\PP(B)-\PP(A)$ y $\PP(A)\le\PP(B)$;
  \item $0\le\PP(A)\le1$;
  \item $\PP(A\cup B)=\PP(A)+\PP(B)-\PP(A\cap B)$.
\end{enumerate}
\end{proposicion}

\begin{proof}
(i) Los sucesos $A_1=\Omega$, $A_i=\emptyset$ ($i\ge2$) son incompatibles dos a dos y su unión es $\Omega$. Por el axioma (iii), $1=\PP(\Omega)=1+\sum_{i\ge2}\PP(\emptyset)$, luego $\sum_{i\ge2}\PP(\emptyset)=0$; como los sumandos son iguales y no negativos, $\PP(\emptyset)=0$.

(ii) Se completa la lista con $A_i=\emptyset$ para $i>n$ y se aplican el axioma (iii) y el apartado~(i).

(iii) $A$ y $A^{c}$ son incompatibles y su unión es $\Omega$: $\PP(A)+\PP(A^{c})=1$.

(iv) Si $A\subseteq B$, $B$ es la unión disjunta de $A$ y $B\setminus A$; por (ii), $\PP(B)=\PP(A)+\PP(B\setminus A)$, y como $\PP(B\setminus A)\ge0$, $\PP(A)\le\PP(B)$.

(v) Por (iv) con $A\subseteq\Omega$, $\PP(A)\le\PP(\Omega)=1$.

(vi) $A\cup B$ es la unión disjunta de $A$ y $B\setminus(A\cap B)$, y $A\cap B\subseteq B$. Por (ii) y (iv), $\PP(A\cup B)=\PP(A)+\PP(B)-\PP(A\cap B)$.
\end{proof}

\begin{teorema}[Desigualdad de Boole o cota de la unión]\label{teo:t3-boole}
Para sucesos cualesquiera $A_1,\dots,A_n$,
\(
\PP\bigl(\bigcup_{i=1}^{n}A_i\bigr)\le\sum_{i=1}^{n}\PP(A_i).
\)
\end{teorema}

\begin{proof}
Por inducción sobre $n$. El caso $n=1$ es trivial. Si es cierto para $n-1$, por la \cref{prop:t3-props}(vi) aplicada a $B=A_1\cup\dots\cup A_{n-1}$ y $A_n$,
\[
\PP(B\cup A_n)=\PP(B)+\PP(A_n)-\PP(B\cap A_n)\le\PP(B)+\PP(A_n)\le\sum_{i=1}^{n-1}\PP(A_i)+\PP(A_n).\qedhere
\]
\end{proof}

La fórmula de inclusión--exclusión del \cref{teo:t2-ie} tiene una versión probabilística con la misma estructura; para tres sucesos,
\[
\PP(A\cup B\cup C)=\PP(A)+\PP(B)+\PP(C)-\PP(A\cap B)-\PP(A\cap C)-\PP(B\cap C)+\PP(A\cap B\cap C),
\]
y se demuestra a partir de (vi) exactamente como el \cref{teo:t1-ie}.

\section{Espacios discretos y modelo de Laplace}

\begin{proposicion}[Espacios discretos]\label{prop:t3-discreto}
Sea $\Omega$ finito o numerable y $p\colon\Omega\to[0,1]$ con $\sum_{\omega\in\Omega}p(\omega)=1$. Entonces $\PP(A)=\sum_{\omega\in A}p(\omega)$ define una ley de probabilidad. Recíprocamente, toda ley de probabilidad sobre $\Omega$ es de esta forma, con $p(\omega)=\PP(\{\omega\})$.
\end{proposicion}

\begin{proof}
La no negatividad y la normalización son inmediatas. Si $A_1,A_2,\dots$ son disjuntos, la suma de $p$ sobre $\bigcup_iA_i$ puede reordenarse agrupando por conjuntos, porque es una serie de términos no negativos, y da $\sum_i\PP(A_i)$. Para el recíproco, todo $A\subseteq\Omega$ es la unión disjunta, finita o numerable, de los sucesos elementales $\{\omega\}$ con $\omega\in A$; por el axioma~(iii), $\PP(A)=\sum_{\omega\in A}\PP(\{\omega\})$.
\end{proof}

\begin{definicion}[Modelo de Laplace]\label{def:t3-laplace}
Si $\Omega$ es finito y todos los resultados son equiprobables, $p(\omega)=1/\abs{\Omega}$, y entonces
\[
\PP(A)=\frac{\abs{A}}{\abs{\Omega}}=\frac{\text{casos favorables}}{\text{casos posibles}}.
\]
\end{definicion}

El modelo de Laplace es una \emph{hipótesis} sobre el experimento, no una consecuencia de los axiomas: es razonable para dados, monedas y cartas bien mezcladas, pero no, por ejemplo, para la suma de dos dados, cuyos 11 valores no son equiprobables. La técnica habitual consiste en elegir un espacio muestral en el que la equiprobabilidad sea evidente, normalmente uno de resultados \emph{ordenados} y distinguibles, aunque los objetos del problema parezcan indistinguibles.

\begin{ejemplo}[Ordenado o no ordenado]\label{ej:t3-ases}
Se extraen dos cartas de una baraja francesa sin reemplazamiento. Con resultados ordenados, $\abs{\Omega}=52\cdot51$ y los pares de ases son $4\cdot3$, luego $\PP(\text{dos ases})=\frac{12}{2652}=\frac{1}{221}$. Con resultados no ordenados, $\abs{\Omega}=\binom{52}{2}=1326$ y los favorables son $\binom{4}{2}=6$: de nuevo $\frac{6}{1326}=\frac1{221}$. Ambos modelos son válidos porque cada subconjunto corresponde exactamente a dos pares ordenados; lo que no es válido es mezclar recuentos ordenados en el numerador con no ordenados en el denominador.
\end{ejemplo}

\section{El problema del cumpleaños}

\begin{proposicion}\label{prop:t3-cumple}
Se asigna a cada una de $n$ personas un día entre $N$, de forma uniforme e independiente (todas las $N^{n}$ asignaciones son equiprobables). La probabilidad de que dos o más personas compartan día es
\[
p_n=1-\frac{N^{\underline{n}}}{N^{n}}=1-\prod_{i=0}^{n-1}\Bigl(1-\frac{i}{N}\Bigr),
\]
y además $1-p_n\le\exp\!\bigl(-\tfrac{n(n-1)}{2N}\bigr)$.
\end{proposicion}

\begin{proof}
El suceso complementario, «todos los días son distintos», está formado por las asignaciones inyectivas, que son $N^{\underline{n}}$ (\cref{teo:t1-variaciones}); se aplica el modelo de Laplace y la regla del complemento. Para la cota se usa la desigualdad $1-x\le e^{-x}$, válida para todo $x$ real (la función $g(x)=e^{-x}-1+x$ tiene $g'(x)=1-e^{-x}$, que se anula solo en $x=0$, donde $g$ alcanza su mínimo $g(0)=0$). Entonces
\[
1-p_n=\prod_{i=0}^{n-1}\Bigl(1-\frac iN\Bigr)\le\prod_{i=0}^{n-1}e^{-i/N}=\exp\Bigl(-\frac1N\sum_{i=0}^{n-1}i\Bigr)=\exp\Bigl(-\frac{n(n-1)}{2N}\Bigr).\qedhere
\]
\end{proof}

\begin{ejemplo}
Con $N=365$ y $n=23$, el producto da $1-p_{23}\approx0{,}4927$, luego $p_{23}\approx0{,}5073$: en un grupo de 23 personas es más probable que haya dos con el mismo cumpleaños que lo contrario. La cota da $1-p_{23}\le e^{-253/365}\approx0{,}49999$, suficiente para asegurar $p_{23}>1/2$. La razón de que baste un número tan pequeño es que lo que importa es el número de \emph{pares}, $\binom{23}{2}=253$, no el de personas~\cite[§16.4]{mcs}.
\end{ejemplo}

\begin{figure}[h]
\centering
\begin{tikzpicture}[x=0.11cm,y=3.2cm]
  \draw[->] (0,0) -- (63,0) node[right] {\small $n$};
  \draw[->] (0,0) -- (0,1.12) node[above] {\small $p_n$};
  \foreach \x in {10,20,30,40,50,60} \draw (\x,0.02) -- (\x,-0.02) node[below] {\scriptsize \x};
  \foreach \y in {0.5,1} \draw (0.6,\y) -- (-0.6,\y) node[left] {\scriptsize \y};
  \draw[dashed,gray] (0,0.5) -- (60,0.5);
  \draw[dashed,gray] (23,0) -- (23,0.5073);
  \draw[thick] plot[smooth] coordinates {(1,0) (5,0.0271) (10,0.1169) (15,0.2529) (20,0.4114) (23,0.5073) (25,0.5687) (30,0.7063) (35,0.8144) (40,0.8912) (45,0.9410) (50,0.9704) (55,0.9863) (60,0.9941)};
  \fill (23,0.5073) circle (0.9pt);
\end{tikzpicture}
\caption{Probabilidad $p_n$ de que al menos dos de $n$ personas compartan cumpleaños ($N=365$).}
\label{fig:t3-cumple}
\end{figure}

\begin{observacion}[Conexión con el curso]
El problema del cumpleaños es el modelo de las \emph{colisiones} de funciones \emph{hash}: si $n$ documentos se asignan uniformemente a $N$ valores, la probabilidad de colisión es $p_n$, y la cota $1-p_n\le e^{-n(n-1)/(2N)}$ muestra que las colisiones aparecen cuando $n$ es del orden de $\sqrt{N}$. Este razonamiento interviene en la deduplicación de corpus con MinHash del módulo~16. El espacio muestral $\V^{T}$ de secuencias de tokens y la regla del complemento son la base del módulo~1, y la cota de la unión reaparece cuando se acota la probabilidad de que falle alguno de muchos sucesos.
\end{observacion}

\begin{observacion}[Errores frecuentes]
(1) Suponer equiprobabilidad donde no la hay, como en las sumas de dos dados. (2) Mezclar recuentos ordenados y no ordenados en el mismo cociente (\cref{ej:t3-ases}). (3) Sumar probabilidades de sucesos que no son incompatibles: $\PP(A\cup B)=\PP(A)+\PP(B)$ solo si $A\cap B=\emptyset$; en general, la suma es solo una cota superior (\cref{teo:t3-boole}). (4) Tratar como indistinguibles objetos físicamente distintos (dos dados iguales siguen siendo dos dados). (5) Confundir «alguien comparte cumpleaños con usted» (\cref{eje:3-cumple}) con «dos personas cualesquiera comparten cumpleaños».
\end{observacion}

\begin{lecturas}
\item Bertsekas y Tsitsiklis~\cite{bt}, §1.1 y §1.2 (clase~1 de MIT 6.041~\cite{mit6041f10}, «Probability models and axioms»).
\item Ross~\cite{ross}, cap.~2 «Axioms of Probability», §2.2 a §2.5.
\item Blitzstein y Hwang~\cite{bh}, §1.2 «Sample spaces and Pebble World», §1.3 «Naive definition of probability» y §1.6 «Non-naive definition of probability».
\item Lehman, Leighton y Meyer~\cite{mcs}, cap.~16 «Events and Probability Spaces», §16.4 «The Birthday Principle» y §16.5 «Set Theory and Probability».
\item Grinstead y Snell~\cite{gs}, §1.2 «Discrete Probability Distributions».
\end{lecturas}

\section*{Ejercicios del tema 3}
\addcontentsline{toc}{section}{Ejercicios del tema 3}

\begin{ejercicio}{1}{eje:3-espacio}
Describa un espacio muestral adecuado y su cardinal para: (a) tres lanzamientos de una moneda; (b) una mano de 5 cartas; (c) la generación de dos tokens con un vocabulario de $K$ tokens; (d) el número de lanzamientos de una moneda hasta la primera cara. ¿En cuáles es razonable el modelo de Laplace?
\end{ejercicio}

\begin{ejercicio}{1}{eje:3-sucesos}
Sean $A$, $B$ y $C$ sucesos. Exprese con uniones, intersecciones y complementarios: (a) ocurre al menos uno; (b) ocurren los tres; (c) no ocurre ninguno; (d) ocurre exactamente uno; (e) ocurre como máximo uno; (f) ocurren $A$ y $B$, pero no $C$.
\fuente{adaptado de MIT 6.041~\cite{mit6041f10}, \emph{Problem Set}~1, problema~1.}
\end{ejercicio}

\begin{ejercicio}{2}{eje:3-monedas}
Se lanza tres veces una moneda equilibrada. Calcule la probabilidad de (a) obtener CCC; (b) obtener CXC; (c) obtener exactamente dos caras; (d) obtener al menos tantas caras como cruces.
\fuente{adaptado de MIT 6.041~\cite{mit6041f10}, \emph{Problem Set}~1, problema~2.}
\end{ejercicio}

\begin{ejercicio}{2}{eje:3-dados}
Se lanzan dos dados equilibrados. Calcule la probabilidad de que (a) la suma sea 7; (b) salga al menos un 6; (c) el máximo de los dos valores sea~4. En (b), compare con el resultado erróneo $\frac16+\frac16$ y explique el error.
\end{ejercicio}

\begin{ejercicio}{2}{eje:3-cuatro-dados}
Se lanzan cuatro dados equilibrados. Sin calcular ninguna de las dos probabilidades por separado, decida cuál es mayor, la de que la suma sea 21 o la de que sea 22. Después calcule ambas.
\fuente{adaptado de Harvard Stat 110~\cite{stat110}, \emph{Strategic Practice}~1, problema~1(a).}
\end{ejercicio}

\begin{ejercicio}{2}{eje:3-horarios}
Una facultad asigna cada asignatura, al azar y de forma independiente, a una de 10 franjas horarias. Un estudiante elige 3 asignaturas. ¿Qué probabilidad hay de que dos o más coincidan en la misma franja?
\fuente{adaptado de Harvard Stat 110~\cite{stat110}, \emph{Homework}~1, problema~3.}
\end{ejercicio}

\begin{ejercicio}{2}{eje:3-robos}
En una ciudad con 6 distritos se producen 6 robos; cada uno ocurre en un distrito elegido al azar, independientemente de los demás. ¿Qué probabilidad hay de que algún distrito haya sufrido más de un robo?
\fuente{adaptado de Harvard Stat 110~\cite{stat110}, \emph{Homework}~1, problema~4.}
\end{ejercicio}

\begin{ejercicio}{2}{eje:3-seis-cartas}
Se reparte una mano de 6 cartas de una baraja francesa. Se llama \emph{doble pareja} a una mano con exactamente dos valores repetidos dos veces y las otras dos cartas de valores distintos entre sí y de los anteriores, y \emph{trío} a una mano con un valor repetido tres veces y las otras tres cartas de valores distintos entre sí y del anterior. ¿Cuál es más probable? Calcule ambas probabilidades.
\fuente{adaptado de MIT 18.05~\cite{mit1805}, \emph{Problem Set}~1, problema~1.}
\end{ejercicio}

\begin{ejercicio}{2}{eje:3-cumple}
Hay $n$ personas, además de usted, con cumpleaños uniformes e independientes entre 365 días. (a) Escriba la probabilidad de que alguna comparta cumpleaños con usted y halle el menor $n$ para el que supera $1/2$. (b) Explique por qué ese $n$ es mucho mayor que 23 y por qué es mayor que $365/2$.
\fuente{adaptado de MIT 18.05~\cite{mit1805}, \emph{Problem Set}~1, problema~3, apartado~(c); el apartado~(b) es nuevo.}
\end{ejercicio}

\begin{ejercicio}{2}{eje:3-urna}
Una urna contiene $r$ bolas rojas y $v$ verdes, $r,v\ge1$. Se extraen dos sin reemplazamiento. (a) Demuestre, con un espacio muestral de resultados ordenados, que la probabilidad de que la segunda sea verde es igual a la de que lo sea la primera. (b) Si $r+v=16$ y es igual de probable que las dos bolas sean del mismo color que de colores distintos, halle $r$ y $v$.
\fuente{adaptado de Harvard Stat 110~\cite{stat110}, \emph{Homework}~1, problema~6.}
\end{ejercicio}

\begin{ejercicio}{2}{eje:3-geometrica}
Se eligen dos números $x$ e $y$ al azar en $[0,1]$, de modo que la probabilidad de cada región del cuadrado $[0,1]^{2}$ es su área. Calcule $\PP(\abs{x-y}>1/3)$ y $\PP(x+y<1/2)$.
\end{ejercicio}

\begin{ejercicio}{3}{eje:3-bonferroni}
Demuestre que, para sucesos cualesquiera, $\PP(A\cap B)\ge\PP(A)+\PP(B)-1$ y, más en general, $\PP(A\cap B\cap C)\ge\PP(A)+\PP(B)+\PP(C)-2$.
\fuente{adaptado de MIT 6.041~\cite{mit6041f10}, \emph{Problem Set}~1, problema~6.}
\end{ejercicio}

\begin{ejercicio}{3}{eje:3-interseccion}
Demuestre que, para sucesos $A_1,\dots,A_n$, $\PP\bigl(\bigcap_{i=1}^{n}A_i\bigr)\ge1-\sum_{i=1}^{n}\PP(A_i^{c})$. Aplíquelo: si cada uno de 20 componentes falla con probabilidad $0{,}001$, sin ninguna hipótesis de independencia, ¿qué puede asegurarse de la probabilidad de que no falle ninguno?
\end{ejercicio}

\begin{ejercicio}{3}{eje:3-distancia}
Sea $d(A,B)=\PP(A\dsim B)$, donde $A\dsim B=(A\setminus B)\cup(B\setminus A)$. (a) Demuestre que $d(A,B)=\PP(A)+\PP(B)-2\PP(A\cap B)$. (b) Demuestre que $d(A,C)\le d(A,B)+d(B,C)$.
\end{ejercicio}

\begin{ejercicio}{3}{eje:3-hash}
Se asignan $n$ claves a $N$ casillas de una tabla \emph{hash}, de forma uniforme e independiente. (a) Usando la \cref{prop:t3-cumple}, demuestre que si $n(n-1)\ge 2N\ln 2$, la probabilidad de alguna colisión es al menos $1/2$. (b) Para $N=2^{32}$, halle el menor $n$ que cumple esa condición y compárelo con $\sqrt{N}$.
\end{ejercicio}

\begin{ejercicio}{3}{eje:3-laplace}
(a) Demuestre que el modelo de Laplace (\cref{def:t3-laplace}) cumple los tres axiomas de Kolmogorov. (b) Demuestre que no existe ninguna ley de probabilidad sobre $\N$ que asigne la misma probabilidad a todos los sucesos elementales $\{n\}$.
\end{ejercicio}

% =====================================================================
\chapter{Probabilidad condicionada, independencia y teorema de Bayes}
\label{tema:4}
% =====================================================================

La probabilidad condicionada formaliza cómo cambia la incertidumbre al recibir información. De ella se deducen la regla de la cadena, el teorema de la probabilidad total y el teorema de Bayes, y con ella se define la independencia. El tema sigue las secciones §1.3 a §1.5 de Bertsekas y Tsitsiklis~\cite{bt} (clases~2 y~3 de MIT 6.041~\cite{mit6041f10}), el capítulo~2 de Blitzstein y Hwang~\cite{bh}, el capítulo~3 de Ross~\cite{ross} y el capítulo~17 de~\cite{mcs}, cuyo tratamiento del problema de Monty Hall y de las pruebas médicas (§17.1 y §17.4.2) se adapta aquí.

\section{Probabilidad condicionada}

\begin{definicion}[Probabilidad condicionada]
Sean $A$ y $B$ sucesos con $\PP(B)>0$. La \emph{probabilidad de $A$ condicionada a $B$} es
\[
\PP(A\mid B)=\frac{\PP(A\cap B)}{\PP(B)}.
\]
\end{definicion}

La idea es que, al saber que ha ocurrido $B$, los resultados fuera de $B$ quedan descartados y los de $B$ conservan sus probabilidades \emph{relativas}; dividir por $\PP(B)$ restablece la normalización. En el modelo de Laplace, $\PP(A\mid B)=\abs{A\cap B}/\abs{B}$: se cuenta dentro de $B$ como si $B$ fuera el nuevo espacio muestral.

\begin{proposicion}\label{prop:t4-esprob}
Si $\PP(B)>0$, la función $A\mapsto\PP(A\mid B)$ es una ley de probabilidad sobre $\Omega$. En consecuencia, cumple todas las propiedades de la \cref{prop:t3-props}; por ejemplo, $\PP(A^{c}\mid B)=1-\PP(A\mid B)$.
\end{proposicion}

\begin{proof}
No negatividad: $\PP(A\cap B)\ge0$ y $\PP(B)>0$. Normalización: $\PP(\Omega\mid B)=\PP(B)/\PP(B)=1$. Aditividad: si $A_1,A_2,\dots$ son incompatibles, también lo son $A_i\cap B$, y $\bigl(\bigcup_iA_i\bigr)\cap B=\bigcup_i(A_i\cap B)$; por el axioma~(iii) aplicado a $\PP$, $\PP\bigl(\bigcup_iA_i\mid B\bigr)=\sum_i\PP(A_i\cap B)/\PP(B)=\sum_i\PP(A_i\mid B)$.
\end{proof}

\begin{ejemplo}\label{ej:t4-hijos}
Una familia tiene dos hijos; los cuatro casos (chico o chica para el mayor y para el menor) son equiprobables. Si se sabe que el mayor es chica, la probabilidad de que ambos lo sean es $\frac{1/4}{1/2}=\frac12$. Si solo se sabe que \emph{al menos uno} es chica, el suceso condicionante tiene tres casos y la probabilidad es $\frac{1/4}{3/4}=\frac13$. Informaciones de apariencia parecida condicionan a sucesos distintos.
\end{ejemplo}

\section{Regla de la cadena}

\begin{teorema}[Regla de la cadena o de la multiplicación]\label{teo:t4-cadena}
Si $A_1,\dots,A_n$ son sucesos con $\PP(A_1\cap\dots\cap A_{n-1})>0$, entonces
\[
\PP(A_1\cap A_2\cap\dots\cap A_n)=\PP(A_1)\,\PP(A_2\mid A_1)\,\PP(A_3\mid A_1\cap A_2)\cdots\PP(A_n\mid A_1\cap\dots\cap A_{n-1}).
\]
\end{teorema}

\begin{proof}
Todas las probabilidades condicionadas están definidas, porque $A_1\cap\dots\cap A_k\supseteq A_1\cap\dots\cap A_{n-1}$ tiene probabilidad positiva para $k\le n-1$. Por inducción sobre $n$: para $n=1$ no hay nada que probar, y si la fórmula vale para $n-1$, la definición de probabilidad condicionada con $B=A_1\cap\dots\cap A_{n-1}$ da $\PP(B\cap A_n)=\PP(B)\,\PP(A_n\mid B)$; basta sustituir $\PP(B)$ por la fórmula para $n-1$ sucesos.
\end{proof}

\begin{ejemplo}
Se extraen tres cartas sin reemplazamiento. Si $A_i$ es «la carta $i$ es de corazones», $\PP(A_1\cap A_2\cap A_3)=\frac{13}{52}\cdot\frac{12}{51}\cdot\frac{11}{50}\approx0{,}0129$. La regla de la cadena sustituye un recuento global por una sucesión de razonamientos locales, que es lo que se representa en un \emph{diagrama de árbol}: la probabilidad de cada hoja es el producto de las probabilidades de las ramas que llevan a ella~\cite[§16.2]{mcs}.
\end{ejemplo}

\section{Probabilidad total y teorema de Bayes}

\begin{definicion}
Una familia de sucesos $B_1,\dots,B_n$ es una \emph{partición} de $\Omega$ si son incompatibles dos a dos y $B_1\cup\dots\cup B_n=\Omega$.
\end{definicion}

\begin{teorema}[Probabilidad total]\label{teo:t4-total}
Si $B_1,\dots,B_n$ es una partición de $\Omega$ con $\PP(B_i)>0$ para todo $i$, entonces, para todo suceso $A$,
\[
\PP(A)=\sum_{i=1}^{n}\PP(A\mid B_i)\,\PP(B_i).
\]
\end{teorema}

\begin{proof}
$A=\bigcup_i(A\cap B_i)$ y los sucesos $A\cap B_i$ son incompatibles dos a dos. Por la aditividad, $\PP(A)=\sum_i\PP(A\cap B_i)$, y $\PP(A\cap B_i)=\PP(A\mid B_i)\PP(B_i)$ por definición.
\end{proof}

\begin{teorema}[Bayes]\label{teo:t4-bayes}
En las condiciones del \cref{teo:t4-total}, si además $\PP(A)>0$,
\[
\PP(B_j\mid A)=\frac{\PP(A\mid B_j)\,\PP(B_j)}{\sum_{i=1}^{n}\PP(A\mid B_i)\,\PP(B_i)}.
\]
\end{teorema}

\begin{proof}
$\PP(B_j\mid A)=\PP(A\cap B_j)/\PP(A)$; el numerador es $\PP(A\mid B_j)\PP(B_j)$ y el denominador se desarrolla con el \cref{teo:t4-total}.
\end{proof}

En la terminología bayesiana, $\PP(B_j)$ es la probabilidad \emph{a priori} de la hipótesis $B_j$, $\PP(A\mid B_j)$ es la \emph{verosimilitud} de la evidencia $A$ bajo esa hipótesis y $\PP(B_j\mid A)$ es la probabilidad \emph{a posteriori}. Con dos hipótesis $B$ y $B^{c}$, dividiendo las dos fórmulas de Bayes se obtiene la \emph{forma de razones} (\emph{odds}):
\[
\frac{\PP(B\mid A)}{\PP(B^{c}\mid A)}=\frac{\PP(B)}{\PP(B^{c})}\cdot\frac{\PP(A\mid B)}{\PP(A\mid B^{c})},
\]
«razón a posteriori = razón a priori $\times$ razón de verosimilitudes».

\begin{ejemplo}[Prueba diagnóstica]\label{ej:t4-diagnostico}
Una enfermedad afecta al $1\,\%$ de la población. Una prueba detecta el $90\,\%$ de los enfermos (\emph{sensibilidad}) y da positivo en el $5\,\%$ de los sanos (tasa de falsos positivos; la \emph{especificidad} es el $95\,\%$). Sean $E$ «estar enfermo» y $+$ «dar positivo». Por Bayes,
\[
\PP(E\mid +)=\frac{0{,}9\cdot0{,}01}{0{,}9\cdot0{,}01+0{,}05\cdot0{,}99}=\frac{0{,}009}{0{,}0585}=\frac{2}{13}\approx0{,}154.
\]
Pese a la aparente calidad de la prueba, solo un $15\,\%$ de los positivos están enfermos. En \emph{frecuencias naturales}: de cada 1000 personas, 10 están enfermas y 9 dan positivo; de las 990 sanas, unas 49,5 dan positivo; de unos 58,5 positivos, solo 9 son enfermos~\cite[§17.4.2 y §17.4.4]{mcs}.
\end{ejemplo}

\begin{ejemplo}[Monty Hall]\label{ej:t4-monty}
Un concursante elige una de tres puertas; tras una hay un coche y tras las otras dos, cabras. El presentador, que sabe dónde está el coche, abre siempre una de las otras dos puertas que oculte una cabra (al azar si hay dos) y ofrece cambiar. Sin pérdida de generalidad, el concursante elige la puerta~1. Sea $C_i$ «el coche está en la puerta $i$», con $\PP(C_i)=1/3$, y $M_3$ «el presentador abre la puerta~3». Entonces $\PP(M_3\mid C_1)=1/2$, $\PP(M_3\mid C_2)=1$ y $\PP(M_3\mid C_3)=0$. Por Bayes,
\[
\PP(C_2\mid M_3)=\frac{1\cdot\frac13}{\frac12\cdot\frac13+1\cdot\frac13+0\cdot\frac13}=\frac23,\qquad\PP(C_1\mid M_3)=\frac13.
\]
Cambiar de puerta duplica la probabilidad de ganar. La intuición «quedan dos puertas, luego $1/2$» ignora que la elección del presentador \emph{depende} de dónde está el coche y, por tanto, aporta información~\cite[§16.1 y §17.1]{mcs}.
\end{ejemplo}

\section{Independencia}

\begin{definicion}[Independencia]
Dos sucesos $A$ y $B$ son \emph{independientes} si $\PP(A\cap B)=\PP(A)\PP(B)$. Si $\PP(B)>0$, esto equivale a $\PP(A\mid B)=\PP(A)$: saber que ha ocurrido $B$ no cambia la probabilidad de $A$. Los sucesos $A_1,\dots,A_n$ son \emph{independientes} (o mutuamente independientes) si para todo subconjunto de índices $S\subseteq[n]$ con $\abs{S}\ge2$,
\[
\PP\Bigl(\bigcap_{i\in S}A_i\Bigr)=\prod_{i\in S}\PP(A_i).
\]
Son \emph{independientes dos a dos} si la condición se cumple al menos para todos los $S$ con $\abs{S}=2$.
\end{definicion}

\begin{proposicion}\label{prop:t4-complementos}
Si $A$ y $B$ son independientes, también lo son $A$ y $B^{c}$, $A^{c}$ y $B$, y $A^{c}$ y $B^{c}$.
\end{proposicion}

\begin{proof}
$A$ es la unión disjunta de $A\cap B$ y $A\cap B^{c}$, luego $\PP(A\cap B^{c})=\PP(A)-\PP(A)\PP(B)=\PP(A)(1-\PP(B))=\PP(A)\PP(B^{c})$. El segundo caso es simétrico, y el tercero resulta de aplicar el primero a la pareja independiente $A^{c}$, $B$.
\end{proof}

\begin{ejemplo}[Independencia dos a dos sin independencia mutua]\label{ej:t4-dosados}
Se lanzan dos monedas equilibradas; sean $A$ «la primera es cara», $B$ «la segunda es cara» y $C$ «las dos coinciden». Cada suceso tiene probabilidad $1/2$ y cada intersección de dos tiene probabilidad $1/4$ (por ejemplo, $A\cap C$ es «dos caras»), luego son independientes dos a dos. Pero $\PP(A\cap B\cap C)=1/4\neq1/8$: conocer $A$ y $B$ determina $C$.
\end{ejemplo}

\begin{observacion}[Independencia e incompatibilidad]
No deben confundirse. Si $A$ y $B$ son incompatibles con probabilidades positivas, \emph{no} son independientes: $\PP(A\cap B)=0<\PP(A)\PP(B)$; saber que ocurre $A$ garantiza que no ocurre $B$. La independencia es, además, una propiedad de la ley $\PP$, no de los sucesos como conjuntos: en la práctica suele ser una \emph{hipótesis} de modelización (lanzamientos sucesivos, ensayos repetidos) más que algo que se demuestre~\cite[§17.7.2]{mcs}.
\end{observacion}

Si $\PP(C)>0$, los sucesos $A$ y $B$ son \emph{condicionalmente independientes dado $C$} si $\PP(A\cap B\mid C)=\PP(A\mid C)\PP(B\mid C)$. Ni la independencia implica la independencia condicional ni al revés (\cref{eje:4-ajedrez}).

\begin{observacion}[Conexión con el curso]
La regla de la cadena es el fundamento de los modelos de lenguaje autorregresivos (módulo~1): $\PP(x_1,\dots,x_T)=\prod_{t=1}^{T}\PP(x_t\mid x_1,\dots,x_{t-1})$, y un Transformer (módulos~13 a~15) calcula cada factor. La hipótesis de Markov del modelo de $n$-gramas es una hipótesis de independencia condicional. El teorema de Bayes reaparece en el módulo~5 (inferencia bayesiana) y en los clasificadores del módulo~8.
\end{observacion}

\begin{observacion}[Errores frecuentes]
(1) Confundir $\PP(A\mid B)$ con $\PP(B\mid A)$ (la \emph{falacia del fiscal}): que una prueba detecte al $90\,\%$ de los enfermos no significa que el $90\,\%$ de los positivos esté enfermo (\cref{ej:t4-diagnostico}). (2) Olvidar la probabilidad a priori, que en enfermedades raras domina el resultado. (3) Confundir independencia con incompatibilidad. (4) Deducir la independencia mutua de la independencia dos a dos (\cref{ej:t4-dosados}). (5) Condicionar al suceso equivocado: «al menos uno es chica» no es lo mismo que «el mayor es chica» (\cref{ej:t4-hijos}).
\end{observacion}

\begin{lecturas}
\item Bertsekas y Tsitsiklis~\cite{bt}, §1.3 a §1.5 (clases~2 y~3 de MIT 6.041~\cite{mit6041f10}).
\item Blitzstein y Hwang~\cite{bh}, cap.~2 «Conditional Probability», en especial §2.3 «Bayes' rule and the law of total probability», §2.5 «Independence of events» y §2.8 «Pitfalls and paradoxes».
\item Ross~\cite{ross}, cap.~3, §3.2 «Conditional Probabilities», §3.3 «Bayes's Formula» y §3.4 «Independent Events».
\item Lehman, Leighton y Meyer~\cite{mcs}, cap.~17 «Conditional Probability», §17.1 a §17.5, §17.7 y §17.8.
\item Grinstead y Snell~\cite{gs}, §4.1 «Discrete Conditional Probability» y §4.3 «Paradoxes».
\end{lecturas}

\section*{Ejercicios del tema 4}
\addcontentsline{toc}{section}{Ejercicios del tema 4}

\begin{ejercicio}{1}{eje:4-tabla}
En una clase de 100 estudiantes, 40 son de Ingeniería y 60 de Matemáticas; de los de Ingeniería, 30 aprueban, y de los de Matemáticas, 36. Se elige un estudiante al azar. Calcule $\PP(\text{aprueba})$, $\PP(\text{aprueba}\mid\text{Ingeniería})$ y $\PP(\text{Ingeniería}\mid\text{aprueba})$. ¿Son independientes «aprobar» y «ser de Ingeniería»?
\end{ejercicio}

\begin{ejercicio}{2}{eje:4-cachorros}
(a) Una perra tiene dos cachorros; se sabe que el mayor es hembra. ¿Qué probabilidad hay de que los dos lo sean? (b) Una gata tiene dos gatitos; se sabe que al menos uno es macho. ¿Qué probabilidad hay de que los dos lo sean? Suponga sexos equiprobables e independientes.
\fuente{adaptado de MIT 18.05~\cite{mit1805}, \emph{Problem Set}~2, problema~1.}
\end{ejercicio}

\begin{ejercicio}{2}{eje:4-spam}
El $80\,\%$ de los correos que recibe una cuenta es correo basura. La expresión «dinero gratis» aparece en el $10\,\%$ de los correos basura y en el $1\,\%$ de los legítimos. Si un correo contiene esa expresión, ¿qué probabilidad hay de que sea basura?
\fuente{adaptado de Harvard Stat 110~\cite{stat110}, \emph{Strategic Practice}~2, sección «Thinking Conditionally», problema~2.}
\end{ejercicio}

\begin{ejercicio}{2}{eje:4-dos-pruebas}
Una enfermedad afecta al $1\,\%$ de la población y una prueba tiene sensibilidad y especificidad del $95\,\%$. Una persona da positivo en dos pruebas que, condicionadas al estado de salud, son independientes. Calcule la probabilidad de que esté enferma (a) actualizando con las dos pruebas a la vez; (b) actualizando primero con una y después con la otra. Compruebe que coinciden.
\fuente{adaptado de Harvard Stat 110~\cite{stat110}, \emph{Homework}~3, problema~2(b).}
\end{ejercicio}

\begin{ejercicio}{2}{eje:4-piezas}
Un almacén tiene 500 piezas antiguas, de las que el $15\,\%$ son defectuosas, y 1500 nuevas, con un $5\,\%$ de defectuosas. Se lanza una moneda para elegir el tipo y se toman dos piezas de ese tipo, con reemplazamiento. (a) ¿Qué probabilidad hay de que ambas sean defectuosas? (b) Si ambas lo son, ¿qué probabilidad hay de que fueran antiguas?
\fuente{adaptado de MIT 6.041~\cite{mit6041f10}, \emph{Problem Set}~2, problema~3 (se ha fijado el muestreo con reemplazamiento).}
\end{ejercicio}

\begin{ejercicio}{2}{eje:4-perro}
Un perro se ha perdido en el bosque A con probabilidad $0{,}4$ o en el B con probabilidad $0{,}6$. Si está en A y se busca allí un día, se encuentra con probabilidad $0{,}25$; si está en B y se busca allí, con probabilidad $0{,}15$. (a) ¿Dónde conviene buscar el primer día? (b) Si se busca en A y no se encuentra, ¿qué probabilidad hay de que esté en A? (c) Si se elige el bosque lanzando una moneda y el perro aparece el primer día, ¿qué probabilidad hay de que se buscara en A?
\fuente{adaptado de MIT 6.041~\cite{mit6041f10}, \emph{Problem Set}~2, problema~4, apartados (a) a~(c).}
\end{ejercicio}

\begin{ejercicio}{2}{eje:4-monty7}
(a) En una versión del concurso con 7 puertas, el concursante elige una y el presentador, que sabe dónde está el coche, abre 3 de las otras 6, todas con cabra. Si el concursante cambia a una de las 3 puertas cerradas restantes elegida al azar, ¿qué probabilidad tiene de ganar? ¿Le conviene cambiar? (b) Generalice a $n\ge3$ puertas, de las que el presentador abre $m$, con $1\le m\le n-2$.
\fuente{adaptado de Harvard Stat 110~\cite{stat110}, \emph{Homework}~3, problema~1.}
\end{ejercicio}

\begin{ejercicio}{2}{eje:4-monty-p}
En el problema de Monty Hall (\cref{ej:t4-monty}), supóngase que, cuando el presentador puede abrir la puerta~2 o la~3, abre la~2 con probabilidad $p\in[\tfrac12,1]$. Calcule la probabilidad de ganar cambiando (a) sin condicionar; (b) sabiendo que abre la puerta~2; (c) sabiendo que abre la puerta~3.
\fuente{adaptado de Harvard Stat 110~\cite{stat110}, \emph{Strategic Practice}~3, sección «Continuing with Conditioning», problema~1.}
\end{ejercicio}

\begin{ejercicio}{2}{eje:4-paridad}
Se lanzan dos dados. Sean $A$ «el primero es impar», $B$ «el segundo es impar» y $C$ «la suma es impar». Estudie si son independientes dos a dos y si son mutuamente independientes.
\fuente{adaptado de MIT 18.05~\cite{mit1805}, \emph{Problem Set}~3, problema~1(a).}
\end{ejercicio}

\begin{ejercicio}{2}{eje:4-ajedrez}
Un jugador se enfrenta a un rival desconocido que, con probabilidad $1/3$ cada caso, es principiante, intermedio o maestro; contra ellos gana cada partida con probabilidad $0{,}9$, $0{,}5$ y $0{,}3$, respectivamente, y las partidas son independientes \emph{dado} el nivel del rival. (a) Calcule la probabilidad de ganar la primera partida. (b) Calcule la probabilidad de ganar la segunda sabiendo que ganó la primera. (c) ¿Son independientes los resultados de las dos partidas? Explique la diferencia entre independencia e independencia condicional.
\fuente{adaptado de Harvard Stat 110~\cite{stat110}, \emph{Strategic Practice}~2, sección «Thinking Conditionally», problema~5.}
\end{ejercicio}

\begin{ejercicio}{3}{eje:4-si-mismo}
Demuestre que un suceso $A$ es independiente de sí mismo si y solo si $\PP(A)\in\{0,1\}$. Demuestre además que un suceso con $\PP(A)\in\{0,1\}$ es independiente de cualquier otro suceso.
\fuente{adaptado de Harvard Stat 110~\cite{stat110}, \emph{Strategic Practice}~2, sección «Independence», problema~1.}
\end{ejercicio}

\begin{ejercicio}{3}{eje:4-union-indep}
(a) Demuestre que si $A$, $B$ y $C$ son mutuamente independientes, entonces $A$ y $B\cup C$ son independientes. (b) Demuestre que si $A_1,\dots,A_n$ son independientes, $\PP\bigl(\bigcup_{i=1}^{n}A_i\bigr)=1-\prod_{i=1}^{n}\bigl(1-\PP(A_i)\bigr)$.
\fuente{adaptado de MIT 6.041~\cite{mit6041f10}, \emph{Problem Set}~2, problema~5.}
\end{ejercicio}

\begin{ejercicio}{3}{eje:4-producto-triple}
Construya tres sucesos $A$, $B$ y $C$ que cumplan $\PP(A\cap B\cap C)=\PP(A)\PP(B)\PP(C)$ pero que no sean mutuamente independientes, y demuestre que su ejemplo cumple ambas propiedades.
\fuente{adaptado de Harvard Stat 110~\cite{stat110}, \emph{Strategic Practice}~2, sección «Independence», problema~4.}
\end{ejercicio}

\begin{ejercicio}{3}{eje:4-razones}
(a) Demuestre la forma de razones del teorema de Bayes. (b) Sean $E_1$ y $E_2$ dos evidencias condicionalmente independientes dadas $B$ y dadas $B^{c}$. Demuestre que actualizar la razón a priori con $E_1$ y después con $E_2$ da el mismo resultado que actualizar con $E_1\cap E_2$ de una vez.
\fuente{el apartado (a) está adaptado de Harvard Stat 110~\cite{stat110}, \emph{Homework}~3, problema~2(a).}
\end{ejercicio}

\begin{ejercicio}{3}{eje:4-simetria}
Sean $A$ y $B$ sucesos con probabilidades estrictamente entre 0 y 1. Demuestre que las tres afirmaciones siguientes son equivalentes: (a) $\PP(A\mid B)>\PP(A)$; (b) $\PP(B\mid A)>\PP(B)$; (c) $\PP(A\mid B)>\PP(A\mid B^{c})$.
\end{ejercicio}

% =====================================================================
\chapter{Variables aleatorias discretas}
\label{tema:5}
% =====================================================================

Una variable aleatoria asigna un número a cada resultado de un experimento; su distribución resume todo lo que interesa de ella. Este tema introduce la función de masa, la esperanza y la varianza, demuestra la linealidad de la esperanza y la explota mediante variables indicadoras, y estudia las distribuciones discretas básicas: Bernoulli, binomial, geométrica, Poisson y categórica. Se siguen los capítulos~3 y~4 de Blitzstein y Hwang~\cite{bh}, el capítulo~2 de Bertsekas y Tsitsiklis~\cite{bt} (clases~5 a~7 de MIT 6.041~\cite{mit6041f10}), el capítulo~4 de Ross~\cite{ross} y el capítulo~18, «Random Variables», de~\cite{mcs}.

\section{Variables aleatorias y función de masa}

\begin{definicion}[Variable aleatoria discreta]
Una \emph{variable aleatoria} es una función $X\colon\Omega\to\R$. Es \emph{discreta} si toma un conjunto finito o numerable de valores. Su \emph{función de masa} es $p_X(x)=\PP(X=x)$, donde $\{X=x\}$ abrevia el suceso $\{\omega\in\Omega: X(\omega)=x\}$, y su \emph{función de distribución} es $F_X(x)=\PP(X\le x)$ para $x\in\R$.
\end{definicion}

La función de masa cumple $p_X(x)\ge0$ y $\sum_xp_X(x)=1$, porque los sucesos $\{X=x\}$ forman una partición de $\Omega$. Para todo conjunto $S\subseteq\R$, $\PP(X\in S)=\sum_{x\in S}p_X(x)$. La función de distribución es no decreciente, tiende a 0 en $-\infty$ y a 1 en $+\infty$, es continua por la derecha y, en el caso discreto, es escalonada, con saltos de altura $p_X(x)$ en cada valor $x$.

\begin{ejemplo}\label{ej:t5-suma}
Sea $S$ la suma de dos dados. Contando en los 36 resultados equiprobables, $p_S(s)=\frac{6-\abs{s-7}}{36}$ para $s=2,\dots,12$; por ejemplo, $p_S(7)=6/36$ y $p_S(2)=1/36$. La suma de todos los valores es $(1+2+3+4+5+6+5+4+3+2+1)/36=1$.
\end{ejemplo}

\section{Esperanza y linealidad}

\begin{definicion}[Esperanza]
La \emph{esperanza} (o media) de una variable discreta $X$ es
\[
\E[X]=\sum_{x}x\,p_X(x),
\]
siempre que la serie converja absolutamente. Es el centro de masas de la distribución.
\end{definicion}

\begin{proposicion}[Esperanza de una función: LOTUS]\label{prop:t5-lotus}
Si $X$ es discreta y $g\colon\R\to\R$, entonces $\E[g(X)]=\sum_xg(x)\,p_X(x)$. Si $\Omega$ es discreto, además $\E[X]=\sum_{\omega\in\Omega}X(\omega)\,\PP(\{\omega\})$.
\end{proposicion}

\begin{proof}
Sea $Y=g(X)$. Para cada valor $y$, $\{Y=y\}$ es la unión disjunta de los $\{X=x\}$ con $g(x)=y$, luego $p_Y(y)=\sum_{x:g(x)=y}p_X(x)$ y
\[
\E[Y]=\sum_yy\sum_{x:g(x)=y}p_X(x)=\sum_y\sum_{x:g(x)=y}g(x)p_X(x)=\sum_xg(x)p_X(x),
\]
reordenando la serie, lo que es lícito por la convergencia absoluta. La segunda afirmación es el mismo argumento con $\Omega$ en el papel de los valores de $X$: $\{X=x\}$ es la unión disjunta de los $\{\omega\}$ con $X(\omega)=x$.
\end{proof}

\begin{teorema}[Linealidad de la esperanza]\label{teo:t5-linealidad}
Si $X$ e $Y$ son variables aleatorias discretas definidas sobre el mismo espacio y con esperanza finita, y $a,b\in\R$, entonces
\[
\E[aX+bY]=a\,\E[X]+b\,\E[Y].
\]
No se requiere ninguna hipótesis de independencia.
\end{teorema}

\begin{proof}
Se demuestra cuando $\Omega$ es discreto; el caso general es análogo usando la distribución conjunta. Por la \cref{prop:t5-lotus},
\[
\E[aX+bY]=\sum_{\omega}\bigl(aX(\omega)+bY(\omega)\bigr)\PP(\{\omega\})=a\sum_\omega X(\omega)\PP(\{\omega\})+b\sum_\omega Y(\omega)\PP(\{\omega\}),
\]
que es $a\E[X]+b\E[Y]$.
\end{proof}

\begin{definicion}[Variable indicadora]
La \emph{indicadora} de un suceso $A$ es la variable $I_A$ que vale 1 si ocurre $A$ y 0 si no.
\end{definicion}

\begin{proposicion}[Puente fundamental]\label{prop:t5-puente}
$\E[I_A]=\PP(A)$. En consecuencia, si $X=I_{A_1}+\dots+I_{A_n}$ cuenta cuántos de los sucesos $A_1,\dots,A_n$ ocurren, $\E[X]=\PP(A_1)+\dots+\PP(A_n)$.
\end{proposicion}

\begin{proof}
$\E[I_A]=1\cdot\PP(A)+0\cdot\PP(A^{c})=\PP(A)$. La segunda parte es la linealidad.
\end{proof}

\begin{ejemplo}[Puntos fijos]\label{ej:t5-fijos}
Se reparten al azar los sombreros de $n$ personas. Sea $X$ el número de personas que reciben el suyo. Con $A_i$ «la persona $i$ recibe su sombrero», $\PP(A_i)=(n-1)!/n!=1/n$, luego $\E[X]=n\cdot\frac1n=1$ para todo $n$. Calcular la distribución de $X$ es mucho más difícil (\cref{ej:t2-desarreglos}); la esperanza, gracias a las indicadoras, es inmediata, aunque los $A_i$ no son independientes~\cite[§4.4]{bh}.
\end{ejemplo}

\section{Varianza e independencia}

\begin{definicion}[Varianza]
La \emph{varianza} de $X$ es $\Var(X)=\E\bigl[(X-\mu)^{2}\bigr]$, con $\mu=\E[X]$, y su \emph{desviación típica} es $\sigma_X=\sqrt{\Var(X)}$. Mide la dispersión de $X$ alrededor de su media, en unidades al cuadrado.
\end{definicion}

\begin{proposicion}\label{prop:t5-var}
$\Var(X)=\E[X^{2}]-\E[X]^{2}$ y, para $a,b\in\R$, $\Var(aX+b)=a^{2}\Var(X)$. En particular $\Var(X)\ge0$ implica $\E[X^{2}]\ge\E[X]^{2}$.
\end{proposicion}

\begin{proof}
Por linealidad, $\E[(X-\mu)^{2}]=\E[X^{2}]-2\mu\E[X]+\mu^{2}=\E[X^{2}]-\mu^{2}$. Como $\E[aX+b]=a\mu+b$, se tiene $\Var(aX+b)=\E[(aX+b-a\mu-b)^{2}]=a^{2}\E[(X-\mu)^{2}]$.
\end{proof}

\begin{definicion}[Independencia de variables]
Las variables discretas $X_1,\dots,X_n$ son \emph{independientes} si $\PP(X_1=x_1,\dots,X_n=x_n)=\prod_i\PP(X_i=x_i)$ para todos los valores $x_1,\dots,x_n$.
\end{definicion}

\begin{proposicion}\label{prop:t5-producto}
Si $X$ e $Y$ son discretas, independientes y con esperanza finita, $\E[XY]=\E[X]\E[Y]$. En consecuencia, si $X_1,\dots,X_n$ son independientes (basta dos a dos), $\Var(X_1+\dots+X_n)=\Var(X_1)+\dots+\Var(X_n)$.
\end{proposicion}

\begin{proof}
Por LOTUS aplicada al par $(X,Y)$ y la independencia,
\[
\E[XY]=\sum_{x,y}xy\,\PP(X=x,Y=y)=\sum_x\sum_yxy\,p_X(x)p_Y(y)=\Bigl(\sum_xx\,p_X(x)\Bigr)\Bigl(\sum_yy\,p_Y(y)\Bigr).
\]
Para la varianza, sean $\mu_i=\E[X_i]$ y $S=\sum_iX_i$. Entonces $S-\E[S]=\sum_i(X_i-\mu_i)$ y
\[
\Var(S)=\E\Bigl[\sum_{i,j}(X_i-\mu_i)(X_j-\mu_j)\Bigr]=\sum_i\Var(X_i)+\sum_{i\neq j}\E[X_i-\mu_i]\,\E[X_j-\mu_j]=\sum_i\Var(X_i),
\]
porque para $i\neq j$ las variables $X_i-\mu_i$ y $X_j-\mu_j$ son independientes y tienen media~0.
\end{proof}

\section{Distribuciones discretas básicas}

\begin{definicion}[Bernoulli y binomial]
$X\sim\Ber(p)$, con $p\in[0,1]$, si $\PP(X=1)=p$ y $\PP(X=0)=1-p$: es la indicadora de un «éxito» de probabilidad $p$. $X\sim\Bin(n,p)$ si $X$ es el número de éxitos en $n$ ensayos de Bernoulli independientes con la misma probabilidad $p$.
\end{definicion}

\begin{teorema}\label{teo:t5-binomial}
Si $X\sim\Bin(n,p)$, entonces $\PP(X=k)=\binom nk p^{k}(1-p)^{n-k}$ para $k=0,\dots,n$, $\E[X]=np$ y $\Var(X)=np(1-p)$.
\end{teorema}

\begin{proof}
Cada sucesión concreta de resultados con $k$ éxitos y $n-k$ fracasos tiene probabilidad $p^{k}(1-p)^{n-k}$ por la independencia, y hay $\binom nk$ sucesiones así (se eligen las posiciones de los éxitos); por la aditividad se obtiene la función de masa, que suma~1 por el binomio de Newton. Para los momentos, $X=I_1+\dots+I_n$ con $I_j\sim\Ber(p)$ independientes, $\E[I_j]=p$ y $\Var(I_j)=\E[I_j^{2}]-p^{2}=p-p^{2}$, porque $I_j^{2}=I_j$. Se aplican la linealidad y la \cref{prop:t5-producto}.
\end{proof}

\begin{definicion}[Geométrica]
$X\sim\Geom(p)$, con $0<p\le1$, si $X$ es el número de ensayos de Bernoulli independientes de parámetro $p$ necesarios hasta obtener el primer éxito (incluido este). Su función de masa es $\PP(X=k)=(1-p)^{k-1}p$ para $k=1,2,\dots$
\end{definicion}

\begin{proposicion}\label{prop:t5-geom}
Si $X\sim\Geom(p)$: (i) $\PP(X>k)=(1-p)^{k}$; (ii) $\E[X]=1/p$ y $\Var(X)=(1-p)/p^{2}$; (iii) \emph{falta de memoria}: $\PP(X>m+k\mid X>m)=\PP(X>k)$ para enteros $m,k\ge0$.
\end{proposicion}

\begin{proof}
(i) $X>k$ si y solo si los $k$ primeros ensayos son fracasos. (ii) Con $q=1-p$, derivando la serie geométrica $\sum_{k\ge0}x^{k}=1/(1-x)$ para $\abs{x}<1$ se obtiene $\sum_{k\ge1}kx^{k-1}=1/(1-x)^{2}$, luego $\E[X]=p\sum_{k\ge1}kq^{k-1}=p/p^{2}=1/p$. Derivando otra vez, $\sum_{k\ge2}k(k-1)x^{k-2}=2/(1-x)^{3}$, de donde $\E[X(X-1)]=pq\cdot2/p^{3}=2q/p^{2}$ y $\Var(X)=\E[X(X-1)]+\E[X]-\E[X]^{2}=\frac{2q}{p^{2}}+\frac1p-\frac1{p^{2}}=\frac{q}{p^{2}}$. (iii) $\PP(X>m+k\mid X>m)=q^{m+k}/q^{m}=q^{k}$.
\end{proof}

\begin{definicion}[Poisson]
$X\sim\Poi(\lambda)$, con $\lambda>0$, si $\PP(X=k)=e^{-\lambda}\lambda^{k}/k!$ para $k=0,1,2,\dots$ Se cumple $\E[X]=\Var(X)=\lambda$ (\cref{eje:5-poisson}).
\end{definicion}

La distribución de Poisson modela el número de sucesos raros entre muchas oportunidades independientes: erratas en un libro, llegadas de peticiones a un servidor en un minuto, apariciones de una palabra poco frecuente en un texto largo. El resultado siguiente lo justifica.

\begin{teorema}[Aproximación de Poisson a la binomial]\label{teo:t5-poisson}
Sea $\lambda>0$ y $p_n=\lambda/n$. Para cada $k\ge0$ fijo,
\(
\lim_{n\to\infty}\binom nk p_n^{k}(1-p_n)^{n-k}=e^{-\lambda}\frac{\lambda^{k}}{k!}.
\)
\end{teorema}

\begin{proof}
Se escribe
\[
\binom nk\Bigl(\frac\lambda n\Bigr)^{k}\Bigl(1-\frac\lambda n\Bigr)^{n-k}=\frac{\lambda^{k}}{k!}\cdot\frac{n(n-1)\cdots(n-k+1)}{n^{k}}\cdot\Bigl(1-\frac\lambda n\Bigr)^{n}\cdot\Bigl(1-\frac\lambda n\Bigr)^{-k}.
\]
El segundo factor es $\prod_{j=0}^{k-1}(1-j/n)\to1$; el tercero tiende a $e^{-\lambda}$, por el límite $(1+x/n)^{n}\to e^{x}$; el cuarto tiende a~1 porque $k$ está fijo.
\end{proof}

\begin{definicion}[Categórica]
Sea $\V=\{v_1,\dots,v_K\}$ un conjunto finito y $\bm p=(p_1,\dots,p_K)$ con $p_i\ge0$ y $\sum_ip_i=1$. Una variable $X$ con valores en $\V$ tiene distribución \emph{categórica}, $X\sim\Cat(\bm p)$, si $\PP(X=v_i)=p_i$. Cuando los valores no son números, en lugar de la esperanza de $X$ se usan las de sus indicadoras: $\E[I_{\{X=v_i\}}]=p_i$.
\end{definicion}

\begin{table}[h]
\centering
\begin{tabular}{@{}lllll@{}}
\toprule
Distribución & Valores & $\PP(X=k)$ & $\E[X]$ & $\Var(X)$ \\
\midrule
$\Ber(p)$ & $\{0,1\}$ & $p^{k}(1-p)^{1-k}$ & $p$ & $p(1-p)$ \\
$\Bin(n,p)$ & $\{0,\dots,n\}$ & $\binom nk p^{k}(1-p)^{n-k}$ & $np$ & $np(1-p)$ \\
$\Geom(p)$ & $\{1,2,\dots\}$ & $(1-p)^{k-1}p$ & $1/p$ & $(1-p)/p^{2}$ \\
$\Poi(\lambda)$ & $\{0,1,\dots\}$ & $e^{-\lambda}\lambda^{k}/k!$ & $\lambda$ & $\lambda$ \\
Uniforme en $[n]$ & $\{1,\dots,n\}$ & $1/n$ & $(n+1)/2$ & $(n^{2}-1)/12$ \\
\bottomrule
\end{tabular}
\caption{Distribuciones discretas básicas.}
\label{tab:t5-discretas}
\end{table}

\begin{observacion}[Conexión con el curso]
La salida de un modelo de lenguaje es, en cada paso, una distribución categórica sobre el vocabulario, obtenida con la softmax (módulos~1 y~5); generar texto es muestrear de ella. La linealidad de la esperanza con indicadoras calcula, por ejemplo, el número esperado de tokens distintos en una muestra (\cref{eje:5-categorica}) o el número esperado de colisiones en una tabla \emph{hash} (\cref{eje:5-hash}, módulo~16). La métrica pass@$k$ del módulo~22 es una probabilidad calculada en un modelo hipergeométrico (\cref{eje:5-passk}).
\end{observacion}

\begin{observacion}[Errores frecuentes]
(1) Creer que la linealidad de la esperanza exige independencia: no la exige; la aditividad de la \emph{varianza} sí. (2) Escribir $\E[g(X)]=g(\E[X])$; en general es falso, por ejemplo $\E[X^{2}]\ge\E[X]^{2}$ con igualdad solo si $X$ es constante. (3) Escribir $\Var(aX)=a\Var(X)$ o $\Var(X-Y)=\Var(X)-\Var(Y)$: para $X$, $Y$ independientes, $\Var(X-Y)=\Var(X)+\Var(Y)$. (4) Confundir las dos versiones de la geométrica: la de este tema cuenta ensayos (valores $1,2,\dots$); otros textos cuentan fracasos (valores $0,1,\dots$), con media $(1-p)/p$. (5) Aplicar la binomial a extracciones \emph{sin} reemplazamiento, que no son independientes.
\end{observacion}

\begin{lecturas}
\item Blitzstein y Hwang~\cite{bh}, cap.~3, §3.1 a §3.3 y §3.8, y cap.~4, §4.1 a §4.4, §4.6 y §4.7.
\item Bertsekas y Tsitsiklis~\cite{bt}, §2.1 a §2.4 (clases~5 y~6 de MIT 6.041~\cite{mit6041f10}).
\item Ross~\cite{ross}, cap.~4, §4.1 a §4.7 y §4.9.
\item Lehman, Leighton y Meyer~\cite{mcs}, cap.~18, §18.1 a §18.5, en especial §18.5 «Linearity of Expectation».
\item Grinstead y Snell~\cite{gs}, §5.1 «Important Distributions» y §6.1 y §6.2.
\end{lecturas}

\section*{Ejercicios del tema 5}
\addcontentsline{toc}{section}{Ejercicios del tema 5}

\begin{ejercicio}{1}{eje:5-masa}
Sea $X$ con $\PP(X=k)=ck$ para $k=1,2,3,4$. Halle $c$, la función de distribución $F_X$ en los puntos $1,2,3,4$, $\E[X]$ y $\Var(X)$.
\end{ejercicio}

\begin{ejercicio}{2}{eje:5-diferencia}
Se lanzan dos dados equilibrados de tres caras (valores 1, 2 y 3) y se define $X$ como el valor del primero menos el del segundo. (a) Halle la función de masa de $X$, $\E[X]$ y $\Var(X)$. (b) Halle la función de masa de $X^{2}$ y compruebe que $\E[X^{2}]=\Var(X)$.
\fuente{adaptado de MIT 6.041~\cite{mit6041f10}, \emph{Problem Set}~3, problema~4.}
\end{ejercicio}

\begin{ejercicio}{2}{eje:5-binomial}
Un modelo resuelve cada uno de 10 problemas independientes con probabilidad $0{,}3$. Sea $X$ el número de problemas resueltos. Calcule $\PP(X=3)$, $\PP(X\ge1)$, $\E[X]$ y $\Var(X)$.
\end{ejercicio}

\begin{ejercicio}{2}{eje:5-geometrica}
Se repite un experimento con probabilidad de éxito $0{,}2$ hasta el primer éxito. Sea $X$ el número de intentos. Calcule $\PP(X=3)$, $\PP(X>5)$, $\E[X]$ y $\Var(X)$. Si ya se han fallado 3 intentos, ¿qué probabilidad hay de necesitar más de 8 en total?
\end{ejercicio}

\begin{ejercicio}{2}{eje:5-poisson-aprox}
En un texto de 1000 tokens, cada token es, independientemente, una palabra rara con probabilidad $0{,}002$. Sea $X$ el número de palabras raras. Calcule $\PP(X=0)$ y $\PP(X\ge2)$ con la binomial exacta y con la aproximación de Poisson, y compare.
\end{ejercicio}

\begin{ejercicio}{2}{eje:5-cumple-esp}
Hay 50 personas con cumpleaños uniformes e independientes entre 365 días. Use variables indicadoras para calcular (a) el número esperado de parejas de personas que cumplen años el mismo día; (b) el número esperado de días del año en que cumplen años al menos dos personas.
\fuente{adaptado de Harvard Stat 110~\cite{stat110}, \emph{Strategic Practice}~4, sección~2, problema~1.}
\end{ejercicio}

\begin{ejercicio}{2}{eje:5-hash}
Se insertan $k$ claves en una tabla \emph{hash} de $n$ casillas; cada clave va a una casilla uniforme, independientemente de las demás. Halle el número esperado de casillas vacías, de casillas con exactamente una clave y de casillas con más de una. Evalúe para $k=n=100$.
\fuente{adaptado de Harvard Stat 110~\cite{stat110}, \emph{Strategic Practice}~4, sección~2, problema~4.}
\end{ejercicio}

\begin{ejercicio}{2}{eje:5-hijos}
Una pareja tiene hijos hasta tener al menos un niño y una niña. Suponiendo sexos equiprobables e independientes, calcule el número esperado de hijos.
\fuente{adaptado de Harvard Stat 110~\cite{stat110}, \emph{Homework}~4, problema~3.}
\end{ejercicio}

\begin{ejercicio}{2}{eje:5-inversiones}
Una \emph{inversión} de una lista de números distintos es un par de posiciones $i<j$ con el elemento en $i$ mayor que el elemento en $j$. Si la lista de $n$ números distintos está en orden aleatorio uniforme, calcule el número esperado de inversiones. Compruébelo enumerando las $4!$ permutaciones de 4 elementos (o un subconjunto representativo).
\fuente{adaptado de Harvard Stat 110~\cite{stat110}, \emph{Homework}~4, problema~6(a).}
\end{ejercicio}

\begin{ejercicio}{2}{eje:5-categorica}
Un modelo genera 10 tokens independientes con distribución categórica $\PP(\texttt{a})=0{,}5$, $\PP(\texttt{b})=0{,}3$, $\PP(\texttt{c})=0{,}2$. (a) Calcule el número esperado de veces que aparece cada token y la varianza del número de apariciones de \texttt{a}. (b) Calcule el número esperado de tokens distintos que aparecen.
\end{ejercicio}

\begin{ejercicio}{2}{eje:5-passk}
Un modelo genera $n=10$ soluciones a un problema, de las que $c=3$ son correctas. Se eligen $k$ de las 10 al azar, sin reemplazamiento. (a) Demuestre que la probabilidad de que al menos una de las $k$ sea correcta es $1-\binom{n-c}{k}\big/\binom{n}{k}$. (b) Evalúe para $k=1,2,5$.
\end{ejercicio}

\begin{ejercicio}{3}{eje:5-colas}
(a) Demuestre que si $X$ toma valores en $\{0,1,2,\dots\}$, entonces $\E[X]=\sum_{k\ge1}\PP(X\ge k)$, escribiendo $X=\sum_{k\ge1}I_{\{X\ge k\}}$ o intercambiando el orden de una suma doble. (b) Úselo para obtener de nuevo la esperanza de una $\Geom(p)$.
\fuente{adaptado de MIT 6.041~\cite{mit6041f10}, \emph{Problem Set}~3, problema~3(a), y de Harvard Stat 110~\cite{stat110}, \emph{Strategic Practice}~4, sección~1, problema~5.}
\end{ejercicio}

\begin{ejercicio}{3}{eje:5-poisson}
Sea $X\sim\Poi(\lambda)$. Demuestre que su función de masa suma~1, que $\E[X]=\lambda$ y que $\E[X(X-1)]=\lambda^{2}$, y deduzca que $\Var(X)=\lambda$.
\end{ejercicio}

\begin{ejercicio}{3}{eje:5-sin-memoria}
Sea $X$ una variable con valores en $\{1,2,\dots\}$ y $\PP(X>1)\in(0,1)$, tal que $\PP(X>m+k\mid X>m)=\PP(X>k)$ para todos los enteros $m,k\ge0$ con $\PP(X>m)>0$. Demuestre que $X$ tiene distribución geométrica. \emph{Indicación:} estudie $g(k)=\PP(X>k)$.
\end{ejercicio}

\begin{ejercicio}{3}{eje:5-geom-cond}
Sean $X$ e $Y$ independientes con distribución $\Geom(p)$. Demuestre que, para $n\ge2$, la distribución de $X$ condicionada a $X+Y=n$ es uniforme en $\{1,\dots,n-1\}$.
\fuente{adaptado de MIT 6.041~\cite{mit6041f10}, \emph{Problem Set}~4, problema~3.}
\end{ejercicio}

% =====================================================================
\chapter{Variables continuas, desigualdades y ley de los grandes números}
\label{tema:6}
% =====================================================================

Muchas magnitudes, como tiempos, longitudes o pesos de una red neuronal, toman valores en un continuo; su distribución se describe con una densidad en lugar de una función de masa. Este tema introduce las variables continuas y tres distribuciones fundamentales (uniforme, exponencial y normal), y termina con las desigualdades de Markov y Chebyshev y la ley débil de los grandes números, que justifica que una media muestral aproxime una esperanza. Se siguen el capítulo~5 de Blitzstein y Hwang~\cite{bh} (§5.1, §5.2, §5.4 y §5.5) y su §10.1--10.2; las secciones §3.1 a §3.3 y §5.1 a §5.2 de Bertsekas y Tsitsiklis~\cite{bt} (clases~8 y~19 de MIT 6.041~\cite{mit6041f10}); los capítulos~5 y~8 de Ross~\cite{ross}; el capítulo~19 de~\cite{mcs} y el capítulo~8 de Grinstead y Snell~\cite{gs}.

\section{Densidad y función de distribución}

\begin{definicion}[Variable continua]
Una variable aleatoria $X$ es \emph{continua} si existe una función $f_X\colon\R\to[0,\infty)$, su \emph{densidad}, tal que para todos $a\le b$
\[
\PP(a\le X\le b)=\int_a^bf_X(x)\,dx.
\]
Necesariamente $\int_{-\infty}^{\infty}f_X(x)\,dx=1$. La \emph{función de distribución} es $F_X(x)=\PP(X\le x)=\int_{-\infty}^{x}f_X(t)\,dt$.
\end{definicion}

Consecuencias inmediatas: $\PP(X=x)=\int_x^xf_X=0$ para todo $x$, de modo que da igual escribir $<$ o $\le$; $F_X$ es continua y no decreciente, y en los puntos donde $f_X$ es continua, $F_X'(x)=f_X(x)$ por el teorema fundamental del cálculo. La densidad \emph{no} es una probabilidad y puede ser mayor que~1: lo que se interpreta es $\PP(x\le X\le x+\delta)\approx f_X(x)\,\delta$ para $\delta$ pequeño.

\begin{definicion}[Esperanza y varianza]
Para $X$ continua, $\E[X]=\int_{-\infty}^{\infty}x\,f_X(x)\,dx$ (si la integral converge absolutamente) y, para una función $g$, $\E[g(X)]=\int g(x)f_X(x)\,dx$. La varianza se define como en el caso discreto, $\Var(X)=\E[(X-\E X)^{2}]=\E[X^{2}]-\E[X]^{2}$, y la linealidad y las propiedades de la \cref{prop:t5-var} siguen siendo válidas.
\end{definicion}

\section{Distribuciones uniforme y exponencial}

\begin{definicion}[Uniforme]
$X\sim\Unif(a,b)$, con $a<b$, si $f_X(x)=\frac{1}{b-a}$ en $[a,b]$ y $0$ fuera. Entonces $F_X(x)=\frac{x-a}{b-a}$ para $x\in[a,b]$: la probabilidad de un subintervalo es proporcional a su longitud.
\end{definicion}

\begin{proposicion}\label{prop:t6-unif}
Si $X\sim\Unif(a,b)$, entonces $\E[X]=\frac{a+b}{2}$ y $\Var(X)=\frac{(b-a)^{2}}{12}$.
\end{proposicion}

\begin{proof}
$\E[X]=\frac1{b-a}\int_a^bx\,dx=\frac{b^{2}-a^{2}}{2(b-a)}=\frac{a+b}2$. Para la varianza basta el caso $\Unif(0,1)$, porque $X=a+(b-a)U$ con $U\sim\Unif(0,1)$ y $\Var(X)=(b-a)^{2}\Var(U)$. Para $U$, $\E[U^{2}]=\int_0^1u^{2}du=\frac13$ y $\Var(U)=\frac13-\frac14=\frac1{12}$.
\end{proof}

\begin{definicion}[Exponencial]
$X\sim\Expo(\lambda)$, con $\lambda>0$, si $f_X(x)=\lambda e^{-\lambda x}$ para $x\ge0$ y $0$ para $x<0$. Integrando, $F_X(x)=1-e^{-\lambda x}$ y $\PP(X>x)=e^{-\lambda x}$ para $x\ge0$. Se cumple $\E[X]=1/\lambda$ y $\Var(X)=1/\lambda^{2}$ (\cref{eje:6-expo}).
\end{definicion}

\begin{proposicion}[Falta de memoria]\label{prop:t6-memoria}
Si $X\sim\Expo(\lambda)$, entonces $\PP(X>s+t\mid X>s)=\PP(X>t)$ para todos $s,t\ge0$.
\end{proposicion}

\begin{proof}
Como $\{X>s+t\}\subseteq\{X>s\}$, $\PP(X>s+t\mid X>s)=\frac{e^{-\lambda(s+t)}}{e^{-\lambda s}}=e^{-\lambda t}$.
\end{proof}

La exponencial es la versión continua de la geométrica: modela tiempos de espera hasta un suceso que «no envejece», como la llegada de la siguiente petición a un servidor.

\section{Distribución normal y tipificación}

\begin{definicion}[Normal]
$Z\sim\Nor(0,1)$ (\emph{normal estándar}) si su densidad es $\varphi(z)=\frac{1}{\sqrt{2\pi}}e^{-z^{2}/2}$; su función de distribución se denota $\Phi$. Para $\mu\in\R$ y $\sigma>0$, $X\sim\Nor(\mu,\sigma^{2})$ si
\[
f_X(x)=\frac{1}{\sigma\sqrt{2\pi}}\exp\Bigl(-\frac{(x-\mu)^{2}}{2\sigma^{2}}\Bigr).
\]
\end{definicion}

Que $\varphi$ integra~1 equivale a la integral gaussiana $\int_{-\infty}^{\infty}e^{-z^{2}/2}dz=\sqrt{2\pi}$, que se demuestra en el módulo~5 con coordenadas polares. Por la simetría de $\varphi$, $\E[Z]=0$ y $\Phi(-z)=1-\Phi(z)$; integrando por partes se obtiene $\E[Z^{2}]=1$, luego $\Var(Z)=1$. $\Phi$ no tiene expresión con funciones elementales y se tabula (\cref{tab:t6-phi}).

\begin{teorema}[Tipificación]\label{teo:t6-tipificacion}
Si $Z\sim\Nor(0,1)$, entonces $X=\mu+\sigma Z\sim\Nor(\mu,\sigma^{2})$. Recíprocamente, si $X\sim\Nor(\mu,\sigma^{2})$, la variable tipificada $Z=(X-\mu)/\sigma$ es $\Nor(0,1)$ y
\[
\PP(X\le x)=\Phi\Bigl(\frac{x-\mu}{\sigma}\Bigr).
\]
En particular, $\E[X]=\mu$ y $\Var(X)=\sigma^{2}$.
\end{teorema}

\begin{proof}
Sea $X=\mu+\sigma Z$. Como $\sigma>0$, $F_X(x)=\PP(Z\le(x-\mu)/\sigma)=\Phi\bigl((x-\mu)/\sigma\bigr)$. Derivando con la regla de la cadena, $f_X(x)=\frac1\sigma\varphi\bigl(\frac{x-\mu}\sigma\bigr)$, que es la densidad $\Nor(\mu,\sigma^{2})$. El recíproco es el mismo cálculo leído al revés: si $X$ tiene esa densidad, $\PP\bigl((X-\mu)/\sigma\le z\bigr)=\PP(X\le\mu+\sigma z)=\int_{-\infty}^{\mu+\sigma z}\frac1\sigma\varphi\bigl(\frac{x-\mu}\sigma\bigr)dx=\int_{-\infty}^{z}\varphi(t)\,dt$ con el cambio $t=(x-\mu)/\sigma$. Por último, $\E[X]=\mu+\sigma\E[Z]=\mu$ y $\Var(X)=\sigma^{2}\Var(Z)=\sigma^{2}$.
\end{proof}

\begin{table}[h]
\centering
\begin{tabular}{@{}lcccccccc@{}}
\toprule
$z$ & $0{,}5$ & $1$ & $1{,}5$ & $1{,}645$ & $1{,}96$ & $2$ & $2{,}5$ & $3$ \\
\midrule
$\Phi(z)$ & $0{,}6915$ & $0{,}8413$ & $0{,}9332$ & $0{,}9500$ & $0{,}9750$ & $0{,}9772$ & $0{,}9938$ & $0{,}9987$ \\
\bottomrule
\end{tabular}
\caption{Valores de la función de distribución normal estándar. Para $z<0$, $\Phi(z)=1-\Phi(-z)$.}
\label{tab:t6-phi}
\end{table}

\begin{ejemplo}[Regla 68--95--99,7]\label{ej:t6-regla}
Si $X\sim\Nor(\mu,\sigma^{2})$, $\PP(\abs{X-\mu}\le k\sigma)=\PP(\abs Z\le k)=2\Phi(k)-1$, que vale $0{,}6827$, $0{,}9545$ y $0{,}9973$ para $k=1,2,3$. Por ejemplo, si las estaturas siguen una $\Nor(170,\,10^{2})$ (en cm), $\PP(\text{estatura}>185)=1-\Phi(1{,}5)=0{,}0668$.
\end{ejemplo}

\section{Desigualdades de Markov y Chebyshev}

\begin{teorema}[Desigualdad de Markov]\label{teo:t6-markov}
Si $X\ge0$ y $a>0$, entonces $\PP(X\ge a)\le\dfrac{\E[X]}{a}$.
\end{teorema}

\begin{proof}
Se cumple la desigualdad entre variables $a\,I_{\{X\ge a\}}\le X$: si $X\ge a$, el lado izquierdo vale $a\le X$; si no, vale $0\le X$. Tomando esperanzas, que conservan las desigualdades por linealidad (si $Y\ge0$, $\E[Y]\ge0$), $a\,\PP(X\ge a)\le\E[X]$.
\end{proof}

\begin{teorema}[Desigualdad de Chebyshev]\label{teo:t6-chebyshev}
Si $X$ tiene media $\mu$ y varianza $\sigma^{2}$ finitas, para todo $c>0$,
\[
\PP(\abs{X-\mu}\ge c)\le\frac{\sigma^{2}}{c^{2}}.
\]
\end{teorema}

\begin{proof}
Se aplica Markov a la variable no negativa $(X-\mu)^{2}$ con $a=c^{2}$, observando que $\{\abs{X-\mu}\ge c\}=\{(X-\mu)^{2}\ge c^{2}\}$: $\PP(\abs{X-\mu}\ge c)\le\E[(X-\mu)^{2}]/c^{2}=\sigma^{2}/c^{2}$.
\end{proof}

Con $c=k\sigma$, Chebyshev dice que $\PP(\abs{X-\mu}\ge k\sigma)\le1/k^{2}$ para \emph{cualquier} distribución: a dos desviaciones típicas hay como mucho un $25\,\%$ de la probabilidad. Para la normal el valor exacto es un $4{,}6\,\%$; la cota es peor porque no usa nada de la forma de la distribución.

\section{Ley débil de los grandes números}

\begin{proposicion}[Media muestral]\label{prop:t6-media}
Sean $X_1,\dots,X_n$ independientes con la misma media $\mu$ y la misma varianza $\sigma^{2}$, y sea $\Xbar_n=\frac1n\sum_{i=1}^{n}X_i$. Entonces $\E[\Xbar_n]=\mu$ y $\Var(\Xbar_n)=\sigma^{2}/n$.
\end{proposicion}

\begin{proof}
Por linealidad, $\E[\Xbar_n]=\frac1n\cdot n\mu=\mu$. Por la \cref{prop:t5-producto} (válida también para variables continuas independientes) y la \cref{prop:t5-var}, $\Var(\Xbar_n)=\frac1{n^{2}}\sum_i\Var(X_i)=\frac{n\sigma^{2}}{n^{2}}$.
\end{proof}

\begin{teorema}[Ley débil de los grandes números]\label{teo:t6-lgn}
En las condiciones de la \cref{prop:t6-media}, para todo $\varepsilon>0$,
\[
\PP\bigl(\abs{\Xbar_n-\mu}\ge\varepsilon\bigr)\le\frac{\sigma^{2}}{n\varepsilon^{2}}\xrightarrow[n\to\infty]{}0.
\]
Se dice que $\Xbar_n$ \emph{converge en probabilidad} a $\mu$.
\end{teorema}

\begin{proof}
Se aplica la desigualdad de Chebyshev a $\Xbar_n$, cuya media es $\mu$ y cuya varianza es $\sigma^{2}/n$.
\end{proof}

\begin{corolario}[Frecuencias relativas]\label{cor:t6-frecuencias}
Si $A$ es un suceso de probabilidad $p$ y $f_n$ es la frecuencia relativa con que ocurre en $n$ repeticiones independientes, entonces $\PP(\abs{f_n-p}\ge\varepsilon)\le\dfrac{p(1-p)}{n\varepsilon^{2}}\le\dfrac{1}{4n\varepsilon^{2}}$.
\end{corolario}

\begin{proof}
$f_n$ es la media de $n$ indicadoras independientes $\Ber(p)$, de varianza $p(1-p)$, y $p(1-p)=\frac14-\bigl(p-\frac12\bigr)^{2}\le\frac14$.
\end{proof}

El corolario reconcilia la definición axiomática con la interpretación frecuentista: la probabilidad de un suceso es el valor al que tiende, en probabilidad, su frecuencia relativa. Con $\varepsilon=0{,}01$ y $n=10^{5}$, la cota es $1/(4\cdot10^{5}\cdot10^{-4})=0{,}025$.

\begin{observacion}[Conexión con el curso]
La normal aparece en la inicialización de los pesos (módulo~5) y en la activación GELU, $x\,\Phi(x)$ (módulo~11). La ley de los grandes números justifica que la entropía cruzada media sobre un corpus estime la pérdida esperada (módulos~1 y~5) y que el gradiente de un minilote estime el gradiente completo (módulo~7); la \cref{prop:t6-media} explica que el error típico de estas estimaciones decrezca como $1/\sqrt n$. El módulo~5 completa este tema con el teorema central del límite.
\end{observacion}

\begin{observacion}[Errores frecuentes]
(1) Interpretar $f_X(x)$ como una probabilidad o exigir que sea $\le1$. (2) Olvidar que $\PP(X=x)=0$ para variables continuas, o escribir que un suceso de probabilidad~0 es imposible. (3) Tipificar con la varianza en lugar de la desviación típica: $Z=(X-\mu)/\sigma$, no $(X-\mu)/\sigma^{2}$. (4) Aplicar Markov a variables que pueden ser negativas. (5) Leer la ley de los grandes números como una «compensación»: tras muchas caras seguidas, la probabilidad de cruz sigue siendo $1/2$; la frecuencia se acerca a $1/2$ porque las desviaciones pasadas quedan diluidas, no corregidas.
\end{observacion}

\begin{lecturas}
\item Blitzstein y Hwang~\cite{bh}, cap.~5, §5.1 «Probability density functions», §5.2 «Uniform», §5.4 «Normal» y §5.5 «Exponential»; cap.~10, §10.1 «Inequalities» y §10.2 «Law of large numbers».
\item Bertsekas y Tsitsiklis~\cite{bt}, §3.1 a §3.3 y §5.1 a §5.3 (clases~8 y~19 de MIT 6.041~\cite{mit6041f10}).
\item Ross~\cite{ross}, cap.~5, §5.1 a §5.5, y cap.~8, §8.2 «Chebyshev's Inequality and the Weak Law of Large Numbers».
\item Lehman, Leighton y Meyer~\cite{mcs}, cap.~19 «Deviation from the Mean», §19.1 a §19.3.
\item Grinstead y Snell~\cite{gs}, §2.2 «Continuous Density Functions», §5.2 «Important Densities» y cap.~8 «Law of Large Numbers».
\end{lecturas}

\section*{Ejercicios del tema 6}
\addcontentsline{toc}{section}{Ejercicios del tema 6}

\begin{ejercicio}{1}{eje:6-conceptos}
Justifique si son verdaderas: (a) una densidad no puede tomar valores mayores que~1; (b) si $X$ es continua, $\PP(X\le 2)=\PP(X<2)$; (c) si $\PP(A)=0$, el suceso $A$ es imposible; (d) toda función de distribución es no decreciente.
\end{ejercicio}

\begin{ejercicio}{2}{eje:6-densidad}
Sea $f(x)=cx^{2}$ para $0\le x\le2$ y $f(x)=0$ en otro caso. (a) Halle $c$ para que $f$ sea una densidad. (b) Halle la función de distribución. (c) Calcule $\PP(X>1)$, $\E[X]$ y $\Var(X)$.
\end{ejercicio}

\begin{ejercicio}{2}{eje:6-pasas}
En una caja de cereales de 30\,cm de altura, la densidad de pasas a la altura $h$ (pasas por centímetro) es $40-h$. (a) ¿Cuántas pasas hay en la caja? (b) Sea $H$ la altura de una pasa elegida al azar. Halle su densidad y su función de distribución. (c) ¿Qué probabilidad hay de que esté en el tercio inferior de la caja?
\fuente{adaptado de MIT 18.05~\cite{mit1805}, \emph{Problem Set}~3, problema~4.}
\end{ejercicio}

\begin{ejercicio}{2}{eje:6-exponencial}
Sea $Y\sim\Expo(\lambda)$. (a) Obtenga la función de distribución integrando la densidad. (b) Calcule $\PP(Y\le1/\lambda)$, es decir, la probabilidad de que $Y$ no supere su media. (c) Halle la mediana de $Y$.
\fuente{adaptado de MIT 18.05~\cite{mit1805}, \emph{Problem Set}~3, problema~5(c).}
\end{ejercicio}

\begin{ejercicio}{2}{eje:6-autobus}
En la ciudad A pasa un autobús exactamente cada 10 minutos, y un viajero llega en un instante uniforme, de modo que su espera es $\Unif(0,10)$. En la ciudad B la espera es exponencial de media 10 minutos. (a) Calcule la espera media en cada ciudad. (b) En cada ciudad, si el viajero ya lleva 6 minutos esperando, ¿qué probabilidad hay de que espere al menos 3 minutos más?
\fuente{adaptado de Harvard Stat 110~\cite{stat110}, \emph{Homework}~6, problema~1, apartados (a) a~(c).}
\end{ejercicio}

\begin{ejercicio}{2}{eje:6-tipificar}
Sea $X\sim\Nor(2,3^{2})$. Usando la \cref{tab:t6-phi}, calcule $\PP(X\le2)$, $\PP(X-\mu>1{,}5\sigma)$ y $\PP(\abs{X-\mu}<1{,}5\sigma)$.
\fuente{adaptado de MIT 18.05~\cite{mit1805}, \emph{Problem Set}~3, problema~5(b).}
\end{ejercicio}

\begin{ejercicio}{2}{eje:6-gestacion}
La duración de un embarazo, en días, se modela como $\Nor(280,\,8{,}5^{2})$. (a) ¿Qué probabilidad hay de que el parto ocurra al menos una semana antes de la fecha prevista (día 280), es decir, como muy tarde el día 273? (b) ¿Qué día $d$ cumple que el parto ocurre después de $d$ con probabilidad $0{,}95$?
\fuente{adaptado de MIT 18.05~\cite{mit1805}, \emph{Problem Set}~3, problema~6 (con otras preguntas).}
\end{ejercicio}

\begin{ejercicio}{2}{eje:6-palo}
Se rompe un palo de longitud~1 por un punto $U\sim\Unif(0,1)$. Sea $L$ la longitud del trozo más largo. Halle la función de distribución de $L$, identifique su distribución y calcule $\E[L]$.
\fuente{adaptado de Harvard Stat 110~\cite{stat110}, \emph{Strategic Practice}~5, sección «Continuous Distributions», problema~4.}
\end{ejercicio}

\begin{ejercicio}{2}{eje:6-markov}
Sea $X\sim\Expo(1)$. (a) Acote $\PP(X\ge3)$ con la desigualdad de Markov y compare con el valor exacto. (b) Acote $\PP(\abs{X-1}\ge2)$ con la de Chebyshev y compare con el valor exacto.
\end{ejercicio}

\begin{ejercicio}{2}{eje:6-n-cheb}
Sean $X_1,\dots,X_n$ independientes con media $\mu$ y varianza $\sigma^{2}$. Use la desigualdad de Chebyshev para hallar un $n$ que garantice que la media muestral dista de $\mu$ menos de $2\sigma$ con probabilidad al menos $0{,}99$.
\fuente{adaptado de Harvard Stat 110~\cite{stat110}, \emph{Strategic Practice}~10, sección «Inequalities», problema~3.}
\end{ejercicio}

\begin{ejercicio}{2}{eje:6-unif-lgn}
Sean $U_1,U_2,\dots$ independientes $\Unif(-1,1)$ y $\overline U_n$ su media muestral. (a) Demuestre con la ley débil de los grandes números que $\overline U_n$ converge en probabilidad a~0. (b) Acote $\PP(\abs{\overline U_{100}}\ge0{,}1)$ con Chebyshev.
\fuente{adaptado de MIT 6.041~\cite{mit6041f10}, \emph{Problem Set}~9, problema~1(a).}
\end{ejercicio}

\begin{ejercicio}{3}{eje:6-encuesta}
En una encuesta a $n$ personas elegidas independientemente, cada una apoya una propuesta con probabilidad $p$ desconocida, y $\hat p$ es la proporción de apoyos en la muestra. (a) Demuestre que $\PP(\abs{\hat p-p}>c)\le\frac{1}{4nc^{2}}$ para todo $c>0$. (b) ¿Qué tamaño de muestra garantiza, sea cual sea $p$, un error menor que $0{,}05$ con probabilidad al menos $0{,}95$?
\fuente{adaptado de Harvard Stat 110~\cite{stat110}, \emph{Penultimate Homework}, problema~6.}
\end{ejercicio}

\begin{ejercicio}{3}{eje:6-cheb-exacta}
Dados $\mu\in\R$, $\sigma>0$ y $c\ge\sigma$, construya una variable aleatoria discreta con tres valores, media $\mu$ y desviación típica $\sigma$, para la que la desigualdad de Chebyshev $\PP(\abs{X-\mu}\ge c)\le\sigma^{2}/c^{2}$ sea una igualdad. Concluya que la cota no puede mejorarse sin más hipótesis.
\fuente{adaptado de MIT 6.041~\cite{mit6041f10}, \emph{Problem Set}~9, problema~2.}
\end{ejercicio}

\begin{ejercicio}{3}{eje:6-cantelli}
Sea $X$ con media $\mu$ y varianza $\sigma^{2}$, y $a>0$. Demuestre la desigualdad de Chebyshev unilateral (de Cantelli): $\PP(X-\mu\ge a)\le\dfrac{\sigma^{2}}{\sigma^{2}+a^{2}}$. \emph{Indicación:} para $t\ge0$, $\PP(X-\mu\ge a)\le\PP\bigl((X-\mu+t)^{2}\ge(a+t)^{2}\bigr)$; aplique Markov y elija $t$ de forma óptima.
\fuente{adaptado de MIT 6.041~\cite{mit6041f10}, \emph{Problem Set}~9, problema~G1.}
\end{ejercicio}

\begin{ejercicio}{3}{eje:6-expo}
Sea $X\sim\Expo(\lambda)$. Demuestre, integrando por partes, que $\E[X]=1/\lambda$ y $\E[X^{2}]=2/\lambda^{2}$, y deduzca $\Var(X)$. Demuestre además que $\lambda X\sim\Expo(1)$.
\end{ejercicio}

\begin{ejercicio}{3}{eje:6-afin-normal}
Sea $X\sim\Nor(\mu,\sigma^{2})$ y $a,b\in\R$ con $a\neq0$. Demuestre que $aX+b\sim\Nor(a\mu+b,\,a^{2}\sigma^{2})$, tratando por separado los casos $a>0$ y $a<0$. Deduzca que si $Z\sim\Nor(0,1)$, también $-Z\sim\Nor(0,1)$.
\end{ejercicio}


\appendix
% =====================================================================
\chapter{Soluciones de los ejercicios}
\label{ap:soluciones}
% =====================================================================

Cada solución se presenta en cinco partes: datos, planteamiento, desarrollo, resultado y comprobación. En los ejercicios de demostración, los datos son las hipótesis y el resultado es el enunciado demostrado; la comprobación consiste en un caso particular o en una verificación independiente. Todos los resultados numéricos se han comprobado con un programa, por enumeración exhaustiva cuando el espacio muestral es pequeño.

% ---------------------------------------------------------------------
\section{Tema 1}

\begin{solucion}{eje:1-reglas}
\datos cuatro situaciones de recuento.
\planteamiento «y» entre etapas sucesivas es la regla del producto; «o» entre casos disjuntos, la de la suma; «no todos» sugiere el complemento; los subconjuntos se cuentan por biyección con cadenas de bits.
\pasos
\begin{enumerate}[label=(\alph*)]
  \item Producto: $4\cdot5\cdot3=60$ menús.
  \item Suma (un plato es primero, segundo o postre, casos disjuntos): $4+5+3=12$.
  \item Complemento: $2^{10}-1=1023$ cadenas.
  \item Biyección con $\{0,1\}^{10}$ (\cref{ej:t1-bits}) y producto: $2^{10}=1024$.
\end{enumerate}
\resultado $60$; $12$; $1023$; $1024$.
\comprobacion en (c) y (d), la única cadena excluida en (c) corresponde al subconjunto vacío en (d): $1024-1=1023$.
\end{solucion}

\begin{solucion}{eje:1-vocabulario}
\datos $K=50\,000$, secuencias de longitud 3.
\planteamiento \cref{teo:t1-variaciones}: con repetición $K^{3}$, sin repetición $K^{\underline 3}$; la fracción con algún repetido se obtiene por complemento.
\pasos
\begin{enumerate}[label=\arabic*.]
  \item $K^{3}=50\,000^{3}=1{,}25\cdot10^{14}$.
  \item $K^{\underline3}=50\,000\cdot49\,999\cdot49\,998=124\,992\,500\,100\,000$.
  \item Con algún repetido: $K^{3}-K^{\underline3}=7\,499\,900\,000$.
  \item Fracción: $7\,499\,900\,000/(1{,}25\cdot10^{14})\approx5{,}99992\cdot10^{-5}$.
\end{enumerate}
\resultado (a) $1{,}25\cdot10^{14}$; (b) $124\,992\,500\,100\,000$; (c) $6{,}00\cdot10^{-5}$.
\comprobacion una secuencia tiene algún repetido si coinciden alguno de los 3 pares de posiciones; cada coincidencia tiene proporción $1/K$, y la cota de la suma da $3/K=6\cdot10^{-5}$, muy cerca del valor exacto porque las coincidencias múltiples son rarísimas.
\end{solucion}

\begin{solucion}{eje:1-contrasenas}
\datos longitudes $n\in\{6,7,8\}$, 36 caracteres posibles, de ellos 26 letras.
\planteamiento regla de la suma sobre la longitud (casos disjuntos) y, para cada longitud, complemento: todas menos las que no tienen dígitos.
\pasos
\begin{enumerate}[label=\arabic*.]
  \item $n=6$: $36^{6}-26^{6}=2\,176\,782\,336-308\,915\,776=1\,867\,866\,560$.
  \item $n=7$: $36^{7}-26^{7}=78\,364\,164\,096-8\,031\,810\,176=70\,332\,353\,920$.
  \item $n=8$: $36^{8}-26^{8}=2\,821\,109\,907\,456-208\,827\,064\,576=2\,612\,282\,842\,880$.
  \item Suma: $1\,867\,866\,560+70\,332\,353\,920+2\,612\,282\,842\,880=2\,684\,483\,063\,360$.
\end{enumerate}
\resultado $2\,684\,483\,063\,360\approx2{,}68\cdot10^{12}$ contraseñas.
\comprobacion las de longitud 8 son el $97\,\%$ del total, como cabe esperar, porque cada carácter adicional multiplica el recuento por casi 36.
\end{solucion}

\begin{solucion}{eje:1-divisibles}
\datos $U=\{1,\dots,1000\}$; $A_d$ = múltiplos de $d$, con $\abs{A_d}=\lfloor1000/d\rfloor$.
\planteamiento inclusión--exclusión para tres conjuntos (\cref{teo:t1-ie}); un entero es múltiplo de dos de los números 2, 3, 5 (primos entre sí) si y solo si lo es de su producto.
\pasos
\begin{enumerate}[label=\arabic*.]
  \item $\abs{A_2}=500$, $\abs{A_3}=333$, $\abs{A_5}=200$.
  \item $\abs{A_2\cap A_3}=\abs{A_6}=166$, $\abs{A_2\cap A_5}=\abs{A_{10}}=100$, $\abs{A_3\cap A_5}=\abs{A_{15}}=66$.
  \item $\abs{A_2\cap A_3\cap A_5}=\abs{A_{30}}=33$.
  \item $\abs{A_2\cup A_3\cup A_5}=500+333+200-166-100-66+33=734$.
  \item Ninguno: $1000-734=266$.
\end{enumerate}
\resultado 734 son divisibles por alguno y 266 por ninguno.
\comprobacion la proporción de no divisibles en cada bloque de 30 enteros consecutivos es $\frac12\cdot\frac23\cdot\frac45=\frac{8}{30}$; en 990 enteros hay $33\cdot8=264$, y entre 991 y 1000 solo 991 y 997 no son divisibles por 2, 3 ni 5: $264+2=266$.
\end{solucion}

\begin{solucion}{eje:1-palindromos}
\datos alfabeto de 26 letras.
\planteamiento un palíndromo queda determinado por su primera mitad (redondeando hacia arriba); la aplicación «palíndromo $\mapsto$ primera mitad» es una biyección.
\pasos
\begin{enumerate}[label=\arabic*.]
  \item Longitud 2: la palabra es $xx$; hay 26.
  \item Longitud 3: la palabra es $xyx$; hay $26^{2}=676$.
  \item Proporciones: $26/26^{2}=1/26$ y $26^{2}/26^{3}=1/26$: son iguales.
  \item Longitud $n$: se eligen libremente las $\lceil n/2\rceil$ primeras letras; hay $26^{\lceil n/2\rceil}$.
\end{enumerate}
\resultado 26 y 676; ambas proporciones valen $1/26$; en general, $26^{\lceil n/2\rceil}$.
\comprobacion la proporción $1/26$ en los dos casos se explica directamente: una palabra de longitud 2 o 3 es palíndromo si y solo si su última letra coincide con la primera.
\end{solucion}

\begin{solucion}{eje:1-mesa}
\datos 8 personas, disposiciones circulares módulo rotación.
\planteamiento \cref{ej:t1-circular}; para (b), se pega a Ana y Blas en un bloque.
\pasos
\begin{enumerate}[label=\arabic*.]
  \item (a) $(8-1)!=7!=5040$.
  \item (b) El bloque Ana--Blas y las otras 6 personas forman 7 unidades: $(7-1)!=720$ disposiciones circulares. Dentro del bloque hay 2 órdenes (Ana a la izquierda o a la derecha de Blas): $2\cdot720=1440$.
  \item (c) Complemento: $5040-1440=3600$.
\end{enumerate}
\resultado 5040; 1440; 3600.
\comprobacion fijada Ana, Blas ocupa uno de los 7 asientos restantes con la misma frecuencia, y 2 de ellos son contiguos a Ana: $5040\cdot\frac27=1440$.
\end{solucion}

\begin{solucion}{eje:1-cifras}
\datos números $\overline{abcd}$ con $a\neq0$ y cifras distintas.
\planteamiento regla del producto generalizada, eligiendo primero las cifras con restricciones; para los pares se separan casos según la última cifra sea 0 o no.
\pasos
\begin{enumerate}[label=\arabic*.]
  \item Total: $a$ admite 9 valores; $b$, 9 (cualquiera salvo $a$, incluido el 0); $c$, 8; $d$, 7: $9\cdot9\cdot8\cdot7=4536$.
  \item Pares con $d=0$: $a$ admite 9 valores, $b$ 8 y $c$ 7: $504$.
  \item Pares con $d\in\{2,4,6,8\}$: $d$ admite 4 valores; $a$, 8 (ni 0 ni $d$); $b$, 8; $c$, 7: $4\cdot8\cdot8\cdot7=1792$.
  \item Pares: $504+1792=2296$.
\end{enumerate}
\resultado 4536 números con cifras distintas, de los que 2296 son pares.
\comprobacion enumeración exhaustiva con un programa: 4536 y 2296.
\end{solucion}

\begin{solucion}{eje:1-funciones}
\datos $\abs{A}=4$, $\abs{B}=2$; en (d), conjuntos de 5 y 3 elementos.
\planteamiento \cref{teo:t1-variaciones}; una función no es sobreyectiva si su imagen omite algún elemento del codominio.
\pasos
\begin{enumerate}[label=(\alph*)]
  \item $2^{4}=16$.
  \item $2^{\underline4}=2\cdot1\cdot0\cdot(-1)=0$: no hay ninguna (palomar).
  \item No sobreyectivas: las constantes, 2. Sobreyectivas: $16-2=14$.
  \item Sea $E_y$ el conjunto de funciones que no alcanzan $y\in B$, $\abs{B}=3$. $\abs{E_y}=2^{5}=32$, $\abs{E_y\cap E_{y'}}=1$ y la intersección de los tres es vacía. No sobreyectivas: $3\cdot32-3\cdot1+0=93$. Sobreyectivas: $3^{5}-93=243-93=150$.
\end{enumerate}
\resultado 16; 0; 14; 150.
\comprobacion enumeración exhaustiva con un programa: 14 sobreyecciones $4\to2$ y 150 sobreyecciones $5\to3$. Además, $150=3!\cdot25$, donde 25 es el número de formas de partir un conjunto de 5 elementos en 3 bloques no vacíos: cada partición, con sus bloques asignados a los 3 elementos de $B$, da una sobreyección.
\end{solucion}

\begin{solucion}{eje:1-sombreros}
\datos $n$ personas, repartos = permutaciones de los $n$ sombreros.
\planteamiento regla del producto sobre las personas, contando cuántos sombreros admite cada una.
\pasos
\begin{enumerate}[label=(\alph*)]
  \item Solo la permutación identidad: 1 reparto.
  \item Las $m$ primeras tienen su sombrero fijado; los $n-m$ restantes se reparten libremente: $(n-m)!$.
  \item Las $m$ primeras personas reciben, en algún orden, los $m$ sombreros de las $m$ últimas: $m!$ formas (biyecciones entre dos conjuntos de $m$ elementos). Los otros $n-m$ sombreros se reparten entre las $n-m$ personas restantes: $(n-m)!$. Total, $m!\,(n-m)!$.
\end{enumerate}
\resultado $1$; $(n-m)!$; $m!\,(n-m)!$.
\comprobacion con $n=4$, $m=2$: (b) da $2$ y (c) da $2!\cdot2!=4$; enumerando las 24 permutaciones de $\{1,2,3,4\}$ se obtienen 2 con $\sigma(1)=1$, $\sigma(2)=2$ y 4 con $\{\sigma(1),\sigma(2)\}=\{3,4\}$.
\end{solucion}

\begin{solucion}{eje:1-encuesta}
\datos $\abs{U}=60$, $\abs{C}=30$, $\abs{A}=25$, $\abs{P}=20$, $\abs{C\cap A}=10$, $\abs{C\cap P}=8$, $\abs{A\cap P}=7$, $\abs{C\cap A\cap P}=3$.
\planteamiento inclusión--exclusión para la unión; para «exactamente una», se resta a cada asignatura lo compartido.
\pasos
\begin{enumerate}[label=\arabic*.]
  \item $\abs{C\cup A\cup P}=30+25+20-10-8-7+3=53$; ninguna: $60-53=7$.
  \item Solo Cálculo: $30-10-8+3=15$; solo Álgebra: $25-10-7+3=11$; solo Programación: $20-8-7+3=8$.
  \item Exactamente una: $15+11+8=34$.
\end{enumerate}
\resultado 7 no cursan ninguna y 34 cursan exactamente una.
\comprobacion exactamente dos: $(10-3)+(8-3)+(7-3)=16$; exactamente tres: 3. Total en la unión: $34+16+3=53$.
\end{solucion}

\begin{solucion}{eje:1-producto-ind}
\datos $\abs{B}=n$; $I(k)$ = número de inyecciones $[k]\to B$.
\planteamiento inducción sobre $k$, descomponiendo cada inyección en su restricción a $[k]$ y el valor en $k+1$, y usando la regla de la suma.
\pasos
\begin{enumerate}[label=\arabic*.]
  \item Caso base, $k=0$: hay una única función $\emptyset\to B$ (la vacía), que es inyectiva; $n^{\underline0}=1$.
  \item Paso: sea $0\le k<n$ y supóngase $I(k)=n^{\underline k}$. Si $f\colon[k+1]\to B$ es inyectiva, $g=f|_{[k]}$ es inyectiva y $f(k+1)\notin g([k])$. Recíprocamente, cada par $(g,b)$ con $g$ inyectiva y $b\in B\setminus g([k])$ define una única inyección $f$. Es una biyección.
  \item Para cada $g$, $\abs{B\setminus g([k])}=n-k$, porque $g$ es inyectiva y $\abs{g([k])}=k$. Por la regla de la suma sobre las $g$, $I(k+1)=I(k)(n-k)=n^{\underline k}(n-k)=n^{\underline{k+1}}$.
  \item Si $k>n$: la restricción a $[n+1]$ de una inyección $[k]\to B$ sería una inyección $[n+1]\to B$, y el paso anterior con $k=n$ da $I(n+1)=n^{\underline n}\cdot(n-n)=0$.
\end{enumerate}
\resultado hay $n^{\underline k}$ inyecciones si $k\le n$ y ninguna si $k>n$: no se pueden colocar $k$ palomas en $n<k$ nidos sin que dos compartan nido.
\comprobacion $n=3$, $k=2$: las inyecciones $[2]\to\{a,b,c\}$ son los 6 pares ordenados de elementos distintos, y $3^{\underline2}=6$.
\end{solucion}

\begin{solucion}{eje:1-fibonacci}
\datos $a_n$ = número de cadenas de $n$ bits sin \texttt{11}.
\planteamiento clasificar las cadenas válidas según su último bit.
\pasos
\begin{enumerate}[label=\arabic*.]
  \item $a_1=2$ (\texttt{0}, \texttt{1}); $a_2=3$ (\texttt{00}, \texttt{01}, \texttt{10}).
  \item Sea $n\ge3$ y $S_0$, $S_1$ las cadenas válidas de longitud $n$ que terminan en \texttt{0} y en \texttt{1}. Son disjuntas y su unión es el conjunto de cadenas válidas (regla de la suma).
  \item $w\mapsto w\texttt{0}$ es una biyección entre las cadenas válidas de longitud $n-1$ y $S_0$: añadir un \texttt{0} no crea \texttt{11}, y quitar el \texttt{0} final de una cadena válida da una válida. Luego $\abs{S_0}=a_{n-1}$.
  \item Si una cadena válida termina en \texttt{1}, su penúltimo bit es \texttt{0}. Por tanto $w\mapsto w\texttt{01}$ es una biyección entre las cadenas válidas de longitud $n-2$ y $S_1$. Luego $\abs{S_1}=a_{n-2}$.
  \item $a_n=a_{n-1}+a_{n-2}$; sucesivamente $a_3=5$, $a_4=8$, $a_5=13$, $a_6=21$, $a_7=34$, $a_8=55$.
\end{enumerate}
\resultado la recurrencia de Fibonacci con $a_1=2$, $a_2=3$, y $a_8=55$.
\comprobacion enumeración exhaustiva de las 256 cadenas de 8 bits: 55 no contienen \texttt{11}.
\end{solucion}

\begin{solucion}{eje:1-ie-cuatro}
\datos $A,B,C,D$ finitos; \cref{teo:t1-ie}.
\planteamiento aplicar la fórmula de dos conjuntos a $X=A\cup B\cup C$ y $D$, y la de tres conjuntos dos veces.
\pasos
\begin{enumerate}[label=\arabic*.]
  \item $\abs{X\cup D}=\abs X+\abs D-\abs{X\cap D}$.
  \item Por la distributiva, $X\cap D=(A\cap D)\cup(B\cap D)\cup(C\cap D)$, y por la fórmula de tres conjuntos, con $(A\cap D)\cap(B\cap D)=A\cap B\cap D$, etc.:
  \[
  \abs{X\cap D}=\abs{A\cap D}+\abs{B\cap D}+\abs{C\cap D}-\abs{A\cap B\cap D}-\abs{A\cap C\cap D}-\abs{B\cap C\cap D}+\abs{A\cap B\cap C\cap D}.
  \]
  \item $\abs X$ se desarrolla con la fórmula de tres conjuntos.
  \item Sustituyendo: suma de los 4 cardinales, menos los $\binom42=6$ de las intersecciones dobles ($AB$, $AC$, $BC$ de $\abs X$ y $AD$, $BD$, $CD$ de $\abs{X\cap D}$), más los $\binom43=4$ de las triples ($ABC$ de $\abs X$ y $ABD$, $ACD$, $BCD$), menos la cuádruple.
\end{enumerate}
\resultado $\abs{A\cup B\cup C\cup D}=\sum\abs{A_i}-\sum_{i<j}\abs{A_i\cap A_j}+\sum_{i<j<k}\abs{A_i\cap A_j\cap A_k}-\abs{A\cap B\cap C\cap D}$.
\comprobacion con $A=B=C=D=\{1\}$: $4-6+4-1=1=\abs{\{1\}}$.
\end{solucion}

\begin{solucion}{eje:1-collares}
\datos (a) $n\ge3$ personas en una mesa redonda, módulo rotación y reflexión; (b) $2n$ personas en $n$ parejas.
\planteamiento regla de la división en ambos casos, identificando cuántos objetos «ordenados» dan el mismo objeto «sin orden».
\pasos
\begin{enumerate}[label=\arabic*.]
  \item (a) Sea $f$ la aplicación que lleva cada disposición circular (de las $(n-1)!$) a su clase módulo reflexión. Cada clase contiene una disposición y su reflejada.
  \item Ambas son distintas si $n\ge3$: en la reflejada, el vecino de la derecha de una persona $x$ es el que en la original era su vecino de la izquierda, y para $n\ge3$ los dos vecinos de $x$ son personas distintas.
  \item $f$ es 2 a~1, luego hay $(n-1)!/2$ clases.
  \item (b) Sea $g$ la aplicación que a cada permutación $(a_1,\dots,a_{2n})$ asocia el emparejamiento $\{\{a_1,a_2\},\dots,\{a_{2n-1},a_{2n}\}\}$; es sobreyectiva.
  \item Un emparejamiento dado se obtiene de las permutaciones que ordenan sus $n$ parejas de cualquier forma ($n!$) y, dentro de cada pareja, ponen primero cualquiera de los dos ($2^{n}$): $g$ es $2^{n}n!$ a~1.
  \item Por la regla de la división, hay $(2n)!/(2^{n}n!)$ emparejamientos.
\end{enumerate}
\resultado $(n-1)!/2$ y $(2n)!/(2^{n}n!)$.
\comprobacion (a) con $n=3$: $2!/2=1$, pues tres personas en círculo, sin distinguir sentido, están siempre en la misma disposición. (b) con $n=2$: $4!/(4\cdot2)=3$, y en efecto $\{a,b,c,d\}$ admite los emparejamientos $ab|cd$, $ac|bd$, $ad|bc$.
\end{solucion}

% ---------------------------------------------------------------------
\section{Tema 2}

\begin{solucion}{eje:2-calcular}
\planteamiento \cref{teo:t2-binom} y \cref{teo:t2-binomio}.
\pasos
\begin{enumerate}[label=\arabic*.]
  \item $\binom{10}{3}=\frac{10\cdot9\cdot8}{3!}=120$ y $\binom{10}{7}=\frac{10!}{7!\,3!}=120$.
  \item $\binom{52}{5}=\frac{52\cdot51\cdot50\cdot49\cdot48}{120}=2\,598\,960$.
  \item El coeficiente de $x^{3}$ en $(1+x)^{7}$ es $\binom73=35$.
  \item Elegir los 3 elementos que \emph{entran} equivale a elegir los 7 que \emph{quedan fuera} (\cref{prop:t2-props}(i)).
\end{enumerate}
\resultado $120$, $120$, $2\,598\,960$ y $35$.
\comprobacion fila 7 del triángulo de Pascal: $1,7,21,35,35,21,7,1$, que suma $128=2^{7}$.
\end{solucion}

\begin{solucion}{eje:2-modelo}
\planteamiento decidir si importa el orden y si se permite repetir (\cref{tab:t2-modelos}).
\pasos
\begin{enumerate}[label=(\alph*)]
  \item Ordenado, sin repetición: $8^{\underline3}=336$.
  \item No ordenado, sin repetición: $\binom83=56$.
  \item Ordenado, con repetición: $8^{3}=512$.
  \item No ordenado, con repetición: $\binom{8+3-1}{3}=\binom{10}{3}=120$.
\end{enumerate}
\resultado $336$, $56$, $512$ y $120$.
\comprobacion $336=56\cdot3!$: cada comisión da $3!$ podios. En cambio $512/3!$ no es entero, lo que confirma que (d) no puede obtenerse dividiendo (c) por $3!$.
\end{solucion}

\begin{solucion}{eje:2-poker}
\datos 52 cartas, 13 valores y 4 palos; $\binom{52}5=2\,598\,960$ manos.
\planteamiento construir cada mano por etapas (valores, después palos) de modo que cada mano se obtenga exactamente una vez~\cite[§14.7]{mcs}.
\pasos
\begin{enumerate}[label=(\alph*)]
  \item Póquer: valor del cuarteto (13) y quinta carta (48 restantes): $13\cdot48=624$.
  \item Full: valor del trío (13), sus palos $\binom43=4$, valor de la pareja (12), sus palos $\binom42=6$: $13\cdot4\cdot12\cdot6=3744$.
  \item Doble pareja: los dos valores de las parejas, sin orden, $\binom{13}{2}=78$; palos de cada pareja, $6\cdot6=36$; quinta carta de otro valor, $52-8=44$: $78\cdot36\cdot44=123\,552$.
  \item Fracciones: $624/2\,598\,960=1/4165\approx0{,}00024$; $3744/2\,598\,960=6/4165\approx0{,}00144$; $123\,552/2\,598\,960=198/4165\approx0{,}0475$.
\end{enumerate}
\resultado 624, 3744 y 123\,552 manos.
\comprobacion si en (c) se eligieran los valores de las parejas en orden ($13\cdot12$) se obtendría el doble, $247\,104$: cada mano se contaría dos veces.
\end{solucion}

\begin{solucion}{eje:2-equipos}
\datos 12 personas.
\planteamiento \cref{teo:t2-multinomial} para equipos con nombre y regla de la división para los equipos del mismo tamaño, que son indistinguibles.
\pasos
\begin{enumerate}[label=\arabic*.]
  \item (a) Con nombres (el de 2, el «primero» de 5 y el «segundo» de 5): $\binom{12}{2,5,5}=\frac{12!}{2!\,5!\,5!}=16\,632$.
  \item Los dos equipos de 5 no tienen nombre: cada reparto se ha contado $2!$ veces. Resultado: $16\,632/2=8316$.
  \item (b) Con nombres: $\frac{12!}{4!\,4!\,4!}=34\,650$; sin nombres, se divide por $3!$: $5775$.
\end{enumerate}
\resultado (a) 8316; (b) 5775.
\comprobacion (b) por etapas: la persona~1 elige 3 compañeros entre 11, $\binom{11}{3}=165$; la persona de menor número entre las restantes elige 3 entre 7, $\binom73=35$; el resto forma el tercer equipo: $165\cdot35=5775$.
\end{solucion}

\begin{solucion}{eje:2-anagramas}
\datos \texttt{MISSISSIPPI}: M~1, I~4, S~4, P~2 (11 letras). \texttt{PROBABILIDAD}: B, A, I, D dos veces cada una y P, R, O, L una vez (12 letras).
\planteamiento \cref{teo:t2-multinomial}; para las \texttt{S} juntas, se trata el bloque \texttt{SSSS} como un solo símbolo.
\pasos
\begin{enumerate}[label=\arabic*.]
  \item $\frac{11!}{1!\,4!\,4!\,2!}=\frac{39\,916\,800}{1152}=34\,650$.
  \item $\frac{12!}{2!^{4}}=\frac{479\,001\,600}{16}=29\,937\,600$.
  \item Con el bloque: símbolos M, I~$\times4$, P~$\times2$ y el bloque: 8 símbolos, $\frac{8!}{4!\,2!}=840$.
\end{enumerate}
\resultado 34\,650; 29\,937\,600; 840.
\comprobacion la proporción de anagramas con las \texttt{S} juntas, $840/34\,650\approx0{,}024$, coincide con la probabilidad de que las 4 posiciones de las \texttt{S}, elegidas al azar entre 11, sean consecutivas: $8/\binom{11}{4}=8/330\approx0{,}024$.
\end{solucion}

\begin{solucion}{eje:2-coeficientes}
\planteamiento binomio con $x\to2x$, $y\to-y$; multinomio de Newton; los términos de $(x+y+z)^{6}$ corresponden a las soluciones de $k_1+k_2+k_3=6$.
\pasos
\begin{enumerate}[label=\arabic*.]
  \item Término general de $(2x-y)^{8}$: $\binom8k(2x)^{k}(-y)^{8-k}$. Con $k=3$: $\binom83\,2^{3}(-1)^{5}=56\cdot8\cdot(-1)=-448$.
  \item $\binom{6}{2,3,1}=\frac{720}{2\cdot6\cdot1}=60$.
  \item Número de términos: soluciones de $k_1+k_2+k_3=6$ con $k_i\ge0$, $\binom{8}{2}=28$ (\cref{teo:t2-estrellas}).
\end{enumerate}
\resultado $-448$; $60$; 28 términos.
\comprobacion la suma de los coeficientes de $(x+y+z)^{6}$ es $3^{6}=729$ (tomar $x=y=z=1$), compatible con 28 términos cuyo mayor coeficiente es $\binom{6}{2,2,2}=90$.
\end{solucion}

\begin{solucion}{eje:2-ecuaciones}
\planteamiento \cref{teo:t2-estrellas} con $n=4$ incógnitas y cambio de variable para las cotas inferiores.
\pasos
\begin{enumerate}[label=(\alph*)]
  \item $\binom{10+3}{3}=\binom{13}3=286$.
  \item $\binom{10-1}{3}=\binom93=84$.
  \item Con $y_1=x_1-2\ge0$: $y_1+x_2+x_3+x_4=8$, $\binom{11}{3}=165$.
\end{enumerate}
\resultado 286, 84 y 165.
\comprobacion por complemento, las soluciones de (a) con $x_1\le1$ son las de $x_2+x_3+x_4=10$ y las de $x_2+x_3+x_4=9$: $\binom{12}2+\binom{11}2=66+55=121$, y $286-121=165$.
\end{solucion}

\begin{solucion}{eje:2-caminos}
\planteamiento \cref{ej:t2-caminos}: de $(0,0)$ a $(a,b)$ hay $\binom{a+b}{a}$ caminos; pasar por un punto intermedio divide el camino en dos tramos independientes (regla del producto).
\pasos
\begin{enumerate}[label=\arabic*.]
  \item (a) $\binom{11}{5}=462$.
  \item (b) De $(0,0)$ a $(2,3)$: $\binom52=10$; de $(2,3)$ a $(5,6)$, es decir, 3 a la derecha y 3 arriba: $\binom63=20$. Total, $200$.
  \item (c) $462-200=262$.
\end{enumerate}
\resultado 462; 200; 262.
\comprobacion la proporción $200/462\approx0{,}43$ es la probabilidad de que un camino uniforme pase por $(2,3)$; es razonable que sea alta porque $(2,3)$ está cerca de la diagonal que va a $(5,6)$.
\end{solucion}

\begin{solucion}{eje:2-estaciones}
\datos 7 personas, 4 estaciones; $4^{7}=16\,384$ asignaciones.
\planteamiento complemento e inclusión--exclusión (\cref{teo:t2-ie}) sobre los sucesos $E_s$ = «nadie nace en la estación $s$».
\pasos
\begin{enumerate}[label=\arabic*.]
  \item $\abs{E_s}=3^{7}=2187$; $\abs{E_s\cap E_t}=2^{7}=128$; tres estaciones vacías: $1^{7}=1$; cuatro: 0.
  \item $\abs{\bigcup E_s}=4\cdot2187-6\cdot128+4\cdot1-0=8748-768+4=7984$.
  \item Con las cuatro estaciones: $16\,384-7984=8400$.
  \item Fracción: $8400/16\,384\approx0{,}5127$.
\end{enumerate}
\resultado 8400 asignaciones, un $51{,}3\,\%$ del total.
\comprobacion es el número de sobreyecciones de un conjunto de 7 en uno de 4, que una enumeración con un programa confirma: 8400.
\end{solucion}

\begin{solucion}{eje:2-desarreglos}
\datos $n=5$.
\planteamiento \cref{ej:t2-desarreglos}.
\pasos
\begin{enumerate}[label=\arabic*.]
  \item $D_5=5!\bigl(1-1+\tfrac12-\tfrac16+\tfrac1{24}-\tfrac1{120}\bigr)=60-20+5-1=44$.
  \item Fracción: $44/120\approx0{,}3667$; $1/e\approx0{,}3679$.
\end{enumerate}
\resultado 44 sorteos válidos, un $36{,}7\,\%$, a menos de $0{,}002$ de $1/e$.
\comprobacion recurrencia $D_n=(n-1)(D_{n-1}+D_{n-2})$ con $D_1=0$, $D_2=1$: $D_3=2$, $D_4=9$, $D_5=4(9+2)=44$.
\end{solucion}

\begin{solucion}{eje:2-ositos}
\datos 5 sabores, $n$ ositos con $30\le n\le50$.
\planteamiento para cada $n$, una bolsa es una solución de $x_1+\dots+x_5=n$ con $x_i\ge0$: $\binom{n+4}{4}$ bolsas. Se suma sobre $n$ y se usa la identidad del palo de \emph{hockey} (\cref{eje:2-palo}).
\pasos
\begin{enumerate}[label=\arabic*.]
  \item Total: $\sum_{n=30}^{50}\binom{n+4}{4}=\sum_{j=34}^{54}\binom j4$.
  \item Por el \cref{eje:2-palo}, $\sum_{j=4}^{54}\binom j4=\binom{55}5$ y $\sum_{j=4}^{33}\binom j4=\binom{34}5$.
  \item Restando: $\binom{55}{5}-\binom{34}{5}=3\,478\,761-278\,256=3\,200\,505$.
\end{enumerate}
\resultado $\binom{55}5-\binom{34}5=3\,200\,505$ bolsas.
\comprobacion otra lectura: añadiendo un sexto «sabor» ficticio $x_6=50-n$, las bolsas con a lo sumo 50 ositos son las soluciones de $x_1+\dots+x_6=50$, $\binom{55}{5}$, y las de a lo sumo 29 son $\binom{34}{5}$; la suma directa con un programa da el mismo valor.
\end{solucion}

\begin{solucion}{eje:2-palo}
\datos enteros $n\ge k\ge0$.
\planteamiento historia: contar los grupos de $k+1$ personas elegidas entre $n+1$ personas de edades distintas, clasificándolos por la persona de mayor edad del grupo.
\pasos
\begin{enumerate}[label=\arabic*.]
  \item Se numeran las personas $1,\dots,n+1$ de menor a mayor edad. El lado derecho, $\binom{n+1}{k+1}$, es el número de grupos de $k+1$.
  \item Si la persona de más edad del grupo es la $j+1$, los otros $k$ miembros se eligen entre las $j$ personas más jóvenes que ella: $\binom jk$ grupos. Para que sea posible, $j\ge k$, y como $j+1\le n+1$, $j\le n$.
  \item Cada grupo tiene exactamente una persona de mayor edad, así que las clases $j=k,\dots,n$ son disjuntas y cubren todos los grupos. Por la regla de la suma, $\binom{n+1}{k+1}=\sum_{j=k}^{n}\binom jk$.
\end{enumerate}
\resultado $\sum_{j=k}^{n}\binom jk=\binom{n+1}{k+1}$.
\comprobacion $k=1$, $n=4$: $1+2+3+4=10=\binom52$.
\end{solucion}

\begin{solucion}{eje:2-sumak}
\datos $n\ge1$ (en (b), $n\ge2$).
\planteamiento (a) absorción y \cref{cor:t2-sumas}; (b) doble conteo de comités con presidente y vicepresidente.
\pasos
\begin{enumerate}[label=\arabic*.]
  \item (a) Para $k\ge1$, $k\binom nk=n\binom{n-1}{k-1}$; el término $k=0$ es nulo. Luego $\sum_kk\binom nk=n\sum_{k=1}^{n}\binom{n-1}{k-1}=n\,2^{n-1}$.
  \item (b) Se cuentan las ternas (comité $C\subseteq[n]$, presidente $p\in C$, vicepresidente $v\in C$, $v\neq p$). Por tamaño del comité: si $\abs C=k$, hay $\binom nk$ comités y $k(k-1)$ elecciones de $(p,v)$; total $\sum_{k=2}^{n}k(k-1)\binom nk$.
  \item Eligiendo primero los cargos: $n$ presidentes, $n-1$ vicepresidentes, y el resto del comité es un subconjunto cualquiera de las $n-2$ personas restantes, $2^{n-2}$: total $n(n-1)2^{n-2}$.
\end{enumerate}
\resultado $\sum_kk\binom nk=n2^{n-1}$ y $\sum_kk(k-1)\binom nk=n(n-1)2^{n-2}$.
\comprobacion $n=3$: $1\cdot3+2\cdot3+3\cdot1=12=3\cdot2^{2}$ y $2\cdot1\cdot3+3\cdot2\cdot1=12=3\cdot2\cdot2^{1}$.
\end{solucion}

\begin{solucion}{eje:2-vandermonde}
\datos $m$ mujeres, $n$ hombres, comités de $k$ personas.
\planteamiento doble conteo: total y clasificación por el número $j$ de mujeres.
\pasos
\begin{enumerate}[label=\arabic*.]
  \item Hay $\binom{m+n}{k}$ comités de $k$ personas.
  \item Los comités con exactamente $j$ mujeres se forman eligiendo $j$ mujeres y $k-j$ hombres: $\binom mj\binom n{k-j}$ (con el convenio $\binom ab=0$ si $b>a$ o $b<0$).
  \item Las clases $j=0,\dots,k$ son disjuntas y lo cubren todo; por la regla de la suma, $\binom{m+n}k=\sum_j\binom mj\binom n{k-j}$.
  \item Con $m=n=k$: $\binom{2n}{n}=\sum_j\binom nj\binom n{n-j}=\sum_j\binom nj^{2}$, por la simetría.
\end{enumerate}
\resultado la identidad de Vandermonde y $\sum_j\binom nj^{2}=\binom{2n}n$.
\comprobacion $n=3$: $1+9+9+1=20=\binom63$.
\end{solucion}

\begin{solucion}{eje:2-pares}
\datos $\abs A=n\ge1$, $a\in A$; $\varphi(X)=X\dsim\{a\}$.
\planteamiento probar que $\varphi$ es una involución que cambia la paridad del cardinal.
\pasos
\begin{enumerate}[label=\arabic*.]
  \item Si $a\notin X$, $\varphi(X)=X\cup\{a\}$ tiene $\abs X+1$ elementos; si $a\in X$, $\varphi(X)=X\setminus\{a\}$ tiene $\abs X-1$. En ambos casos cambia la paridad.
  \item $\varphi(\varphi(X))=X$: añadir y quitar $a$ (o al revés) devuelve $X$. Luego $\varphi$ es biyectiva, con inversa ella misma.
  \item Restringida a los subconjuntos de cardinal par, $\varphi$ toma valores en los de cardinal impar y es biyectiva sobre ellos (su inversa es $\varphi$ restringida a los impares). Hay tantos subconjuntos pares como impares.
  \item $\sum_k(-1)^{k}\binom nk=\#\text{pares}-\#\text{impares}=0$. Como $\#\text{pares}+\#\text{impares}=2^{n}$, hay $2^{n-1}$ de cada clase.
\end{enumerate}
\resultado $\sum_k(-1)^{k}\binom nk=0$ y hay $2^{n-1}$ subconjuntos de cardinal par.
\comprobacion $n=3$: pares $\emptyset$, $\{a,b\}$, $\{a,c\}$, $\{b,c\}$, que son 4 $=2^{2}$.
\end{solucion}

% ---------------------------------------------------------------------
\section{Tema 3}

\begin{solucion}{eje:3-espacio}
\planteamiento elegir resultados lo bastante detallados para describir el experimento y, si es posible, equiprobables.
\pasos
\begin{enumerate}[label=(\alph*)]
  \item $\Omega=\{\mathrm C,\mathrm X\}^{3}$, con $2^{3}=8$ resultados; Laplace es razonable si la moneda es equilibrada.
  \item $\Omega$ = subconjuntos de 5 cartas, $\binom{52}5=2\,598\,960$; Laplace es razonable si la baraja está bien mezclada.
  \item $\Omega=\V^{2}$, con $K^{2}$ resultados; Laplace \emph{no} es razonable: el modelo asigna probabilidades muy distintas a unas secuencias y otras.
  \item $\Omega=\{1,2,3,\dots\}$, infinito numerable; Laplace no es aplicable (no hay ley uniforme sobre un conjunto infinito numerable, \cref{eje:3-laplace}), y además $\PP(\{k\})=2^{-k}$ con una moneda equilibrada.
\end{enumerate}
\resultado $8$; $2\,598\,960$; $K^{2}$; infinito numerable. Laplace solo en (a) y (b).
\comprobacion en (d), $\sum_{k\ge1}2^{-k}=1$, de modo que la ley indicada es válida (\cref{prop:t3-discreto}).
\end{solucion}

\begin{solucion}{eje:3-sucesos}
\planteamiento traducir «o», «y» y «no» en unión, intersección y complementario; «como máximo uno» es el complementario de «al menos dos».
\pasos
\begin{enumerate}[label=(\alph*)]
  \item $A\cup B\cup C$.
  \item $A\cap B\cap C$.
  \item $A^{c}\cap B^{c}\cap C^{c}=(A\cup B\cup C)^{c}$.
  \item $(A\cap B^{c}\cap C^{c})\cup(A^{c}\cap B\cap C^{c})\cup(A^{c}\cap B^{c}\cap C)$.
  \item $(A^{c}\cap B^{c})\cup(A^{c}\cap C^{c})\cup(B^{c}\cap C^{c})$: entre cada dos, al menos uno no ocurre.
  \item $A\cap B\cap C^{c}$.
\end{enumerate}
\resultado las expresiones anteriores.
\comprobacion (e) es el complementario de «al menos dos», $(A\cap B)\cup(A\cap C)\cup(B\cap C)$, por las leyes de De Morgan; y es la unión disjunta de (c) y (d).
\end{solucion}

\begin{solucion}{eje:3-monedas}
\datos $\Omega=\{\mathrm C,\mathrm X\}^{3}$, 8 resultados equiprobables.
\planteamiento modelo de Laplace.
\pasos
\begin{enumerate}[label=(\alph*)]
  \item $\{\mathrm{CCC}\}$: $1/8$.
  \item $\{\mathrm{CXC}\}$: $1/8$.
  \item Exactamente dos caras: $\binom32=3$ resultados, $3/8$.
  \item Al menos tantas caras como cruces equivale a 2 o 3 caras: $3+1=4$ resultados, $1/2$.
\end{enumerate}
\resultado $1/8$, $1/8$, $3/8$, $1/2$.
\comprobacion (d) vale $1/2$ por simetría: intercambiar caras y cruces lleva «más caras» a «más cruces», y con 3 lanzamientos no puede haber empate.
\end{solucion}

\begin{solucion}{eje:3-dados}
\datos 36 resultados equiprobables $(i,j)$.
\planteamiento Laplace; complemento en (b); en (c), $\{\max=4\}=\{\max\le4\}\setminus\{\max\le3\}$.
\pasos
\begin{enumerate}[label=(\alph*)]
  \item Suma 7: 6 pares, $6/36=1/6$.
  \item Ningún 6: $5^{2}=25$; al menos un 6: $1-25/36=11/36$.
  \item $\abs{\{\max\le4\}}=16$, $\abs{\{\max\le3\}}=9$: $(16-9)/36=7/36$.
\end{enumerate}
\resultado $1/6$, $11/36$, $7/36$.
\comprobacion $\frac16+\frac16=\frac{12}{36}$ cuenta dos veces el resultado $(6,6)$, que está en ambos sucesos: $\frac{12}{36}-\frac1{36}=\frac{11}{36}$ (\cref{prop:t3-props}(vi)).
\end{solucion}

\begin{solucion}{eje:3-cuatro-dados}
\datos 4 dados, $6^{4}=1296$ resultados.
\planteamiento la aplicación $(x_1,\dots,x_4)\mapsto(7-x_1,\dots,7-x_4)$ es una biyección de $\Omega$ que lleva suma $s$ a suma $28-s$.
\pasos
\begin{enumerate}[label=\arabic*.]
  \item Por la biyección, $\PP(S=21)=\PP(S=7)$ y $\PP(S=22)=\PP(S=6)$.
  \item Sumas pequeñas: con $y_i=x_i-1\ge0$, la suma 7 equivale a $y_1+\dots+y_4=3$, con $\binom63=20$ soluciones (todas con $y_i\le5$); la suma 6 equivale a $y_1+\dots+y_4=2$, con $\binom53=10$.
  \item $\PP(S=21)=20/1296\approx0{,}0154$ y $\PP(S=22)=10/1296\approx0{,}0077$.
\end{enumerate}
\resultado es mayor la probabilidad de 21, el doble que la de 22.
\comprobacion enumeración de los 1296 resultados: 20 con suma 21 y 10 con suma 22.
\end{solucion}

\begin{solucion}{eje:3-horarios}
\datos 3 asignaturas, 10 franjas; $10^{3}$ asignaciones equiprobables.
\planteamiento complemento: todas en franjas distintas.
\pasos
\begin{enumerate}[label=\arabic*.]
  \item Franjas distintas: $10\cdot9\cdot8=720$.
  \item $\PP(\text{coincidencia})=1-720/1000=0{,}28$.
\end{enumerate}
\resultado $7/25=0{,}28$.
\comprobacion por la cota de la unión sobre los 3 pares, $\PP\le3/10$; el valor exacto es algo menor porque se descuentan las coincidencias triples.
\end{solucion}

\begin{solucion}{eje:3-robos}
\datos 6 robos, 6 distritos, $6^{6}=46\,656$ asignaciones equiprobables.
\planteamiento complemento: ningún distrito repite, es decir, la asignación es una biyección.
\pasos
\begin{enumerate}[label=\arabic*.]
  \item Asignaciones sin repetición: $6!=720$.
  \item $\PP=1-720/46\,656=1-5/324=319/324$.
\end{enumerate}
\resultado $319/324\approx0{,}9846$.
\comprobacion es el problema del cumpleaños (\cref{prop:t3-cumple}) con $N=n=6$: $1-\prod_{i=0}^{5}(1-i/6)=1-\frac{6!}{6^{6}}$.
\end{solucion}

\begin{solucion}{eje:3-seis-cartas}
\datos manos de 6 cartas, $\binom{52}6=20\,358\,520$.
\planteamiento construir cada mano por valores y palos, como en el \cref{eje:2-poker}.
\pasos
\begin{enumerate}[label=\arabic*.]
  \item Doble pareja: valores de las parejas $\binom{13}2=78$; palos $\binom42^{2}=36$; valores de las dos cartas sueltas, distintos entre sí y de los anteriores, $\binom{11}2=55$; palos de esas dos, $4^{2}=16$. Total: $78\cdot36\cdot55\cdot16=2\,471\,040$.
  \item Trío: valor del trío 13, palos $\binom43=4$; valores de las tres sueltas $\binom{12}3=220$, palos $4^{3}=64$. Total: $13\cdot4\cdot220\cdot64=732\,160$.
  \item Probabilidades: $2\,471\,040/20\,358\,520\approx0{,}1214$ y $732\,160/20\,358\,520\approx0{,}0360$.
\end{enumerate}
\resultado la doble pareja es más probable: $0{,}1214$ frente a $0{,}0360$.
\comprobacion un programa que cuenta manos por su patrón de multiplicidades de valores, $(2,2,1,1)$ y $(3,1,1,1)$, da los mismos dos números.
\end{solucion}

\begin{solucion}{eje:3-cumple}
\datos $n$ personas además de usted; 365 días.
\planteamiento complemento: ninguna de las $n$ cumple el día de usted, cada una con probabilidad $364/365$, de forma independiente.
\pasos
\begin{enumerate}[label=\arabic*.]
  \item $\PP(A_n)=1-(364/365)^{n}$.
  \item $\PP(A_n)>1/2\iff n>\ln2/(-\ln(364/365))\approx252{,}65$. Luego $n=253$: $\PP(A_{252})\approx0{,}4991$ y $\PP(A_{253})\approx0{,}5005$.
  \item (b) En el problema clásico cuenta cualquiera de los $\binom{n}{2}$ pares; aquí solo los $n$ pares que le incluyen a usted. Por la cota de la unión, $\PP(A_n)\le n/365$, luego hace falta $n\ge183>365/2$; y $n$ debe ser aún mayor porque las otras personas coinciden a veces entre sí y cubren menos de $n$ días distintos.
\end{enumerate}
\resultado $\PP(A_n)=1-(364/365)^{n}$ y $n=253$.
\comprobacion con $n=253$ hay tantos pares que le incluyen como pares hay en un grupo de 23 personas, $\binom{23}2=253$; por eso ambos problemas tienen probabilidades cercanas a $1/2$.
\end{solucion}

\begin{solucion}{eje:3-urna}
\datos $r$ rojas, $v$ verdes, dos extracciones sin reemplazamiento.
\planteamiento resultados ordenados $(b_1,b_2)$ de bolas distintas, todos equiprobables: $(r+v)(r+v-1)$.
\pasos
\begin{enumerate}[label=\arabic*.]
  \item Primera verde: $v(r+v-1)$ resultados; probabilidad $v/(r+v)$.
  \item Segunda verde: primera roja y segunda verde, $rv$; o ambas verdes, $v(v-1)$. Total $v(r+v-1)$: misma probabilidad $v/(r+v)$.
  \item (b) Mismo color: $\frac{r(r-1)+v(v-1)}{16\cdot15}=\frac12$, luego $r^{2}+v^{2}-16=120$ y $r^{2}+v^{2}=136$.
  \item Con $v=16-r$: $r^{2}+(16-r)^{2}=136\iff r^{2}-16r+60=0\iff r\in\{6,10\}$.
\end{enumerate}
\resultado (a) ambas valen $v/(r+v)$; (b) $(r,v)=(6,10)$ o $(10,6)$.
\comprobacion con $r=6$, $v=10$: $\frac{30+90}{240}=\frac12$.
\end{solucion}

\begin{solucion}{eje:3-geometrica}
\datos $(x,y)$ uniforme en $[0,1]^{2}$.
\planteamiento calcular áreas de regiones del cuadrado.
\pasos
\begin{enumerate}[label=\arabic*.]
  \item $\{y>x+1/3\}$ es el triángulo de vértices $(0,1/3)$, $(0,1)$, $(2/3,1)$, de catetos $2/3$: área $\frac12\cdot\frac49=\frac29$. Por simetría, $\{x>y+1/3\}$ tiene la misma área. Total: $4/9$.
  \item $\{x+y<1/2\}$ es el triángulo de vértices $(0,0)$, $(1/2,0)$, $(0,1/2)$: área $\frac12\cdot\frac14=\frac18$.
\end{enumerate}
\resultado $4/9$ y $1/8$.
\comprobacion el mismo razonamiento con un umbral $d\in[0,1]$ da $\PP(\abs{x-y}>d)=(1-d)^{2}$, que vale 1 para $d=0$ y 0 para $d=1$, como debe ser; con $d=1/3$ se recupera $4/9$.
\end{solucion}

\begin{solucion}{eje:3-bonferroni}
\datos sucesos $A$, $B$, $C$.
\planteamiento \cref{prop:t3-props}(vi) y (v); después, aplicar el caso de dos sucesos dos veces.
\pasos
\begin{enumerate}[label=\arabic*.]
  \item $\PP(A\cap B)=\PP(A)+\PP(B)-\PP(A\cup B)\ge\PP(A)+\PP(B)-1$, porque $\PP(A\cup B)\le1$.
  \item Con $A\cap B$ y $C$: $\PP(A\cap B\cap C)\ge\PP(A\cap B)+\PP(C)-1$.
  \item Sustituyendo el paso 1: $\PP(A\cap B\cap C)\ge\PP(A)+\PP(B)+\PP(C)-2$.
\end{enumerate}
\resultado las dos desigualdades.
\comprobacion si $\PP(A)=\PP(B)=\PP(C)=0{,}9$, la cota garantiza $\PP(A\cap B\cap C)\ge0{,}7$; con $A$, $B$, $C$ independientes el valor sería $0{,}729$, compatible con la cota.
\end{solucion}

\begin{solucion}{eje:3-interseccion}
\datos sucesos $A_1,\dots,A_n$.
\planteamiento De Morgan y la cota de la unión (\cref{teo:t3-boole}).
\pasos
\begin{enumerate}[label=\arabic*.]
  \item $\bigl(\bigcap_iA_i\bigr)^{c}=\bigcup_iA_i^{c}$, luego $\PP\bigl(\bigcap_iA_i\bigr)=1-\PP\bigl(\bigcup_iA_i^{c}\bigr)$.
  \item Por la cota de la unión, $\PP\bigl(\bigcup_iA_i^{c}\bigr)\le\sum_i\PP(A_i^{c})$, de donde la desigualdad.
  \item Aplicación: $A_i$ = «el componente $i$ no falla», $\PP(A_i^{c})=0{,}001$. Entonces $\PP(\text{ninguno falla})\ge1-20\cdot0{,}001=0{,}98$.
\end{enumerate}
\resultado $\PP\bigl(\bigcap A_i\bigr)\ge1-\sum\PP(A_i^{c})$; sin hipótesis de independencia, la probabilidad de que no falle ninguno es al menos $0{,}98$.
\comprobacion con independencia sería $0{,}999^{20}\approx0{,}9802$, apenas por encima de la cota.
\end{solucion}

\begin{solucion}{eje:3-distancia}
\datos sucesos $A,B,C$; $d(A,B)=\PP(A\dsim B)$.
\planteamiento (a) descomponer en partes disjuntas; (b) probar $A\dsim C\subseteq(A\dsim B)\cup(B\dsim C)$ y aplicar monotonía y la cota de la unión.
\pasos
\begin{enumerate}[label=\arabic*.]
  \item (a) $A\dsim B$ es la unión disjunta de $A\setminus B$ y $B\setminus A$, y $\PP(A\setminus B)=\PP(A)-\PP(A\cap B)$ por la \cref{prop:t3-props}(iv) aplicada a $A\cap B\subseteq A$. Análogamente para $B\setminus A$. Sumando: $\PP(A)+\PP(B)-2\PP(A\cap B)$.
  \item (b) Sea $\omega\in A\setminus C$. Si $\omega\in B$, entonces $\omega\in B\setminus C\subseteq B\dsim C$; si $\omega\notin B$, entonces $\omega\in A\setminus B\subseteq A\dsim B$. El caso $\omega\in C\setminus A$ es simétrico.
  \item Así $A\dsim C\subseteq(A\dsim B)\cup(B\dsim C)$ y $d(A,C)\le\PP\bigl((A\dsim B)\cup(B\dsim C)\bigr)\le d(A,B)+d(B,C)$.
\end{enumerate}
\resultado las dos afirmaciones; $d$ se comporta como una distancia entre sucesos.
\comprobacion con dos dados, $A$ «primer dado par», $B$ «segundo dado par»: $d(A,B)=\frac12+\frac12-2\cdot\frac14=\frac12$, y en efecto $A\dsim B$ = «exactamente un dado par», de probabilidad $1/2$.
\end{solucion}

\begin{solucion}{eje:3-hash}
\datos $n$ claves, $N$ casillas, asignación uniforme e independiente.
\planteamiento \cref{prop:t3-cumple}: $\PP(\text{ninguna colisión})\le\exp(-n(n-1)/(2N))$.
\pasos
\begin{enumerate}[label=\arabic*.]
  \item (a) Si $n(n-1)\ge2N\ln2$, entonces $\exp\bigl(-\frac{n(n-1)}{2N}\bigr)\le e^{-\ln2}=\frac12$, luego $\PP(\text{alguna colisión})\ge\frac12$.
  \item (b) $2N\ln2=2\cdot2^{32}\ln2\approx5\,954\,088\,944$. Para $n=77\,163$, $n(n-1)=5\,954\,051\,406$, insuficiente; para $n=77\,164$, $n(n-1)=5\,954\,205\,732$, suficiente.
  \item $\sqrt N=65\,536$ y $77\,164/65\,536\approx1{,}177\approx\sqrt{2\ln2}$.
\end{enumerate}
\resultado $n=77\,164$ claves bastan para una probabilidad de colisión de al menos $1/2$, del orden de $\sqrt N$.
\comprobacion el producto exacto $\prod_{i<n}(1-i/N)$ con $n=77\,164$ vale $0{,}49999$, de modo que la cota es casi exacta en este régimen.
\end{solucion}

\begin{solucion}{eje:3-laplace}
\datos (a) $\Omega$ finito, $\PP(A)=\abs A/\abs\Omega$; (b) $\Omega=\N$.
\planteamiento (a) verificar los tres axiomas; (b) reducción al absurdo con la aditividad numerable.
\pasos
\begin{enumerate}[label=\arabic*.]
  \item (a) $\abs A\ge0$ da la no negatividad; $\PP(\Omega)=\abs\Omega/\abs\Omega=1$.
  \item Si $A_1,A_2,\dots$ son disjuntos dos a dos, como $\Omega$ es finito solo un número finito de ellos es no vacío; por la regla de la suma, $\abs{\bigcup_iA_i}=\sum_i\abs{A_i}$ (los vacíos aportan 0). Dividiendo por $\abs\Omega$ se obtiene la aditividad.
  \item (b) Supóngase $\PP(\{n\})=c$ para todo $n\in\N$. Por la aditividad numerable, $1=\PP(\N)=\sum_{n\in\N}c$.
  \item Si $c=0$, la serie vale $0\neq1$; si $c>0$, diverge. En ambos casos, contradicción.
\end{enumerate}
\resultado el modelo de Laplace es una ley de probabilidad, y no existe una ley «uniforme» sobre $\N$.
\comprobacion sobre $\N$ sí existen leyes que dan probabilidad positiva a todos los puntos, como $\PP(\{n\})=2^{-(n+1)}$, pero no pueden ser constantes.
\end{solucion}

% ---------------------------------------------------------------------
\section{Tema 4}

\begin{solucion}{eje:4-tabla}
\datos 100 estudiantes: 40 de Ingeniería (30 aprueban), 60 de Matemáticas (36 aprueban).
\planteamiento modelo de Laplace y definición de probabilidad condicionada.
\pasos
\begin{enumerate}[label=\arabic*.]
  \item $\PP(\text{ap})=(30+36)/100=0{,}66$.
  \item $\PP(\text{ap}\mid I)=30/40=0{,}75$.
  \item $\PP(I\mid\text{ap})=30/66=5/11\approx0{,}4545$.
  \item $\PP(\text{ap}\mid I)=0{,}75\neq0{,}66=\PP(\text{ap})$: no son independientes.
\end{enumerate}
\resultado $0{,}66$; $0{,}75$; $5/11$; no independientes.
\comprobacion Bayes: $\PP(I\mid\text{ap})=\PP(\text{ap}\mid I)\PP(I)/\PP(\text{ap})=0{,}75\cdot0{,}4/0{,}66=5/11$.
\end{solucion}

\begin{solucion}{eje:4-cachorros}
\datos dos crías, sexos equiprobables e independientes: $\{\mathrm{HH},\mathrm{HM},\mathrm{MH},\mathrm{MM}\}$ equiprobables (la primera letra, la cría mayor).
\planteamiento definición de probabilidad condicionada (\cref{ej:t4-hijos}).
\pasos
\begin{enumerate}[label=(\alph*)]
  \item Condicionante «mayor hembra» $=\{\mathrm{HH},\mathrm{HM}\}$; $\PP(\mathrm{HH}\mid\cdot)=\frac{1/4}{1/2}=\frac12$.
  \item Condicionante «al menos un macho» $=\{\mathrm{HM},\mathrm{MH},\mathrm{MM}\}$; $\PP(\mathrm{MM}\mid\cdot)=\frac{1/4}{3/4}=\frac13$.
\end{enumerate}
\resultado $1/2$ y $1/3$.
\comprobacion en (a) la información se refiere a una cría concreta y la otra sigue siendo equiprobable; en (b) no, y el suceso condicionante tiene tres casos.
\end{solucion}

\begin{solucion}{eje:4-spam}
\datos $\PP(S)=0{,}8$, $\PP(F\mid S)=0{,}1$, $\PP(F\mid S^{c})=0{,}01$.
\planteamiento teorema de Bayes con la partición $\{S,S^{c}\}$.
\pasos
\begin{enumerate}[label=\arabic*.]
  \item $\PP(F)=0{,}1\cdot0{,}8+0{,}01\cdot0{,}2=0{,}08+0{,}002=0{,}082$.
  \item $\PP(S\mid F)=0{,}08/0{,}082=40/41$.
\end{enumerate}
\resultado $40/41\approx0{,}9756$.
\comprobacion forma de razones: razón a priori $0{,}8/0{,}2=4$, razón de verosimilitudes $0{,}1/0{,}01=10$, razón a posteriori $40$, es decir, probabilidad $40/41$.
\end{solucion}

\begin{solucion}{eje:4-dos-pruebas}
\datos $\PP(E)=0{,}01$; $\PP(+\mid E)=0{,}95$; $\PP(+\mid E^{c})=0{,}05$; independencia condicional de las pruebas.
\planteamiento Bayes con la evidencia conjunta $++$, y Bayes dos veces sucesivas.
\pasos
\begin{enumerate}[label=\arabic*.]
  \item (a) $\PP(++\mid E)=0{,}95^{2}=0{,}9025$ y $\PP(++\mid E^{c})=0{,}05^{2}=0{,}0025$.
  \item $\PP(E\mid++)=\frac{0{,}01\cdot0{,}9025}{0{,}01\cdot0{,}9025+0{,}99\cdot0{,}0025}=\frac{0{,}009025}{0{,}0115}\approx0{,}7848$.
  \item (b) Tras la primera prueba: $p_1=\frac{0{,}0095}{0{,}0095+0{,}0495}=\frac{0{,}0095}{0{,}059}\approx0{,}1610$.
  \item Tomando $p_1$ como nueva probabilidad a priori: $\frac{0{,}95\,p_1}{0{,}95\,p_1+0{,}05(1-p_1)}\approx\frac{0{,}15297}{0{,}15297+0{,}04195}\approx0{,}7848$.
\end{enumerate}
\resultado $361/460\approx0{,}7848$ por ambas vías.
\comprobacion razones: a priori $1/99$; cada prueba multiplica por $0{,}95/0{,}05=19$; a posteriori $361/99$, es decir, probabilidad $361/460$.
\end{solucion}

\begin{solucion}{eje:4-piezas}
\datos tipo antiguo $V$ o nuevo $N$ con probabilidad $1/2$; $\PP(\text{def}\mid V)=0{,}15$, $\PP(\text{def}\mid N)=0{,}05$; dos extracciones con reemplazamiento, independientes dado el tipo.
\planteamiento probabilidad total y Bayes.
\pasos
\begin{enumerate}[label=\arabic*.]
  \item $\PP(DD\mid V)=0{,}15^{2}=0{,}0225$ y $\PP(DD\mid N)=0{,}05^{2}=0{,}0025$.
  \item (a) $\PP(DD)=\frac12(0{,}0225+0{,}0025)=0{,}0125$.
  \item (b) $\PP(V\mid DD)=\frac{0{,}5\cdot0{,}0225}{0{,}0125}=\frac{0{,}01125}{0{,}0125}=0{,}9$.
\end{enumerate}
\resultado $0{,}0125$ y $0{,}9$.
\comprobacion razón de verosimilitudes $0{,}0225/0{,}0025=9$ con razón a priori 1: razón a posteriori 9, probabilidad $9/10$. El número de piezas de cada tipo no interviene, porque el tipo se elige con la moneda.
\end{solucion}

\begin{solucion}{eje:4-perro}
\datos $\PP(A)=0{,}4$, $\PP(B)=0{,}6$; $\PP(\text{hallar}\mid A,\text{buscar en }A)=0{,}25$; $\PP(\text{hallar}\mid B,\text{buscar en }B)=0{,}15$.
\planteamiento regla de la multiplicación, Bayes y probabilidad total.
\pasos
\begin{enumerate}[label=\arabic*.]
  \item (a) Buscando en A: $0{,}4\cdot0{,}25=0{,}10$; en B: $0{,}6\cdot0{,}15=0{,}09$. Conviene A.
  \item (b) Sea $F$ «no se encuentra buscando en A»: $\PP(F\mid A)=0{,}75$, $\PP(F\mid B)=1$. $\PP(A\mid F)=\frac{0{,}4\cdot0{,}75}{0{,}4\cdot0{,}75+0{,}6\cdot1}=\frac{0{,}3}{0{,}9}=\frac13$.
  \item (c) Sea $H$ «se halla el primer día». $\PP(H\mid\text{buscar en A})=0{,}1$, $\PP(H\mid\text{buscar en B})=0{,}09$. $\PP(\text{buscó en A}\mid H)=\frac{0{,}5\cdot0{,}1}{0{,}5\cdot0{,}1+0{,}5\cdot0{,}09}=\frac{10}{19}$.
\end{enumerate}
\resultado (a) en A; (b) $1/3$; (c) $10/19\approx0{,}526$.
\comprobacion en (b), la búsqueda fallida en A reduce la probabilidad de A de $0{,}4$ a $0{,}33$, como cabe esperar.
\end{solucion}

\begin{solucion}{eje:4-monty7}
\datos $n$ puertas, el presentador abre $m$ con cabra entre las $n-1$ no elegidas; caso (a): $n=7$, $m=3$.
\planteamiento probabilidad total según la elección inicial sea correcta o no.
\pasos
\begin{enumerate}[label=\arabic*.]
  \item Si la elección inicial es correcta (probabilidad $1/n$), cambiar pierde siempre.
  \item Si es incorrecta (probabilidad $(n-1)/n$), el coche está tras una de las $n-1-m$ puertas cerradas restantes, y cambiando a una de ellas al azar se acierta con probabilidad $1/(n-m-1)$.
  \item $\PP(\text{ganar cambiando})=\frac{n-1}{n}\cdot\frac{1}{n-m-1}$.
  \item (a) $n=7$, $m=3$: $\frac67\cdot\frac13=\frac27$, frente a $\frac17$ sin cambiar.
  \item (b) Como $n-1>n-m-1$, $\frac{n-1}{n(n-m-1)}>\frac1n$: siempre conviene cambiar.
\end{enumerate}
\resultado $2/7$ con 7 puertas; en general $\frac{n-1}{n(n-m-1)}>\frac1n$.
\comprobacion $n=3$, $m=1$: $\frac23\cdot1=\frac23$, el valor del \cref{ej:t4-monty}.
\end{solucion}

\begin{solucion}{eje:4-monty-p}
\datos elección inicial puerta 1; $\PP(C_i)=1/3$; $\PP(M_2\mid C_1)=p$, $\PP(M_3\mid C_1)=1-p$; $\PP(M_3\mid C_2)=1$; $\PP(M_2\mid C_3)=1$.
\planteamiento Bayes con la partición $\{C_1,C_2,C_3\}$.
\pasos
\begin{enumerate}[label=\arabic*.]
  \item (a) Cambiar gana si y solo si el coche no está en la puerta~1: $2/3$.
  \item (b) $\PP(M_2)=\frac13p+0+\frac13=\frac{1+p}3$; cambiar gana si $C_3$: $\PP(C_3\mid M_2)=\frac{1/3}{(1+p)/3}=\frac1{1+p}$.
  \item (c) $\PP(M_3)=\frac13(1-p)+\frac13=\frac{2-p}3$; $\PP(C_2\mid M_3)=\frac1{2-p}$.
\end{enumerate}
\resultado $2/3$; $1/(1+p)$; $1/(2-p)$.
\comprobacion probabilidad total: $\frac{1+p}3\cdot\frac1{1+p}+\frac{2-p}3\cdot\frac1{2-p}=\frac23$. Con $p=1/2$, (b) y (c) valen $2/3$; con $p=1$, (b) vale $1/2$ y (c) vale~1.
\end{solucion}

\begin{solucion}{eje:4-paridad}
\datos dos dados; $A$, $B$, $C$ como en el enunciado.
\planteamiento la suma es impar si y solo si exactamente uno de los dos es impar.
\pasos
\begin{enumerate}[label=\arabic*.]
  \item $\PP(A)=\PP(B)=\frac12$; $\PP(C)=\PP(\text{impar, par})+\PP(\text{par, impar})=\frac14+\frac14=\frac12$.
  \item $\PP(A\cap B)=\frac14$; $A\cap C$ = «primero impar, segundo par», $\frac14$; $B\cap C$ = «primero par, segundo impar», $\frac14$. Independientes dos a dos.
  \item $A\cap B\cap C=\emptyset$: si ambos son impares, la suma es par. $\PP(A\cap B\cap C)=0\neq\frac18$.
\end{enumerate}
\resultado independientes dos a dos, pero no mutuamente.
\comprobacion es el mismo fenómeno que el \cref{ej:t4-dosados}: $C$ es una función de las paridades de los dos dados.
\end{solucion}

\begin{solucion}{eje:4-ajedrez}
\datos niveles $P$, $I$, $M$ con probabilidad $1/3$; $\PP(G\mid P)=0{,}9$, $\PP(G\mid I)=0{,}5$, $\PP(G\mid M)=0{,}3$; partidas independientes dado el nivel.
\planteamiento probabilidad total, condicionando al nivel.
\pasos
\begin{enumerate}[label=\arabic*.]
  \item (a) $\PP(G_1)=\frac13(0{,}9+0{,}5+0{,}3)=\frac{1{,}7}3\approx0{,}5667$.
  \item $\PP(G_1\cap G_2)=\frac13(0{,}9^{2}+0{,}5^{2}+0{,}3^{2})=\frac{1{,}15}3\approx0{,}3833$.
  \item (b) $\PP(G_2\mid G_1)=\frac{1{,}15}{1{,}7}=\frac{23}{34}\approx0{,}6765$.
  \item (c) $\PP(G_2\mid G_1)\neq\PP(G_2)=0{,}5667$: no son independientes. Ganar la primera hace más probable que el rival sea débil y, por tanto, que se gane la segunda. Son independientes \emph{condicionalmente} al nivel, pero no incondicionalmente.
\end{enumerate}
\resultado $17/30$; $23/34$; dependientes, aunque condicionalmente independientes dado el nivel.
\comprobacion $\PP(G_1\cap G_2)=0{,}3833>\PP(G_1)^{2}=0{,}3211$: los resultados están positivamente relacionados.
\end{solucion}

\begin{solucion}{eje:4-si-mismo}
\datos un suceso $A$.
\planteamiento la independencia de $A$ consigo mismo es $\PP(A\cap A)=\PP(A)^{2}$.
\pasos
\begin{enumerate}[label=\arabic*.]
  \item $\PP(A)=\PP(A)^{2}\iff\PP(A)(1-\PP(A))=0\iff\PP(A)\in\{0,1\}$.
  \item Si $\PP(A)=0$: $0\le\PP(A\cap B)\le\PP(A)=0=\PP(A)\PP(B)$.
  \item Si $\PP(A)=1$: $\PP(A^{c})=0$, luego $\PP(A^{c}\cap B)=0$ y $\PP(A\cap B)=\PP(B)-\PP(A^{c}\cap B)=\PP(B)=\PP(A)\PP(B)$.
\end{enumerate}
\resultado $A$ es independiente de sí mismo si y solo si $\PP(A)\in\{0,1\}$, y entonces es independiente de cualquier suceso.
\comprobacion $\Omega$ y $\emptyset$ son independientes de todo suceso, como cabe esperar: saber que ha ocurrido $\Omega$ no informa de nada.
\end{solucion}

\begin{solucion}{eje:4-union-indep}
\datos (a) $A$, $B$, $C$ mutuamente independientes; (b) $A_1,\dots,A_n$ independientes.
\planteamiento (a) distributiva e inclusión--exclusión; (b) independencia de los complementarios y De Morgan.
\pasos
\begin{enumerate}[label=\arabic*.]
  \item (a) $\PP(A\cap(B\cup C))=\PP((A\cap B)\cup(A\cap C))=\PP(A\cap B)+\PP(A\cap C)-\PP(A\cap B\cap C)$.
  \item Por la independencia mutua, esto es $\PP(A)\bigl[\PP(B)+\PP(C)-\PP(B)\PP(C)\bigr]=\PP(A)\PP(B\cup C)$, usando $\PP(B\cap C)=\PP(B)\PP(C)$.
  \item (b) Si $A_1,\dots,A_n$ son independientes, sustituir uno de ellos, $A_i$, por $A_i^{c}$ conserva la independencia: para todo $S\ni i$, con $T=S\setminus\{i\}$, $\PP\bigl(\bigcap_{T}A_j\cap A_i^{c}\bigr)=\PP\bigl(\bigcap_TA_j\bigr)-\PP\bigl(\bigcap_SA_j\bigr)=\prod_T\PP(A_j)\,(1-\PP(A_i))$. Repitiendo, $A_1^{c},\dots,A_n^{c}$ son independientes.
  \item $\PP\bigl(\bigcup A_i\bigr)=1-\PP\bigl(\bigcap A_i^{c}\bigr)=1-\prod_i(1-\PP(A_i))$.
\end{enumerate}
\resultado las dos afirmaciones.
\comprobacion $n$ lanzamientos de una moneda, $A_i$ «cara en el $i$»: $\PP(\text{alguna cara})=1-2^{-n}$, el resultado conocido.
\end{solucion}

\begin{solucion}{eje:4-producto-triple}
\planteamiento basta que la intersección triple tenga la probabilidad «correcta» mientras falla alguna de las condiciones dos a dos.
\pasos
\begin{enumerate}[label=\arabic*.]
  \item Ejemplo 1: $\Omega=[8]$ con ley uniforme, $A=\{1,2,3,4\}$, $B=\{1,2,3,5\}$, $C=\{1,6,7,8\}$. Cada uno tiene probabilidad $1/2$.
  \item $A\cap B\cap C=\{1\}$, con probabilidad $1/8=\PP(A)\PP(B)\PP(C)$.
  \item Pero $A\cap B=\{1,2,3\}$ tiene probabilidad $3/8\neq1/4$: no son mutuamente independientes (ni siquiera dos a dos).
  \item Ejemplo 2, degenerado: $C=\emptyset$ y $A=B$ con $\PP(A)=1/2$. La triple vale $0=\PP(A)\PP(B)\cdot0$, pero $\PP(A\cap B)=1/2\neq1/4$.
\end{enumerate}
\resultado la condición sobre la intersección triple no implica la independencia mutua.
\comprobacion en el ejemplo 1, $\PP(A\cap C)=\PP(\{1\})=1/8\neq1/4$ y $\PP(B\cap C)=1/8\neq1/4$: fallan las tres condiciones dobles.
\end{solucion}

\begin{solucion}{eje:4-razones}
\datos sucesos $B$, $E_1$, $E_2$ con probabilidades positivas; en (b), independencia condicional dada $B$ y dada $B^{c}$.
\planteamiento (a) cociente de dos fórmulas de Bayes; (b) comparar el producto de las dos razones de verosimilitudes con la conjunta.
\pasos
\begin{enumerate}[label=\arabic*.]
  \item (a) $\PP(B\mid A)=\PP(A\mid B)\PP(B)/\PP(A)$ y $\PP(B^{c}\mid A)=\PP(A\mid B^{c})\PP(B^{c})/\PP(A)$; al dividir se cancela $\PP(A)$.
  \item (b) Una sola actualización: $O_{12}=O_0\cdot\frac{\PP(E_1\cap E_2\mid B)}{\PP(E_1\cap E_2\mid B^{c})}=O_0\cdot\frac{\PP(E_1\mid B)\PP(E_2\mid B)}{\PP(E_1\mid B^{c})\PP(E_2\mid B^{c})}$, por la independencia condicional.
  \item Dos actualizaciones: tras $E_1$, $O_1=O_0\frac{\PP(E_1\mid B)}{\PP(E_1\mid B^{c})}$. Al incorporar $E_2$, las verosimilitudes son ahora $\PP(E_2\mid B\cap E_1)$ y $\PP(E_2\mid B^{c}\cap E_1)$, que por la independencia condicional valen $\PP(E_2\mid B)$ y $\PP(E_2\mid B^{c})$.
  \item Por tanto $O_2=O_1\frac{\PP(E_2\mid B)}{\PP(E_2\mid B^{c})}=O_{12}$.
\end{enumerate}
\resultado ambas formas de actualizar coinciden.
\comprobacion el \cref{eje:4-dos-pruebas} es un caso numérico: $\frac1{99}\cdot19\cdot19=\frac{361}{99}$ por ambas vías.
\end{solucion}

\begin{solucion}{eje:4-simetria}
\datos $0<\PP(A),\PP(B)<1$.
\planteamiento reescribir cada afirmación en términos de $\PP(A\cap B)$ y $\PP(A)\PP(B)$.
\pasos
\begin{enumerate}[label=\arabic*.]
  \item (a) $\iff\PP(A\cap B)/\PP(B)>\PP(A)\iff\PP(A\cap B)>\PP(A)\PP(B)$.
  \item (b) $\iff\PP(A\cap B)/\PP(A)>\PP(B)\iff\PP(A\cap B)>\PP(A)\PP(B)$. Luego (a)$\iff$(b).
  \item Por la probabilidad total, $\PP(A)=\PP(A\mid B)\PP(B)+\PP(A\mid B^{c})(1-\PP(B))$, de donde
  \[
  \PP(A\mid B)-\PP(A)=(1-\PP(B))\bigl(\PP(A\mid B)-\PP(A\mid B^{c})\bigr).
  \]
  \item Como $1-\PP(B)>0$, (a)$\iff$(c).
\end{enumerate}
\resultado las tres afirmaciones son equivalentes: la «correlación positiva» entre sucesos es simétrica.
\comprobacion en el \cref{eje:4-tabla}: $\PP(\text{ap}\mid I)=0{,}75>0{,}66$ y $\PP(I\mid\text{ap})=5/11\approx0{,}45>0{,}4=\PP(I)$; además $\PP(\text{ap}\mid I^{c})=0{,}6<0{,}75$.
\end{solucion}

% ---------------------------------------------------------------------
\section{Tema 5}

\begin{solucion}{eje:5-masa}
\datos $\PP(X=k)=ck$, $k=1,\dots,4$.
\planteamiento normalización, acumulación y definiciones de esperanza y varianza (\cref{prop:t5-var}).
\pasos
\begin{enumerate}[label=\arabic*.]
  \item $c(1+2+3+4)=10c=1$, luego $c=1/10$.
  \item $F_X(1)=0{,}1$, $F_X(2)=0{,}3$, $F_X(3)=0{,}6$, $F_X(4)=1$.
  \item $\E[X]=\sum k\cdot\frac{k}{10}=\frac{1+4+9+16}{10}=3$.
  \item $\E[X^{2}]=\sum k^{2}\cdot\frac k{10}=\frac{1+8+27+64}{10}=10$; $\Var(X)=10-9=1$.
\end{enumerate}
\resultado $c=1/10$; $F_X=0{,}1;\,0{,}3;\,0{,}6;\,1$; $\E[X]=3$; $\Var(X)=1$.
\comprobacion $\Var(X)=\sum(k-3)^{2}\frac k{10}=\frac{4+2+0+4}{10}=1$.
\end{solucion}

\begin{solucion}{eje:5-diferencia}
\datos dos dados de tres caras; $X=D_1-D_2$; 9 resultados equiprobables.
\planteamiento contar los pares que dan cada diferencia; LOTUS para $X^{2}$.
\pasos
\begin{enumerate}[label=\arabic*.]
  \item Valores: $-2$ (1 par), $-1$ (2), $0$ (3), $1$ (2), $2$ (1); $p_X=\frac19,\frac29,\frac39,\frac29,\frac19$.
  \item $\E[X]=0$ por simetría ($X$ y $-X$ tienen la misma distribución).
  \item $\Var(X)=\E[X^{2}]=\frac{4\cdot1+1\cdot2+0+1\cdot2+4\cdot1}{9}=\frac{12}9=\frac43$.
  \item $X^{2}$ toma $0$ con probabilidad $\frac39$, $1$ con $\frac49$ y $4$ con $\frac29$; $\E[X^{2}]=\frac49+\frac89=\frac43$.
\end{enumerate}
\resultado $p_X$ como arriba; $\E[X]=0$; $\Var(X)=4/3$; $p_{X^{2}}=(\frac13,\frac49,\frac29)$ en $(0,1,4)$.
\comprobacion por independencia, $\Var(D_1-D_2)=\Var(D_1)+\Var(D_2)=2\cdot\frac{3^{2}-1}{12}=\frac43$ (\cref{tab:t5-discretas}).
\end{solucion}

\begin{solucion}{eje:5-binomial}
\datos $X\sim\Bin(10;0{,}3)$.
\planteamiento \cref{teo:t5-binomial}.
\pasos
\begin{enumerate}[label=\arabic*.]
  \item $\PP(X=3)=\binom{10}3\,0{,}3^{3}\,0{,}7^{7}=120\cdot0{,}027\cdot0{,}0823543\approx0{,}2668$.
  \item $\PP(X\ge1)=1-0{,}7^{10}\approx1-0{,}0282=0{,}9718$.
  \item $\E[X]=10\cdot0{,}3=3$; $\Var(X)=10\cdot0{,}3\cdot0{,}7=2{,}1$.
\end{enumerate}
\resultado $0{,}2668$; $0{,}9718$; $3$; $2{,}1$.
\comprobacion la moda de la binomial es $\lfloor(n+1)p\rfloor=3$, coherente con que $\PP(X=3)$ sea el valor más alto de la función de masa.
\end{solucion}

\begin{solucion}{eje:5-geometrica}
\datos $X\sim\Geom(0{,}2)$, $q=0{,}8$.
\planteamiento \cref{prop:t5-geom}.
\pasos
\begin{enumerate}[label=\arabic*.]
  \item $\PP(X=3)=q^{2}p=0{,}64\cdot0{,}2=0{,}128$.
  \item $\PP(X>5)=q^{5}=0{,}32768$.
  \item $\E[X]=1/0{,}2=5$; $\Var(X)=0{,}8/0{,}04=20$.
  \item Falta de memoria: $\PP(X>8\mid X>3)=\PP(X>5)=0{,}32768$.
\end{enumerate}
\resultado $0{,}128$; $0{,}32768$; $5$; $20$; $0{,}32768$.
\comprobacion directamente, $\PP(X>8)/\PP(X>3)=0{,}8^{8}/0{,}8^{3}=0{,}8^{5}$.
\end{solucion}

\begin{solucion}{eje:5-poisson-aprox}
\datos $X\sim\Bin(1000;\,0{,}002)$; aproximación $\Poi(2)$.
\planteamiento \cref{teo:t5-poisson} con $\lambda=np=2$.
\pasos
\begin{enumerate}[label=\arabic*.]
  \item Binomial: $\PP(X=0)=0{,}998^{1000}\approx0{,}13506$; $\PP(X=1)=1000\cdot0{,}002\cdot0{,}998^{999}\approx0{,}27067$; $\PP(X\ge2)\approx1-0{,}13506-0{,}27067=0{,}59427$.
  \item Poisson: $\PP(X=0)=e^{-2}\approx0{,}13534$; $\PP(X=1)=2e^{-2}\approx0{,}27067$; $\PP(X\ge2)=1-3e^{-2}\approx0{,}59399$.
\end{enumerate}
\resultado la aproximación de Poisson acierta en las tres primeras cifras decimales.
\comprobacion las dos distribuciones tienen media 2 y varianzas $1{,}996$ y $2$, casi iguales.
\end{solucion}

\begin{solucion}{eje:5-cumple-esp}
\datos 50 personas, 365 días.
\planteamiento linealidad con indicadoras (\cref{prop:t5-puente}): una por pareja en (a) y una por día en (b).
\pasos
\begin{enumerate}[label=\arabic*.]
  \item (a) Para cada una de las $\binom{50}2=1225$ parejas, $\PP(\text{mismo día})=1/365$. $\E=1225/365\approx3{,}356$.
  \item (b) Para el día $d$, sea $N_d\sim\Bin(50,1/365)$ el número de personas nacidas ese día. $\PP(N_d\ge2)=1-\bigl(\tfrac{364}{365}\bigr)^{50}-50\cdot\tfrac1{365}\bigl(\tfrac{364}{365}\bigr)^{49}$.
  \item $\E=365\,\PP(N_d\ge2)\approx3{,}076$.
\end{enumerate}
\resultado $\approx3{,}356$ parejas y $\approx3{,}076$ días.
\comprobacion (b) es algo menor que (a) porque un día con tres personas aporta 3 parejas pero un solo día.
\end{solucion}

\begin{solucion}{eje:5-hash}
\datos $k$ claves, $n$ casillas.
\planteamiento indicadoras por casilla; el número de claves en una casilla es $\Bin(k,1/n)$.
\pasos
\begin{enumerate}[label=\arabic*.]
  \item Vacías: $\PP(\text{casilla vacía})=(1-1/n)^{k}$; esperanza $n(1-1/n)^{k}$.
  \item Con exactamente una: $\PP=k\cdot\frac1n(1-\frac1n)^{k-1}$; esperanza $k(1-1/n)^{k-1}$.
  \item Con más de una: $n-n(1-\frac1n)^{k}-k(1-\frac1n)^{k-1}$, por linealidad, pues las tres indicadoras de cada casilla suman~1.
  \item $k=n=100$: $100\cdot0{,}99^{100}\approx36{,}60$; $100\cdot0{,}99^{99}\approx36{,}97$; $100-36{,}60-36{,}97\approx26{,}42$.
\end{enumerate}
\resultado $36{,}60$ vacías, $36{,}97$ con una clave y $26{,}42$ con más de una.
\comprobacion la fracción de casillas vacías, $0{,}366$, está próxima a $e^{-1}\approx0{,}368$, el límite de $(1-1/n)^{n}$.
\end{solucion}

\begin{solucion}{eje:5-hijos}
\datos sexos equiprobables e independientes.
\planteamiento el primer hijo fija un sexo; a partir de ahí se espera el primer hijo del sexo opuesto, lo que es una geométrica.
\pasos
\begin{enumerate}[label=\arabic*.]
  \item Sea $Y$ el número de hijos tras el primero hasta el primero del sexo contrario, incluido. $Y\sim\Geom(1/2)$, independientemente del sexo del primero.
  \item $X=1+Y$, $\E[X]=1+2=3$.
\end{enumerate}
\resultado 3 hijos en promedio.
\comprobacion directamente: $\PP(X=k)=2\cdot(1/2)^{k}$ para $k\ge2$ (los $k-1$ primeros del mismo sexo, cualquiera de los dos, y el $k$-ésimo distinto), y $\sum_{k\ge2}k\,2^{1-k}=3$.
\end{solucion}

\begin{solucion}{eje:5-inversiones}
\datos permutación uniforme de $n$ números distintos.
\planteamiento una indicadora por cada par de posiciones $i<j$.
\pasos
\begin{enumerate}[label=\arabic*.]
  \item Para $i<j$, intercambiar los elementos de las posiciones $i$ y $j$ es una biyección entre las permutaciones con el par invertido y las que no lo tienen; luego $\PP(\text{invertido})=1/2$.
  \item Por linealidad, $\E[\text{inversiones}]=\binom n2\cdot\frac12=\frac{n(n-1)}4$.
\end{enumerate}
\resultado $n(n-1)/4$.
\comprobacion para $n=4$, las 24 permutaciones suman 72 inversiones, de media $3=4\cdot3/4$.
\end{solucion}

\begin{solucion}{eje:5-categorica}
\datos 10 tokens i.i.d. con $p=(0{,}5;\,0{,}3;\,0{,}2)$.
\planteamiento el número de apariciones de cada token es binomial; para (b), una indicadora por token.
\pasos
\begin{enumerate}[label=\arabic*.]
  \item $N_{\texttt a}\sim\Bin(10;0{,}5)$, $N_{\texttt b}\sim\Bin(10;0{,}3)$, $N_{\texttt c}\sim\Bin(10;0{,}2)$: medias $5$, $3$ y $2$.
  \item $\Var(N_{\texttt a})=10\cdot0{,}5\cdot0{,}5=2{,}5$.
  \item (b) $\PP(\text{el token }i\text{ aparece})=1-(1-p_i)^{10}$. Distintos esperados: $3-(0{,}5^{10}+0{,}7^{10}+0{,}8^{10})\approx3-(0{,}00098+0{,}02825+0{,}10737)=2{,}8634$.
\end{enumerate}
\resultado medias $5$, $3$, $2$; varianza $2{,}5$; $\approx2{,}863$ tokens distintos.
\comprobacion las medias suman 10, el número de tokens generados, como exige la linealidad.
\end{solucion}

\begin{solucion}{eje:5-passk}
\datos $n=10$ soluciones, $c=3$ correctas, se eligen $k$ sin reemplazamiento.
\planteamiento modelo de Laplace sobre los $\binom nk$ subconjuntos; complemento.
\pasos
\begin{enumerate}[label=\arabic*.]
  \item (a) Ninguna correcta equivale a elegir las $k$ entre las $n-c$ incorrectas: $\binom{n-c}k$ subconjuntos. $\PP(\text{alguna correcta})=1-\binom{n-c}k\big/\binom nk$.
  \item (b) $k=1$: $1-7/10=0{,}3$. $k=2$: $1-21/45=8/15\approx0{,}533$. $k=5$: $1-21/252=11/12\approx0{,}917$.
\end{enumerate}
\resultado $0{,}3$; $0{,}533$; $0{,}917$.
\comprobacion con $k=1$ se recupera la proporción de correctas, $c/n$; y si $k>n-c=7$ la probabilidad es 1, porque $\binom7k=0$.
\end{solucion}

\begin{solucion}{eje:5-colas}
\datos $X$ con valores en $\{0,1,2,\dots\}$.
\planteamiento intercambiar el orden en una suma doble de términos no negativos.
\pasos
\begin{enumerate}[label=\arabic*.]
  \item $\sum_{k\ge1}\PP(X\ge k)=\sum_{k\ge1}\sum_{j\ge k}p_X(j)$.
  \item Los pares $(k,j)$ con $1\le k\le j$ pueden sumarse primero en $k$: para cada $j\ge1$ hay $j$ valores de $k$. Como todos los términos son no negativos, el reordenamiento es lícito: $\sum_{j\ge1}j\,p_X(j)=\E[X]$.
  \item Equivalentemente, $X=\sum_{k\ge1}I_{\{X\ge k\}}$ (si $X=j$, valen 1 exactamente las indicadoras con $k\le j$).
  \item Geométrica: $\PP(X\ge k)=q^{k-1}$, luego $\E[X]=\sum_{k\ge1}q^{k-1}=\frac1{1-q}=\frac1p$.
\end{enumerate}
\resultado $\E[X]=\sum_{k\ge1}\PP(X\ge k)$ y, para la geométrica, $1/p$.
\comprobacion para un dado, $\sum_{k=1}^{6}\frac{7-k}{6}=\frac{21}6=3{,}5$.
\end{solucion}

\begin{solucion}{eje:5-poisson}
\datos $p(k)=e^{-\lambda}\lambda^{k}/k!$.
\planteamiento serie de la exponencial, $\sum_{k\ge0}\lambda^{k}/k!=e^{\lambda}$, y desplazamiento de índices.
\pasos
\begin{enumerate}[label=\arabic*.]
  \item $\sum_kp(k)=e^{-\lambda}e^{\lambda}=1$.
  \item $\E[X]=\sum_{k\ge1}k\,e^{-\lambda}\frac{\lambda^{k}}{k!}=\lambda\sum_{k\ge1}e^{-\lambda}\frac{\lambda^{k-1}}{(k-1)!}=\lambda$.
  \item $\E[X(X-1)]=\sum_{k\ge2}k(k-1)e^{-\lambda}\frac{\lambda^{k}}{k!}=\lambda^{2}\sum_{k\ge2}e^{-\lambda}\frac{\lambda^{k-2}}{(k-2)!}=\lambda^{2}$.
  \item $\Var(X)=\E[X(X-1)]+\E[X]-\E[X]^{2}=\lambda^{2}+\lambda-\lambda^{2}=\lambda$.
\end{enumerate}
\resultado la función de masa suma 1 y $\E[X]=\Var(X)=\lambda$.
\comprobacion la binomial $\Bin(n,\lambda/n)$ tiene media $\lambda$ y varianza $\lambda(1-\lambda/n)\to\lambda$, coherente con el \cref{teo:t5-poisson}.
\end{solucion}

\begin{solucion}{eje:5-sin-memoria}
\datos $X\in\{1,2,\dots\}$ sin memoria; $q=\PP(X>1)\in(0,1)$; $g(k)=\PP(X>k)$.
\planteamiento la falta de memoria se escribe $g(m+k)=g(m)g(k)$; se resuelve por inducción.
\pasos
\begin{enumerate}[label=\arabic*.]
  \item $g(0)=1$ porque $X\ge1$. Si $g(m)>0$, $\PP(X>m+k\mid X>m)=g(m+k)/g(m)$, luego $g(m+k)=g(m)g(k)$.
  \item Inducción: $g(k)=q^{k}$. Cierto para $k=0,1$. Si $g(k)=q^{k}>0$, la relación con $m=k$ y $1$ da $g(k+1)=g(k)g(1)=q^{k+1}$.
  \item $\PP(X=k)=g(k-1)-g(k)=q^{k-1}(1-q)$ para $k\ge1$.
\end{enumerate}
\resultado $X\sim\Geom(p)$ con $p=1-q$.
\comprobacion la recíproca es la \cref{prop:t5-geom}(iii): la geométrica cumple la propiedad.
\end{solucion}

\begin{solucion}{eje:5-geom-cond}
\datos $X$, $Y$ independientes $\Geom(p)$, $q=1-p$, $n\ge2$.
\planteamiento probabilidad condicionada e independencia.
\pasos
\begin{enumerate}[label=\arabic*.]
  \item Para $1\le i\le n-1$: $\PP(X=i,\,X+Y=n)=\PP(X=i)\PP(Y=n-i)=q^{i-1}p\cdot q^{n-i-1}p=p^{2}q^{n-2}$.
  \item No depende de $i$. Sumando en $i$: $\PP(X+Y=n)=(n-1)p^{2}q^{n-2}$.
  \item $\PP(X=i\mid X+Y=n)=\frac{p^{2}q^{n-2}}{(n-1)p^{2}q^{n-2}}=\frac1{n-1}$.
\end{enumerate}
\resultado uniforme en $\{1,\dots,n-1\}$.
\comprobacion interpretación: dado que el segundo éxito de una sucesión de ensayos ocurre en el ensayo $n$, el primero está en cualquiera de las $n-1$ posiciones anteriores con igual probabilidad, porque todas las sucesiones con dos éxitos en $n$ ensayos y el último éxito en $n$ tienen la misma probabilidad $p^{2}q^{n-2}$.
\end{solucion}

% ---------------------------------------------------------------------
\section{Tema 6}

\begin{solucion}{eje:6-conceptos}
\pasos
\begin{enumerate}[label=(\alph*)]
  \item Falsa: la densidad de $\Unif(0,\tfrac12)$ vale 2. Lo que debe ser $\le1$ es la integral sobre cualquier intervalo.
  \item Verdadera: $\PP(X\le2)=\PP(X<2)+\PP(X=2)$ y $\PP(X=2)=0$.
  \item Falsa: si $X$ es continua, $\{X=x\}$ tiene probabilidad 0 para todo $x$, y sin embargo $X$ toma algún valor.
  \item Verdadera: si $x\le y$, $\{X\le x\}\subseteq\{X\le y\}$ y se aplica la monotonía.
\end{enumerate}
\resultado (a) F; (b) V; (c) F; (d) V.
\comprobacion en (c), si todos los sucesos de probabilidad nula fueran imposibles, $\Omega=\bigcup_x\{X=x\}$ sería imposible; no hay contradicción con la aditividad porque la unión no es numerable.
\end{solucion}

\begin{solucion}{eje:6-densidad}
\datos $f(x)=cx^{2}$ en $[0,2]$.
\planteamiento normalización e integrales de potencias.
\pasos
\begin{enumerate}[label=\arabic*.]
  \item $\int_0^2cx^{2}dx=\frac{8c}3=1$, luego $c=3/8$.
  \item $F(x)=0$ si $x<0$; $F(x)=x^{3}/8$ si $0\le x\le2$; $F(x)=1$ si $x>2$.
  \item $\PP(X>1)=1-F(1)=7/8$.
  \item $\E[X]=\frac38\int_0^2x^{3}dx=\frac38\cdot4=\frac32$; $\E[X^{2}]=\frac38\int_0^2x^{4}dx=\frac38\cdot\frac{32}5=\frac{12}5$; $\Var(X)=\frac{12}5-\frac94=\frac3{20}$.
\end{enumerate}
\resultado $c=3/8$; $F(x)=x^{3}/8$ en $[0,2]$; $7/8$; $3/2$; $3/20$.
\comprobacion $F(2)=1$, y la media está a la derecha del punto medio, como corresponde a una densidad creciente.
\end{solucion}

\begin{solucion}{eje:6-pasas}
\datos densidad de pasas $40-h$ en $[0,30]$.
\planteamiento total por integración; la densidad de $H$ es la densidad de pasas normalizada.
\pasos
\begin{enumerate}[label=\arabic*.]
  \item (a) $\int_0^{30}(40-h)\,dh=1200-450=750$ pasas.
  \item (b) $g(h)=\frac{40-h}{750}$ en $[0,30]$; $G(h)=\frac{40h-h^{2}/2}{750}$ en $[0,30]$, 0 antes y 1 después.
  \item (c) $G(10)=\frac{400-50}{750}=\frac7{15}\approx0{,}467$.
\end{enumerate}
\resultado 750 pasas; $g$ y $G$ como arriba; $7/15$.
\comprobacion $G(30)=\frac{1200-450}{750}=1$; y $7/15>1/3$ porque las pasas se acumulan en el fondo.
\end{solucion}

\begin{solucion}{eje:6-exponencial}
\datos $f(y)=\lambda e^{-\lambda y}$, $y\ge0$.
\planteamiento integrar la densidad.
\pasos
\begin{enumerate}[label=\arabic*.]
  \item $F(y)=\int_0^y\lambda e^{-\lambda t}dt=\bigl[-e^{-\lambda t}\bigr]_0^y=1-e^{-\lambda y}$ para $y\ge0$.
  \item $\PP(Y\le1/\lambda)=1-e^{-1}\approx0{,}632$.
  \item Mediana: $1-e^{-\lambda m}=\frac12\iff m=\frac{\ln2}\lambda$.
\end{enumerate}
\resultado $F(y)=1-e^{-\lambda y}$; $1-e^{-1}$; $\ln2/\lambda$.
\comprobacion la mediana, $0{,}693/\lambda$, es menor que la media, $1/\lambda$, como corresponde a una distribución con cola larga a la derecha.
\end{solucion}

\begin{solucion}{eje:6-autobus}
\datos A: $W\sim\Unif(0,10)$; B: $W\sim\Expo(1/10)$.
\planteamiento \cref{prop:t6-unif} y \cref{prop:t6-memoria}; probabilidad condicionada.
\pasos
\begin{enumerate}[label=\arabic*.]
  \item (a) A: $\E[W]=5$; B: $\E[W]=10$.
  \item (b) A: $\PP(W\ge9\mid W\ge6)=\frac{1/10}{4/10}=\frac14$.
  \item B: $\PP(W\ge9\mid W\ge6)=\PP(W\ge3)=e^{-3/10}\approx0{,}741$.
\end{enumerate}
\resultado medias 5 y 10 minutos; $1/4$ en A y $0{,}741$ en B.
\comprobacion en A, tras 6 minutos la espera restante es uniforme en $[0,4]$ y $\PP(\ge3)=1/4$; en B, la espera restante vuelve a ser exponencial de media 10, por la falta de memoria.
\end{solucion}

\begin{solucion}{eje:6-tipificar}
\datos $X\sim\Nor(2,9)$, $\mu=2$, $\sigma=3$.
\planteamiento \cref{teo:t6-tipificacion}: $Z=(X-2)/3\sim\Nor(0,1)$.
\pasos
\begin{enumerate}[label=\arabic*.]
  \item $\PP(X\le2)=\Phi(0)=0{,}5$.
  \item $\PP(X-\mu>1{,}5\sigma)=\PP(Z>1{,}5)=1-0{,}9332=0{,}0668$.
  \item $\PP(\abs{X-\mu}<1{,}5\sigma)=2\Phi(1{,}5)-1=0{,}8664$.
\end{enumerate}
\resultado $0{,}5$; $0{,}0668$; $0{,}8664$.
\comprobacion los resultados no dependen de $\mu$ ni de $\sigma$: solo de cuántas desviaciones típicas se alejan los umbrales de la media.
\end{solucion}

\begin{solucion}{eje:6-gestacion}
\datos $X\sim\Nor(280;\,8{,}5^{2})$.
\planteamiento tipificar y usar $\Phi$ (tabla o software) y su inversa.
\pasos
\begin{enumerate}[label=\arabic*.]
  \item (a) $z=(273-280)/8{,}5\approx-0{,}824$; $\PP(X\le273)=\Phi(-0{,}824)=1-\Phi(0{,}824)\approx0{,}205$.
  \item (b) Se busca $d$ con $\PP(X\le d)=0{,}05$, es decir, $\frac{d-280}{8{,}5}=-1{,}645$ (\cref{tab:t6-phi}); $d=280-1{,}645\cdot8{,}5\approx266{,}0$.
\end{enumerate}
\resultado (a) $\approx0{,}205$; (b) $d\approx266$, unos 14 días antes de la fecha prevista.
\comprobacion (a) está entre $\Phi(-1)=0{,}159$ y $\Phi(-0{,}5)=0{,}309$, como corresponde a $z=-0{,}82$.
\end{solucion}

\begin{solucion}{eje:6-palo}
\datos $U\sim\Unif(0,1)$; $L=\max(U,1-U)\in[\tfrac12,1]$.
\planteamiento calcular $F_L(x)=\PP(L\le x)$ para $x\in[\tfrac12,1]$.
\pasos
\begin{enumerate}[label=\arabic*.]
  \item $L\le x\iff U\le x$ y $1-U\le x\iff1-x\le U\le x$.
  \item $F_L(x)=x-(1-x)=2x-1$ para $\frac12\le x\le1$ (0 antes y 1 después).
  \item Es la función de distribución de $\Unif(\tfrac12,1)$; $\E[L]=\frac34$.
\end{enumerate}
\resultado $L\sim\Unif(\tfrac12,1)$ y $\E[L]=3/4$.
\comprobacion el trozo corto, $1-L$, es $\Unif(0,\tfrac12)$ con media $1/4$, y $\frac34+\frac14=1$.
\end{solucion}

\begin{solucion}{eje:6-markov}
\datos $X\sim\Expo(1)$: $\E[X]=1$, $\Var(X)=1$.
\planteamiento \cref{teo:t6-markov,teo:t6-chebyshev}.
\pasos
\begin{enumerate}[label=\arabic*.]
  \item (a) Markov: $\PP(X\ge3)\le1/3$. Exacto: $e^{-3}\approx0{,}0498$.
  \item (b) Chebyshev: $\PP(\abs{X-1}\ge2)\le1/4$. Como $X\ge0$, $\abs{X-1}\ge2\iff X\ge3$; exacto $e^{-3}\approx0{,}0498$.
\end{enumerate}
\resultado cotas $1/3$ y $1/4$ frente al valor exacto $0{,}0498$.
\comprobacion ambas cotas son válidas pero poco ajustadas: solo usan la media o la varianza, no la forma exponencial de la cola.
\end{solucion}

\begin{solucion}{eje:6-n-cheb}
\datos media $\mu$, varianza $\sigma^{2}$, $\Var(\Xbar_n)=\sigma^{2}/n$.
\planteamiento Chebyshev para $\Xbar_n$ con $c=2\sigma$.
\pasos
\begin{enumerate}[label=\arabic*.]
  \item $\PP(\abs{\Xbar_n-\mu}\ge2\sigma)\le\frac{\sigma^{2}/n}{4\sigma^{2}}=\frac1{4n}$.
  \item $\frac1{4n}\le0{,}01\iff n\ge25$.
\end{enumerate}
\resultado $n=25$ basta, sea cual sea la distribución.
\comprobacion con $n=25$, $\sigma_{\Xbar}=\sigma/5$ y el umbral $2\sigma$ son 10 desviaciones típicas de la media muestral; Chebyshev da $1/100$.
\end{solucion}

\begin{solucion}{eje:6-unif-lgn}
\datos $U_i\sim\Unif(-1,1)$ independientes.
\planteamiento \cref{prop:t6-unif} y \cref{teo:t6-lgn}.
\pasos
\begin{enumerate}[label=\arabic*.]
  \item $\mu=0$ y $\sigma^{2}=\frac{2^{2}}{12}=\frac13$, finitas; por la ley débil, $\overline U_n\to0$ en probabilidad.
  \item $\PP(\abs{\overline U_{100}}\ge0{,}1)\le\frac{1/3}{100\cdot0{,}01}=\frac13$.
\end{enumerate}
\resultado convergencia a 0; cota $1/3$.
\comprobacion la desviación típica de $\overline U_{100}$ es $\sqrt{1/300}\approx0{,}058$, de modo que $0{,}1$ son unas $1{,}7$ desviaciones típicas y la cota $1/1{,}7^{2}\approx0{,}33$ es coherente.
\end{solucion}

\begin{solucion}{eje:6-encuesta}
\datos $\hat p$ = media de $n$ variables $\Ber(p)$ independientes.
\planteamiento \cref{cor:t6-frecuencias}.
\pasos
\begin{enumerate}[label=\arabic*.]
  \item (a) $\E[\hat p]=p$ y $\Var(\hat p)=p(1-p)/n$. Por Chebyshev, $\PP(\abs{\hat p-p}>c)\le\PP(\abs{\hat p-p}\ge c)\le\frac{p(1-p)}{nc^{2}}$.
  \item $p(1-p)=\frac14-(p-\frac12)^{2}\le\frac14$, luego la cota es $\le\frac1{4nc^{2}}$.
  \item (b) $\frac1{4n\cdot0{,}05^{2}}\le0{,}05\iff n\ge\frac{1}{4\cdot0{,}0025\cdot0{,}05}=2000$.
\end{enumerate}
\resultado la desigualdad del enunciado; $n=2000$ basta.
\comprobacion (ampliación) la aproximación normal del teorema central del límite, que se estudia en el módulo~5, permite reducir el tamaño a $n\approx(1{,}96/(2\cdot0{,}05))^{2}\approx385$; Chebyshev es más conservador porque vale para cualquier distribución.
\end{solucion}

\begin{solucion}{eje:6-cheb-exacta}
\datos $\mu$, $\sigma>0$, $c\ge\sigma$.
\planteamiento colocar la masa en $\mu-c$, $\mu$ y $\mu+c$, de modo simétrico.
\pasos
\begin{enumerate}[label=\arabic*.]
  \item Sea $\PP(X=\mu\pm c)=\frac{\sigma^{2}}{2c^{2}}$ cada uno y $\PP(X=\mu)=1-\frac{\sigma^{2}}{c^{2}}$. Es una ley válida porque $\sigma^{2}/c^{2}\le1$.
  \item $\E[X]=\mu$ por simetría.
  \item $\Var(X)=c^{2}\cdot\frac{\sigma^{2}}{2c^{2}}\cdot2=\sigma^{2}$.
  \item $\PP(\abs{X-\mu}\ge c)=\PP(X=\mu\pm c)=\frac{\sigma^{2}}{c^{2}}$: igualdad en Chebyshev.
\end{enumerate}
\resultado la variable construida alcanza la cota; ninguna cota que dependa solo de $\sigma$ y $c$ puede ser menor.
\comprobacion con $\sigma=1$, $c=2$: masas $\frac18,\frac34,\frac18$; $\PP(\abs{X-\mu}\ge2)=\frac14=\frac1{2^{2}}$.
\end{solucion}

\begin{solucion}{eje:6-cantelli}
\datos media $\mu$, varianza $\sigma^{2}$, $a>0$.
\planteamiento desplazar por $t\ge0$, elevar al cuadrado, aplicar Markov y optimizar en $t$.
\pasos
\begin{enumerate}[label=\arabic*.]
  \item Si $X-\mu\ge a$, entonces $X-\mu+t\ge a+t>0$ y $(X-\mu+t)^{2}\ge(a+t)^{2}$.
  \item Markov: $\PP(X-\mu\ge a)\le\frac{\E[(X-\mu+t)^{2}]}{(a+t)^{2}}=\frac{\sigma^{2}+t^{2}}{(a+t)^{2}}$, porque $\E[X-\mu]=0$.
  \item Sea $h(t)=\frac{\sigma^{2}+t^{2}}{(a+t)^{2}}$. $h'(t)=\frac{2t(a+t)-2(\sigma^{2}+t^{2})}{(a+t)^{3}}=\frac{2(at-\sigma^{2})}{(a+t)^{3}}$, que se anula en $t^{*}=\sigma^{2}/a$, mínimo de $h$ en $[0,\infty)$.
  \item $h(t^{*})=\frac{\sigma^{2}+\sigma^{4}/a^{2}}{(a+\sigma^{2}/a)^{2}}=\frac{\sigma^{2}(a^{2}+\sigma^{2})/a^{2}}{(a^{2}+\sigma^{2})^{2}/a^{2}}=\frac{\sigma^{2}}{\sigma^{2}+a^{2}}$.
\end{enumerate}
\resultado $\PP(X-\mu\ge a)\le\sigma^{2}/(\sigma^{2}+a^{2})$.
\comprobacion con $a=\sigma$ la cota es $1/2$, mientras que Chebyshev bilateral solo daría $\PP(\abs{X-\mu}\ge\sigma)\le1$.
\end{solucion}

\begin{solucion}{eje:6-expo}
\datos $f(x)=\lambda e^{-\lambda x}$, $x\ge0$.
\planteamiento integración por partes con $u=x^{k}$, $dv=\lambda e^{-\lambda x}dx$, $v=-e^{-\lambda x}$; los términos de borde se anulan porque $x^{k}e^{-\lambda x}\to0$ cuando $x\to\infty$.
\pasos
\begin{enumerate}[label=\arabic*.]
  \item $\E[X]=\bigl[-xe^{-\lambda x}\bigr]_0^\infty+\int_0^\infty e^{-\lambda x}dx=0+\frac1\lambda$.
  \item $\E[X^{2}]=\bigl[-x^{2}e^{-\lambda x}\bigr]_0^\infty+\int_0^\infty2xe^{-\lambda x}dx=\frac2\lambda\E[X]=\frac2{\lambda^{2}}$.
  \item $\Var(X)=\frac2{\lambda^{2}}-\frac1{\lambda^{2}}=\frac1{\lambda^{2}}$.
  \item $\PP(\lambda X>t)=\PP(X>t/\lambda)=e^{-t}$ para $t\ge0$: $\lambda X\sim\Expo(1)$.
\end{enumerate}
\resultado $\E[X]=1/\lambda$, $\Var(X)=1/\lambda^{2}$ y $\lambda X\sim\Expo(1)$.
\comprobacion por el último apartado, $\E[X]=\frac1\lambda\E[\lambda X]=\frac1\lambda\cdot1$, coherente con el primero.
\end{solucion}

\begin{solucion}{eje:6-afin-normal}
\datos $X\sim\Nor(\mu,\sigma^{2})$, $Y=aX+b$, $a\neq0$.
\planteamiento calcular $F_Y$ y reconocer la función de distribución de una normal, $\Phi((y-m)/s)$ (\cref{teo:t6-tipificacion}).
\pasos
\begin{enumerate}[label=\arabic*.]
  \item $a>0$: $F_Y(y)=\PP\bigl(X\le\frac{y-b}a\bigr)=\Phi\Bigl(\frac{(y-b)/a-\mu}{\sigma}\Bigr)=\Phi\Bigl(\frac{y-(a\mu+b)}{a\sigma}\Bigr)$.
  \item $a<0$: al dividir por $a$ se invierte la desigualdad: $F_Y(y)=\PP\bigl(X\ge\frac{y-b}a\bigr)=1-\Phi\Bigl(\frac{(y-b)/a-\mu}\sigma\Bigr)=\Phi\Bigl(\frac{y-(a\mu+b)}{\abs a\sigma}\Bigr)$, usando $1-\Phi(z)=\Phi(-z)$ y $-\frac{(y-b)/a-\mu}{\sigma}=\frac{y-b-a\mu}{-a\sigma}$.
  \item En ambos casos $F_Y(y)=\Phi\bigl((y-m)/s\bigr)$ con $m=a\mu+b$ y $s=\abs a\sigma$: $Y\sim\Nor(a\mu+b,a^{2}\sigma^{2})$.
  \item Con $a=-1$, $b=0$ y $X=Z\sim\Nor(0,1)$: $-Z\sim\Nor(0,1)$.
\end{enumerate}
\resultado $aX+b\sim\Nor(a\mu+b,\,a^{2}\sigma^{2})$.
\comprobacion la media y la varianza coinciden con las que da la linealidad: $\E[Y]=a\mu+b$ y $\Var(Y)=a^{2}\sigma^{2}$.
\end{solucion}


\backmatter
% =====================================================================
\begin{thebibliography}{99}
\addcontentsline{toc}{chapter}{Bibliografía}
\markboth{Bibliografía}{}
\small

\bibitem{bt}
D.~P. Bertsekas y J.~N. Tsitsiklis.
\newblock \emph{Introduction to Probability}, 2.ª ed.
\newblock Athena Scientific, 2008.
\newblock Texto del curso 6.041 del MIT.

\bibitem{bh}
J.~K. Blitzstein y J.~Hwang.
\newblock \emph{Introduction to Probability}, 2.ª ed.
\newblock Chapman \& Hall/CRC Press, 2019.
\newblock Texto del curso Stat~110 de Harvard.

\bibitem{stat110}
J.~K. Blitzstein.
\newblock Stat 110: Probability.
\newblock Harvard University. \emph{Strategic Practice and Homework}~1 a~6 y~10, otoño de 2011.
\newblock \url{https://stat110.hsites.harvard.edu/strategic-practice-problems}

\bibitem{chen}
M.~Chen \emph{et al.}
\newblock Evaluating large language models trained on code.
\newblock arXiv:2107.03374, 2021.
\newblock \url{https://arxiv.org/abs/2107.03374}

\bibitem{gs}
C.~M. Grinstead y J.~L. Snell.
\newblock \emph{Introduction to Probability}, 2.ª ed. revisada.
\newblock American Mathematical Society. Versión libre del proyecto CHANCE de 4 de julio de 2006, bajo licencia GNU FDL.
\newblock \url{https://chance.dartmouth.edu/teaching_aids/books_articles/probability_book/book.html}

\bibitem{mcs}
E.~Lehman, F.~T. Leighton y A.~R. Meyer.
\newblock \emph{Mathematics for Computer Science}.
\newblock Edición del curso 6.042J del MIT, primavera de 2015 (MIT OpenCourseWare). Se citan los capítulos~14 y~16 a~19 de esa edición.
\newblock \url{https://ocw.mit.edu/courses/6-042j-mathematics-for-computer-science-spring-2015/pages/readings/}

\bibitem{mit6042}
A.~R. Meyer y A.~Chlipala.
\newblock 6.042J: Mathematics for Computer Science.
\newblock Massachusetts Institute of Technology, MIT OpenCourseWare, primavera de 2015.
\newblock \url{https://ocw.mit.edu/courses/6-042j-mathematics-for-computer-science-spring-2015/}

\bibitem{mit1805}
J.~Orloff y J.~Bloom.
\newblock 18.05: Introduction to Probability and Statistics.
\newblock Massachusetts Institute of Technology, MIT OpenCourseWare, primavera de 2022. \emph{Problem Sets}~1 a~3.
\newblock \url{https://ocw.mit.edu/courses/18-05-introduction-to-probability-and-statistics-spring-2022/}

\bibitem{ross}
S.~Ross.
\newblock \emph{A First Course in Probability}, 10.ª ed.
\newblock Pearson, 2024.

\bibitem{mit6041f10}
J.~N. Tsitsiklis.
\newblock 6.041: Probabilistic Systems Analysis and Applied Probability.
\newblock Massachusetts Institute of Technology, MIT OpenCourseWare, otoño de 2010. \emph{Problem Sets}~1 a~4 y~9.
\newblock \url{https://ocw.mit.edu/courses/6-041-probabilistic-systems-analysis-and-applied-probability-fall-2010/}

\end{thebibliography}


\end{document}
